\chapter*{Summary of Part~\ref{part:standard-landscape}}
\label{ch:landscape-census}\label{chap:landscape-census}
\addcontentsline{toc}{chapter}{Summary of Part~\ref{part:standard-landscape}}

% Cross-volume notation: \kappaChHodge is reserved for the Hodge
% supertrace Xi(X) of the CY construction.  The OPE-derived genus-one
% scalar in the Open quadrant is denoted \kappa_{\mathrm{mod}} below.
% The macro \kappaChHodge is shared with Vol II and Vol III.
\providecommand{\kappaChHodge}{\kappa^{\mathrm{Hodge}}_{\mathrm{ch}}}
\providecommand{\kappaBKM}{\kappa_{\mathrm{BKM}}}

The complementarity sum $c + c' = 2d$ holds for every non-critical
Kac--Moody algebra on the Feigin--Frenkel level-shift locus. The
universal generating function $x/\sin x$ governs the scalar
Faber--Pandharipande genus series on the proved uniform-weight locus
through the Wick-rotated identity
$(x/2)/\sin(x/2)=\widehat{A}(ix)$; multi-weight families carry
additional cross-channel corrections. The rank-one branch divisor
\[
\Delta(x) = 1 - 2x - 3x^2=(1-3x)(1+x)
\]
governs the Virasoro and $\beta\gamma$ bar generating functions and is
the Drinfeld--Sokolov branch core adjacent to
$\widehat{\mathfrak{sl}}_2$.  The corrected affine
$\widehat{\mathfrak{sl}}_2$ series has coefficient denominator
$(1-x)^2$.  The census displays the two denominators together with
their respective bar constructions.

Each chirally Koszul vertex algebra in the currently computed and
Koszul-closed rows of the standard-landscape census carries a specified
shadow datum: the scalar $\kappa_{\mathrm{mod}}$, the component on which it is measured,
the cubic and quartic coefficients $S_3,S_4$, and the first
non-vanishing higher obstruction after the quartic term. These data
stratify the completed vertex-algebra rows into the four archetype
classes G/L/C/M: $S_3=S_4=0$ characterizes Gaussian algebras
(class~G); $S_3\ne0$ and $S_4=0$ singles out the affine
Kac--Moody tree component (class~L); $S_3=0$, $S_4\ne0$, and
$S_r=0$ for $r\ge5$ is the contact/quartic component
(class~C); class~M has nonzero stress-tensor or higher-spin
coefficients in arbitrarily high degrees.  The critical discriminant
$\Delta=8\kappa_{\mathrm{mod}} S_4$ is a derived scalar on a chosen component; the
archetype is determined by the full sequence
\((S_3,S_4,S_5,\ldots)\).  A frontier row records its quadratic
comparison \(q_A:A^{\mathrm i}\to\bar B(A)\) and its unresolved
cohomology explicitly.

\paragraph{Distinguished chart on the Koszulness moduli.}
Each shadow-class stratum selects a distinguished chart on the
Koszulness moduli scheme
$\cM_{\mathrm{Kosz}}(\cA)$ of
Theorem~\ref{thm:koszulness-moduli-kp}
(see~\S\ref{subsec:koszulness-moduli}):
class~G sits at $\Phi_{\mathrm{KZ}}$ (the core ten charts
$U_1,\ldots,U_{10}$); the generic affine Kac--Moody class~L locus sits
at $\Phi_{\mathrm{ell}}$ (Enriquez elliptic chart $U_{13}$); class~C at
$\Phi_{\mathrm{AT}}$ (Alekseev--Torossian chart $U_{11}$);
class~M at $\Phi_\infty = \Phi_{\mathrm{Kon}}$ (Kontsevich
integral chart $U_{14}$); and the Feigin--Frenkel critical
stratum FF at $\Phi_{\mathrm{dR,B}}$ (Hodge--Betti chart
$U_{12}$). Family-by-family distinguished charts are recorded in
the last column of Table~\ref{tab:master-invariants-chart}
below.

% Regimes I--III : census spans all families (Convention~\ref{conv:regime-tags}).
The table below records the computed invariant rows and the
conjectural rows treated in Part~\ref{part:physics-bridges}.  A
computed row gives one algebra on its stated parameter locus.  A
frontier row gives the specific formula and comparison morphism whose
status is displayed in that row.  The columns place the family
computations in a common normalization.
The scalar modular characteristic
$\kappa_{\mathrm{mod}}(\widehat{\fg}_k) = td/2h^\vee$
and $\kappa_{\mathrm{mod}}(\mathrm{Vir}_c) = c/2$ both descend from the universality
theorem (Theorem~\ref{thm:genus-universality}).

\noindent\emph{The Heisenberg invariants are derived in
Chapter~\ref{ch:heisenberg-frame}.  Each subsequent computed row
cites the proposition carrying its formula; each frontier row names
the comparison morphism that remains to be constructed.}

\begin{table}[ht]
\centering
\caption{Distinguished chart on
 $\cM_{\mathrm{Kosz}}(\cA)$ per family
 (Theorem~\ref{thm:koszulness-moduli-kp}).}
\label{tab:master-invariants-chart}
\renewcommand{\arraystretch}{1.3}
\begin{tabular}{lccl}
\toprule
Family & Class & Distinguished chart & Rationale \\
\midrule
Heisenberg $\cH_k$ & G & $\Phi_{\mathrm{KZ}}$ ($U_1$)
 & abelian shadow; KZ-flat \\
Lattice $V_L$ & G & $\Phi_{\mathrm{KZ}}$ ($U_1$--$U_{10}$)
 & abelian shadow \\
$\widehat{\fg}_k$ generic $k$ & L & $\Phi_{\mathrm{ell}}$ ($U_{13}$)
 & genus-1 Enriquez; $S_3\neq 0$ \\
$\widehat{\fg}_k$ at $k=-h^\vee$ & FF & $\Phi_{\mathrm{dR,B}}$ ($U_{12}$)
 & Feigin--Frenkel center; Hodge--Betti \\
$\beta\gamma$ system & C & $\Phi_{\mathrm{AT}}$ ($U_{11}$)
 & quartic bracket; AT Lagrangian \\
$\Vir_c$ generic & M & $\Phi_{\mathrm{Kon}} = \Phi_\infty$ ($U_{14}$)
 & infinite tower; chord-diagram \\
$\cW_N$ principal & M & $\Phi_{\mathrm{Kon}} = \Phi_\infty$ ($U_{14}$)
 & DS-transported from Virasoro \\
$\cW^{(2)}_k(\mathfrak{sl}_3)$ & M & $\Phi_{\mathrm{Kon}} = \Phi_\infty$ ($U_{14}$)
 & minimal nilpotent DS \\
\bottomrule
\end{tabular}
\end{table}

\paragraph{Standard families: objects, formulas, and sources.}
\label{par:standard-family-source-data}
For a family \(A_U\) over a parameter locus \(U\), the census records
the singular products defining \(A_U\), the formulas derived from
those products, the category in which the derivation takes place, and
the theorem that supplies it.  The first table gives the local
formulas and their domains.  The next two tables record the comparison
morphisms that carry these formulas from the OPE presentation to the
bar, quadratic-dual, Verdier-dual, and derived-centre constructions.
Every later use of a row inherits the domain and status displayed
here.

For the principal family, write
\[
 H_{W_N}^{\mathrm{mod/comp}}
 :=H_{\mathrm{diag}}^{g=1}+H_{W_N}^{\mathrm{DS/bar}}.
\]
The first summand constructs the normalized diagonal genus-one trace;
the second identifies the reflected principal DS level with the
completed bar--Verdier companion.  The specialization at \(N=3\) is
\(H_{W_3}^{\mathrm{mod/comp}}\).
The rational-function identity
\(K_N^c=4N^3-2N-2\), and hence \(K_3^c=100\), is exact arithmetic.
The formulas for \(\kappa_{\mathrm{mod}}(\mathcal W_N)\),
\(K_N^{\kappa_{\mathrm{mod}}}\), and the value \(250/3\) are conditional
on the explicitly displayed package
\(H_{\mathrm{diag}}^{g=1}+H_{W_N}^{\mathrm{DS/bar}}\).

{\scriptsize
\renewcommand{\arraystretch}{1.18}
\begin{longtable}{@{}p{0.135\textwidth}p{0.16\textwidth}p{0.385\textwidth}p{0.205\textwidth}p{0.09\textwidth}@{}}
\caption{Defining formulas, parameter loci, and mathematical sources}
\label{tab:standard-family-status}\\
\toprule
Family & Parameter locus & Defining formula and derived invariant &
Source & Status \\
\midrule
\endfirsthead
\toprule
Family & Parameter locus & Defining formula and derived invariant &
Source & Status \\
\midrule
\endhead
Heisenberg \(\cH_k\) &
\(k\in\mathbb C\) &
\(J(z)J(w)\sim k(z-w)^{-2}\),
\(d[J|J]=k\mathbf1\),
\(r^{\mathrm{coll}}(z)=k/z\), and
\(\kappa_{\mathrm{mod}}(\cH_k)=k\). &
Proposition~\ref{prop:heisenberg-standard-family-ledger};
Corollary~\ref{cor:heisenberg-closed-sector-ledger} for the centre. &
proved here; centre conditional \\

Affine \(V_k(\fg)\) &
\(\fg\) simple, \(k\ne-h^\vee\) &
\(J^a(z)J^b(w)\sim
k\kappa_\fg^{ab}(z-w)^{-2}
+\,
f^{ab}{}_cJ^c(w)(z-w)^{-1}\);
\(\mathfrak C(x,y,z)=\kappa_\fg(x,[y,z])\),
\(\{\mathfrak C,\mathfrak C\}=0\),
\(r_{\max}=3\), and
\(\kappa_{\mathrm{mod}}=d(k+h^\vee)/(2h^\vee)\). &
Proposition~\ref{prop:affine-standard-family-ledger}. &
proved here on the current sector \\

\(\beta\gamma_\lambda\), \(bc_\lambda\) &
\(\lambda\in\mathbb C\); standard contact locus &
\(\beta(z)\gamma(w)\sim(z-w)^{-1}\),
\(\kappa_{\mathrm{mod}}(\beta\gamma_\lambda)=6\lambda^2-6\lambda+1\),
\(\kappa_{\mathrm{mod}}(bc_\lambda)
=-\kappa_{\mathrm{mod}}(\beta\gamma_\lambda)\),
\(S_3=0\), \(S_4=-5/12\), \(S_r=0\) for \(r\ge5\).
At \(\lambda=\tfrac12\), \(c_{\beta\gamma}=-1\). &
Proposition~\ref{prop:betagamma-standard-family-ledger}. &
conditional \\

\(\Vir_c\) &
\(c(5c+22)\ne0\) for the quartic formula &
\(T(z)T(w)\sim
\frac{c/2}{(z-w)^4}+\frac{2T(w)}{(z-w)^2}
+\frac{\partial T(w)}{z-w}\);
\(\kappa_{\mathrm{mod}}=c/2\), \(S_3=2\), and
\(S_4=10/[c(5c+22)]\).
The same-family comparison has target \(\Vir_{26-c}\). &
Proposition~\ref{prop:virasoro-standard-family-ledger}. &
conditional \\

Principal \(\mathcal W_N\) &
\(N\ge3\); generic principal DS locus &
\(\Lambda={:}TT{:}-(3/10)\partial^2T\),
\(\alpha(c)=16/(22+5c)\),
the exact identity \(K_N^c=4N^3-2N-2\), and, on
\(H_{\mathrm{diag}}^{g=1}+H_{W_N}^{\mathrm{DS/bar}}\),
\(\kappa_{\mathrm{mod}}(\mathcal W_N)=c(H_N-1)\) and
\(K_N^{\kappa_{\mathrm{mod}}}=(H_N-1)K_N^c\).
Thus \(K_3^c=100\) exactly and
\(K_3^{\kappa_{\mathrm{mod}}}=250/3\) conditionally. &
Proposition~\ref{prop:principal-WN-standard-family-ledger}. &
exact central arithmetic; modular statement conditional on
\(H_{\mathrm{diag}}^{g=1}+H_{W_N}^{\mathrm{DS/bar}}\) \\

Critical affine \(V_{-h^\vee}(\fg)\) &
\(k=-h^\vee\) &
The current OPE has level \(-h^\vee\), the shifted bar curvature is
zero, and
\(\mathfrak z(\widehat{\fg}_{-h^\vee})
\cong\mathcal O(\mathrm{Op}_{\fg^\vee}(D))\). &
Proposition~\ref{prop:km-critical-separation};
Theorem~\ref{thm:critical-level-structure}. &
proved here / proved elsewhere \\

Subregular \(\mathcal W_n^{(2)}(k)\) &
\(n\ge3\), generic \(k\), Li/PBW subregular locus &
\[
C_{n,r}(h_n,v)
=\sum_{t=0}^{r}\binom{n-r+t}{t}e_{r-t}(v)h_n^t,
\qquad 2\le r\le n-1,
\]
and the canonical completed dual has degree \(n-1\). &
Theorems~\ref{thm:w-subregular-appell} and
\ref{thm:w-subregular-classification}. &
proved here on the stated locus \\

\(\mathcal N=2\) SCA &
\(c=3k/(k+2)\); universal OPE and unitary Kazama--Suzuki loci &
Generators \(\{T,J,G^+,G^-\}\) carry the standard OPE packet and
integer spectral flow.  The modular formula
\(\kappa=(k+4)/4=(6-c)/[2(3-c)]\) is conditional on the
trace-compatible genus-one coset package
\(H_{\mathcal N=2}^{\mathrm{mod}}\); quadratic recognition is
conditional on \(H_{\mathcal N=2}^{\mathrm{bar}}\). &
Propositions~\ref{prop:n2-standard-family-ledger} and
\ref{prop:n2-kappa}. &
OPE/coset proved on the stated loci; modular/bar comparisons conditional \\

Lattice \(V_\Lambda^\varepsilon\) and
\(V_\Lambda^{N,q}\) &
\(\Lambda\) even; \(\varepsilon\) normalized;
\(N\ge3\), \(2q\ne0\) on the ordered deformation &
\(\varepsilon(\alpha,\beta)
=(-1)^{\langle\alpha,\beta\rangle}\varepsilon(\beta,\alpha)\)
is the vertex-locality equation, while
\(d[e^\alpha|e^\beta]
=\varepsilon(\alpha,\beta)e^{\alpha+\beta}\)
for \(\langle\alpha,\beta\rangle=-1\).
The dual deformation carries \(-q\). &
Proposition~\ref{prop:lattice-standard-cocycle-ledger};
Theorem~\ref{thm:quantum-lattice-structure}. &
proved here \\

Moonshine \(V^\natural\) &
Virasoro/Griess/Monster presentation at \(c=24\) &
\(\dim V^\natural_1=0\), and the Virasoro-sector scalar is
\(\kappa_T=c/2=12\).  McKay--Thompson coefficients belong to the
Monster trace functions \(T_g\). &
Proposition~\ref{prop:moonshine-kappa};
Proposition~\ref{prop:canonical-monster-atlas-values}. &
conditional scalar / proved trace data \\

Symmetric orbifold &
\(N\ge0\); seed elliptic genus with finite graded coefficients &
\[
\sum_{N\ge0}p^NZ(\operatorname{Sym}^N X)
=\prod_{\substack{n\ge0,\;m\ge1\\ \ell}}
(1-q^ny^\ell p^m)^{-c_X(mn,\ell)}.
\]
The ordered fixed-point question is the Reynolds comparison
\(\bar B(X^{\otimes N})^{S_N}\to
\bar B((X^{\otimes N})^{S_N})\). &
Theorem~\ref{thm:symn-dmvv-product};
Remark~\ref{rem:symn-orbifold-functor}. &
proved character identity; bar comparison open \\

Triplet \(\mathcal W(p)\) &
\(p\ge2\) &
\(c_p=1-6(p-1)^2/p\);
\(\mathcal W(p)/C_2(\mathcal W(p))\) is finite-dimensional.
The quadratic comparison
\(q_p:\mathcal W(p)^{\mathrm i}\to\bar B(\mathcal W(p))\)
is the open Koszul construction. &
Proposition~\ref{prop:wp-c2-cofinite};
Conjecture~\ref{conj:wp-koszulness}. &
proved \(C_2\)-finiteness; Koszulness conjectured \\

Admissible affine quotient \(L_k(\fg)\) &
\(k+h^\vee=p/q>0\), \((p,q)=1\) &
The quotient \(V_k(\fg)\twoheadrightarrow L_k(\fg)\) induces
\[
\bar B(V_k(\fg))\longrightarrow\bar B(L_k(\fg)).
\]
Its cone is the quotient-bar obstruction; each fixed-weight sector
is finite on the non-degenerate admissible locus. &
Proposition~\ref{prop:bar-admissible};
Corollary~\ref{cor:bar-admissible-finiteness}. &
proved finite windows; global comparison open \\

Super-Yangian \(Y_\hbar(\mathfrak{gl}_{m|n})\) &
\(m,n\ge0\), parity convention fixed &
The graded RTT equation
\(R(u-v)T_1(u)T_2(v)=T_2(v)T_1(u)R(u-v)\)
defines the relation space; the quantum Berezinian and supertrace fix
the central and scalar normalizations. &
Proposition~\ref{prop:super-yangian-parity-sign-ledger}. &
conditional \\

Exceptional affine and Yangian rows &
\(\fg\in\{G_2,F_4,E_6,E_7,E_8\}\);
finite RTT window \(N\) &
The affine row uses the current formula above with the root data of
\(\fg\).  The Yangian row is determined in window \(N\) by
\(A_{\fg,\le N}=T(V_{\le N})/\langle R_{\fg,\le N}\rangle\)
and the inverse-relation comparison
\(\rho_{\fg,N}:R_{\fg,\le N}^{\perp}\to
R_{\fg,\le N}(-\hbar)\). &
Proposition~\ref{prop:complete-exceptional-shadow};
Theorem~\ref{thm:exceptional-yangian-koszul-duality-all-five-types}. &
affine formulas proved; Yangian comparison conditional \\
\bottomrule
\end{longtable}
}

{\scriptsize
\renewcommand{\arraystretch}{1.18}
\begin{longtable}{@{}p{0.17\textwidth}p{0.24\textwidth}p{0.25\textwidth}p{0.30\textwidth}@{}}
\caption{Presentations, finite calculations, and comparison morphisms}
\label{tab:standard-family-ambient}\\
\toprule
Family & OPE or algebraic presentation & Finite calculation &
Ambient construction or comparison morphism \\
\midrule
\endfirsthead
\toprule
Family & OPE or algebraic presentation & Finite calculation &
Ambient construction or comparison morphism \\
\midrule
\endhead
Heisenberg &
rank-one current OPE &
ordered bar through degree \(3\); finite Fock weights &
completed Fock category; full bar--cobar reconstruction and the
quadratic comparison \(q_{\cH}\) \\

Affine \(V_k(\fg)\) &
current OPE with invariant form \(\kappa_\fg\) &
bracket differential, KZ windows, fixed conformal weights &
weight-completed current bar; PBW comparison
\(q_{V_k(\fg)}:V_k(\fg)^{\mathrm i}\to\bar B(V_k(\fg))\) \\

Critical affine &
critical current OPE &
fixed conformal weights and Feigin--Frenkel generators &
completed critical bar and the oper-valued degree-zero centre \\

\(\beta\gamma_\lambda\), \(bc_\lambda\) &
Weyl and Clifford OPEs &
ordered degrees \(2\) and \(3\), parity, contact direction &
finite-type Verdier duality and continuous duals \\

\(\Vir_c\), \(\mathcal W_N\) &
stress-tensor and DS strong-generator OPEs &
Kac and quasi-primary windows; finitely many spin channels &
weight-completed bars; the \(\mathcal W_N\) comparison is transported
along the DS/bar morphism \\

Subregular \(\mathcal W_n^{(2)}\) &
KRW strong generators and Miura product &
Slodowy/PBW symbols \(C_{n,r}(h_n,v)\) &
completed subregular bar on the Li/PBW locus \\

Lattice and quantum lattice &
even lattice and bimultiplicative \(2\)-cocycle &
charge sectors and adjacent-root \(A_2\) complex &
sectorwise completed ordered bar; inverse-cocycle duality \\

Moonshine &
Monster VOA, Griess product, and McKay--Thompson traces &
Virasoro channel, weight-\(2\) Griess algebra, finite trace
coefficients &
Monster/Borcherds completion; a K3/Mukai comparison is a separate
functorial construction \\

Symmetric orbifold &
\((X^{\otimes N})^{S_N}\) and \(S_N\)-twisted sectors &
finite \(N\), fixed-point bar, DMVV coefficients &
Reynolds bar map for the fixed-point algebra; \(G\)-crossed
completion for twist-field OPEs \\

Triplet and admissible quotients &
screening kernel \(\mathcal W(p)\) or quotient \(L_k(\fg)\) &
null-vector, \(C_2\), Li, and fixed-weight bar complexes &
the cones of \(q_p\) and of
\(\bar B(V_k(\fg))\to\bar B(L_k(\fg))\) are the open obstruction
complexes \\

Super and exceptional Yangians &
parity-signed or exceptional RTT relations &
finite module, Berezinian, supertrace, and relation windows &
finite-window maps \(\rho_{\fg,N}\); the inverse limit requires
Mittag--Leffler compatibility and completed PBW \\

Exceptional affine Kac--Moody &
current OPE for \(G_2,F_4,E_6,E_7,E_8\) &
finite root and current sectors &
the universal affine completion, governed by the same PBW map as the
generic simple-current family \\
\bottomrule
\end{longtable}
}

The full bar coalgebra and the quadratic Koszul dual coalgebra are
different constructions:
\[
\bar B(A)=T^c(s^{-1}\bar A),\qquad
q_A:A^{\mathrm i}\longrightarrow\bar B(A).
\]
Koszulness means that \(q_A\) is a quasi-isomorphism in the ambient
category displayed above.  On the dualizable Ran locus,
\(A^!=\mathbb D_{\Ran}(A^{\mathrm i})\), while
\(Z_{\mathrm{ch}}^{\mathrm{der}}(A)
=C^\bullet_{\mathrm{ch}}(A,A)\) is the derived chiral centre.

{\scriptsize
\renewcommand{\arraystretch}{1.18}
\begin{longtable}{@{}p{0.15\textwidth}p{0.18\textwidth}p{0.26\textwidth}p{0.20\textwidth}p{0.19\textwidth}@{}}
\caption{Bar, quadratic-dual, Verdier-dual, and centre constructions}
\label{tab:standard-family-five-object-package}\\
\toprule
Family & \(A\) & \(\bar B(A)\), \(A^{\mathrm i}\), and \(q_A\) &
\(A^!\) & \(Z^{\mathrm{der}}_{\mathrm{ch}}(A)\) \\
\midrule
\endfirsthead
\toprule
Family & \(A\) & \(\bar B(A)\), \(A^{\mathrm i}\), and \(q_A\) &
\(A^!\) & \(Z^{\mathrm{der}}_{\mathrm{ch}}(A)\) \\
\midrule
\endhead
Heisenberg &
\(\cH_k\) &
full tensor bar and the Gaussian quadratic comparison \(q_{\cH}\) &
\((\operatorname{Sym}^{\mathrm{ch}}(V^*[1]),m_0=-k\omega)\) &
\(C^\bullet_{\mathrm{ch}}(\cH_k,\cH_k)\) \\

Affine, non-critical &
\(V_k(\fg)\) &
completed current bar and the PBW/CE comparison \(q_{k,\fg}\) &
\(\mathrm{CE}^{\mathrm{ch}}
(\widehat{\fg}_{-k-2h^\vee})\) on the stated dualizable locus &
\(C^\bullet_{\mathrm{ch}}(V_k(\fg),V_k(\fg))\) \\

Affine, critical &
\(V_{-h^\vee}(\fg)\) &
critical completed bar and its critical quadratic comparison &
\(\mathbb D_{\Ran}(V_{-h^\vee}(\fg)^{\mathrm i})\) &
\(\mathfrak z(\widehat{\fg}_{-h^\vee})
\hookrightarrow H^0Z^{\mathrm{der}}_{\mathrm{ch}}(A)\), with
\(\mathfrak z\cong\mathcal O(\mathrm{Op}_{\fg^\vee}(D))\) \\

\(\beta\gamma_\lambda\) &
bosonic Weyl chiral algebra &
ordered full bar and the free-field quadratic comparison &
\(bc_\lambda\) under parity-changing Verdier duality &
\(C^\bullet_{\mathrm{ch}}(\beta\gamma_\lambda,\beta\gamma_\lambda)\) \\

\(\Vir_c\) &
\(\Vir_c\) &
completed Virasoro bar and the same-family comparison \(q_c\) &
\(\Vir_{26-c}\) on the conditional same-family locus &
\(C^\bullet_{\mathrm{ch}}(\Vir_c,\Vir_c)\) \\

Principal \(\mathcal W_N^k\) &
principal DS algebra &
completed DS bar and the transported comparison \(q_{N,k}^{\mathrm{DS}}\) &
\(\mathcal W_N^{k'}\) on the DS/bar transport locus &
\(C^\bullet_{\mathrm{ch}}(\mathcal W_N^k,\mathcal W_N^k)\) \\

Lattice &
\(V_\Lambda^\varepsilon\) &
sectorwise full bar; \(A^{\mathrm i}\) carries the inverse cocycle &
\(\mathbb D_{\Ran}(A^{\mathrm i})\);
\(A^!\simeq A\) on the unimodular locus &
sectorwise \(\Eone\)-chiral Hochschild cochains \\

Moonshine &
\(V^\natural\) &
Virasoro/Griess/Borcherds sector bars; the full quadratic comparison
is an open construction &
\(\Vir_2\) on the Virasoro sector; full VOA dual open &
sectorwise centre; full Monster comparison open \\

Symmetric orbifold &
\((A^{\otimes N})^{S_N}\) &
fixed-point full bar with Reynolds comparison from
\(\bar B(A^{\otimes N})^{S_N}\) &
Verdier dual after the fixed-point comparison &
fixed-point chiral Hochschild cochains; twisted centre after
\(G\)-crossed completion \\

Triplet / admissible quotient &
\(\mathcal W(p)\) or \(L_k(\fg)\) &
\(q_p\) or the quotient-bar map displayed above &
dual algebra defined by the corresponding quadratic comparison &
the chiral Hochschild complex, including its positive-depth
obstruction groups \\

Super / exceptional Yangian &
\(A_{\fg,\le N}=T(V_{\le N})/\langle R_{\fg,\le N}\rangle\) &
finite full bar and inverse-relation map \(\rho_{\fg,N}\) &
finite-window inverse-RTT dual on the locus where \(\rho_{\fg,N}\)
is invertible &
finite-window Hochschild cochains; pro-centre after compatible inverse
limit \\

Exceptional affine Kac--Moody &
\(V_k(\fg)\), \(\fg\in\{G_2,F_4,E_6,E_7,E_8\}\) &
the universal completed current bar and affine PBW comparison &
\(\mathrm{CE}^{\mathrm{ch}}
(\widehat{\fg}_{-k-2h^\vee})\) on the affine dualizable locus &
\(C^\bullet_{\mathrm{ch}}(V_k(\fg),V_k(\fg))\) \\
\bottomrule
\end{longtable}
}

\providecommand{\PhiFA}{\Phi^{\mathrm{FA}}}
\providecommand{\SpCh}{\mathrm{Sp}^{\mathrm{ch}}}
\providecommand{\EdHolFA}{E_d\text{-}\mathrm{HolFA}}
\providecommand{\EnHolFA}{E_n\text{-}\mathrm{HolFA}}

\section*{The $5\times 5$ $\kappa$-stratification matrix on
G\,/\,L\,/\,C\,/\,M\,/\,B}
\label{sec:five-by-five-kappa-matrix}
\index{$\kappa$-stratification matrix|textbf}
\index{archetype B@archetype $\mathbf B$ (K3$\times E$)|textbf}

The five columns of the matrix are five functions with different
domains:
\[
 \kappa_{\mathrm{cat}},\qquad
 \kappa^{\mathrm{Hodge}}_{\mathrm{ch}},\qquad
 \kappa^{\mathrm{Heis}}_{\mathrm{ch}},\qquad
 \kappa_{\mathrm{BKM}},\qquad
 \kappa_{\mathrm{fiber}}.
\]
For a compact CY host \(X\),
\(\kappa_{\mathrm{cat}}(X)=\chi(\cO_X)\) and
\(\kappa^{\mathrm{Hodge}}_{\mathrm{ch}}(X)
=\Xi(X)=\sum_q(-1)^qh^{0,q}(X)\).
A chosen Heisenberg specialisation defines
\(\kappa^{\mathrm{Heis}}_{\mathrm{ch}}\);
a Borcherds input
\(\Phi\) defines
\(\kappa_{\mathrm{BKM}}(\Phi)=c^\Phi(0)/2\);
the chosen transverse fibre defines
\(\kappa_{\mathrm{fiber}}\) by its stated rank or Euler
characteristic.  A row therefore consists of a host, a specialisation
functor, and five evaluations, each on its own source object.
The Heisenberg coordinate is level-weighted throughout the census:
the rank-one Heisenberg algebra at level~$k$ carries
\begin{equation}\label{eq:heisenberg-kappa}
\kappa(\cH_k) = k,
\end{equation}
so a rank-$r$ Heisenberg specialisation at level $k$ carries
\(\kappa^{\mathrm{Heis}}_{\mathrm{ch}}=rk\), a level-one
specialisation carries \(\kappa^{\mathrm{Heis}}_{\mathrm{ch}}=r\),
and the Cartan-oscillator specialisation of \(V_k(\fg)\) carries
\(\kappa^{\mathrm{Heis}}_{\mathrm{ch}}(V_k(\fg))
=\operatorname{rank}(\fg)\cdot k\).
On a compact complex host, Dolbeault theory gives the canonical
numerical equality
\[
\kappa_{\mathrm{cat}}(X)
=\sum_q(-1)^q\dim H^q(X,\cO_X)
=\sum_q(-1)^qh^{0,q}(X)
=\kappa_{\mathrm{ch}}^{\mathrm{Hodge}}
\bigl(\PhiFA_d(\Perf X)\bigr).
\]

\begin{theorem}[The five coordinates of
\texorpdfstring{$K3\times E$}{K3 x E};
\ClaimStatusConditional]
\label{thm:k3e-five-coordinate-computation}
\index{K3 constants!five coordinates}
\index{kappa-stratification matrix@$\kappa$-stratification matrix!K3 coordinates}
\textbf{Type signature.} (CY quadrant; presentation
\(\Phi_3^{(\Sigma_2,C)}
=\SpCh_{\Sigma_2,C}\circ\PhiFA_3\);
levels \(0\), \(1\), and \(5\) according to the coordinate;
hypothesis package: compact host \(K3\times E\),
specified transverse cycle \(\Sigma_2\), curve \(C\), Gritsenko
Jacobi input for \(\Delta_5\), and the Mukai lattice of the K3
fibre.)
The five evaluations are
\[
\begin{aligned}
\kappa_{\mathrm{cat}}(K3\times E)&=0,&
\kappa_{\mathrm{ch}}^{\mathrm{Hodge}}(K3\times E)&=0,\\
\kappa_{\mathrm{ch}}^{\mathrm{Heis}}(K3\times E)&=3,&
\kappa_{\mathrm{BKM}}(\Delta_5)&=5,&
\kappa_{\mathrm{fiber}}(K3)&=24.
\end{aligned}
\]
The first, second, fourth, and fifth equalities follow from their
definitions.  Volume~III defines the third coordinate by the additive
normalization
\[
\kappa_{\mathrm{ch}}^{\mathrm{Heis}}(K3\times E)
:=
\kappa_{\mathrm{ch}}^{\mathrm{Heis}}(K3)
+\kappa_{\mathrm{ch}}^{\mathrm{Heis}}(E)
=2+1.
\]
Its realization in the Open quadrant is the trace-preserving
comparison morphism
\[
\eta_{K3,E}\colon
\SpCh_{\Sigma_2,C}
\bigl(\PhiFA_3(\Perf(K3\times E))\bigr)
\longrightarrow \mathcal H_1^{\oplus3}.
\]
The equality \(\kappa_{\mathrm{ch}}^{\mathrm{Heis}}=3\) in Volume~I
is conditional on constructing \(\eta_{K3,E}\) and proving that it
carries the specialization trace to the rank-three level-one
Heisenberg trace.
\end{theorem}

\begin{proof}
Each coordinate is evaluated from its definition.
\begin{enumerate}[label=\textup{(\roman*)},leftmargin=*]
\item K\"unneth multiplicativity gives
\[
\kappa_{\mathrm{cat}}(K3\times E)
=\chi(\cO_{K3})\chi(\cO_E)
=(1+1)(1-1)=0.
\]
\item K\"unneth gives
\(h^{0,\bullet}(K3\times E)=(1,1,1,1)\).  Hence
\[
\kappa_{\mathrm{ch}}^{\mathrm{Hodge}}(K3\times E)
=\Xi(K3\times E)=1-1+1-1=0,
\]
in agreement with Volume~III
Theorem~\ref{V3-thm:cy-d-tri-stratum}.
\item The additive Heisenberg normalization in Volume~III
Remark~\ref{V3-rem:four-kappa-stage-assignment-cy-d-strat} assigns
\[
\kappa_{\mathrm{ch}}^{\mathrm{Heis}}(K3)=2,\qquad
\kappa_{\mathrm{ch}}^{\mathrm{Heis}}(E)=1,
\]
and therefore assigns \(3\) to their product specialization.
The morphism \(\eta_{K3,E}\) is the remaining realization theorem.
\item The Borcherds coordinate is defined by the constant Fourier
coefficient of the Jacobi input.  For Gritsenko's form,
\(c_1(0)=10\), so
\[
\kappa_{\mathrm{BKM}}(\Delta_5)=\frac{c_1(0)}2=5
\]
by the Borcherds--Gritsenko weight theorem
\cite{Borcherds1998,Gritsenko1999}.
\item The fibre coordinate is the rank of the Mukai lattice:
\[
\kappa_{\mathrm{fiber}}(K3)
=\operatorname{rk}\widetilde H(K3,\mathbb Z)
=b_0(K3)+b_2(K3)+b_4(K3)=1+22+1=24
\]
\cite{Mukai1987}.
\end{enumerate}
\end{proof}

The Mukai lattice supplies the proved arithmetic datum
\[
2c_+\bigl(\widetilde H(K3,\mathbb Z)\bigr)
=2\cdot4=8,
\]
obtained from its rank~$24$ and signature~$(4,20)$.  Its interpretation
as the chiral conductor $K^{\kappa_{\mathrm{ch}}}_{\mathbf B}$ has
status \ClaimStatusConjectured under
\[
H_{\mathsf B}
:=(H_{\mathrm{chart}},H_{\mathrm{KD}},H_{\mathrm{scalar}},
H_{\mathrm{mod}},H_{\mathrm{quantum}}),
\]
where the five hypotheses construct the indefinite-lattice chiral chart,
its Koszul partner, the scalar comparison, the modular comparison, and
the quantum-group specialization, respectively.
An exceptional RTT Yangian comparison begins instead with a
finite-dimensional representation \(V_0\), the coefficient space
\[
V_{\le N}
=\End(V_0)\otimes
\operatorname{span}\{u^{-1},\ldots,u^{-N-1}\},
\]
and the relation space
\(R_{\fg,\le N}=\operatorname{im}\Phi_{R_\hbar,\le N}\).
Its open construction consists of the PBW comparison
\[
p_{\fg,N}\colon
U\bigl(\fg[t]/t^{N+1}\fg[t]\bigr)\longrightarrow
\operatorname{gr}A^{\mathrm{RTT}}_{\fg,\le N},
\]
the inverse-relation morphism
\[
\rho_{\fg,N}\colon
R_{\fg,\le N}^{\perp}
\longrightarrow
\operatorname{im}\Phi_{R_\hbar^{-1},\le N},
\]
and the presentation morphism
\[
\delta_{\fg,N}\colon
A^{\mathrm{RTT}}_{\le N}(R_\hbar^{-1})
\longrightarrow
Y_{-\hbar}(\fg)_{\le N}.
\]
Exceptional Yangian duality requires these three morphisms to be
isomorphisms compatibly in \(N\).

\medskip

Each row has a fixed quadrant, presentation, Beilinson level, and
hypothesis package.  A CY row includes the data
\((X,\Sigma_{d-1},C)\) and the two-stage functor
\[
\Phi_d^{(\Sigma_{d-1},C)}
=\SpCh_{\Sigma_{d-1},C}\circ\PhiFA_d.
\]
For an Open witness \(A\), a CY coordinate begins with a
trace-compatible realization
\[
\iota_A\colon
A\xrightarrow{\ \sim\ }
\SpCh_{\Sigma_{d-1},C}
\bigl(\PhiFA_d(\mathcal C_X)\bigr).
\]

\begin{table}[ht]
\centering
\caption{The $5\times 5$ $\kappa$-stratification matrix.
The symbol \(\diamond_A\) denotes the open construction of a
trace-compatible CY realization \(\iota_A\) followed by the indicated
coordinate evaluation.  The five cells of row \(\mathbf B\) are
computed from \(K3\times E\).}
\label{tab:five-by-five-kappa-stratification}
\renewcommand{\arraystretch}{1.4}
{\small
\begin{tabular}{|c|l|c|c|c|c|c|}
\hline
\textbf{Row} & \textbf{Witness}
 & $\boldsymbol{\kappa_{\mathrm{cat}}}$
 & $\boldsymbol{\kappa^{\mathrm{Hodge}}_{\mathrm{ch}}}$
 & $\boldsymbol{\kappa^{\mathrm{Heis}}_{\mathrm{ch}}}$
 & $\boldsymbol{\kappa_{\mathrm{BKM}}}$
 & $\boldsymbol{\kappa_{\mathrm{fiber}}}$ \\
\hline
$\mathbf G$ & $\cH_1$ & \(\diamond_{\cH_1}\) &
\(\diamond_{\cH_1}\) & \(\diamond_{\cH_1}\) &
\(\diamond_{\cH_1}\) & \(\diamond_{\cH_1}\) \\
\hline
$\mathbf L$ & $V_k(\fg)$, generic $k$ & \(\diamond_{V_k(\fg)}\) &
\(\diamond_{V_k(\fg)}\) & \(\diamond_{V_k(\fg)}\) &
\(\diamond_{V_k(\fg)}\) & \(\diamond_{V_k(\fg)}\) \\
\hline
$\mathbf C$ & $\beta\gamma_{1/2}$ & \(\diamond_{\beta\gamma}\) &
\(\diamond_{\beta\gamma}\) & \(\diamond_{\beta\gamma}\) &
\(\diamond_{\beta\gamma}\) & \(\diamond_{\beta\gamma}\) \\
\hline
$\mathbf M$ & $\mathrm{Vir}_c$, $c(5c+22)\ne0$ &
\(\diamond_{\mathrm{Vir}_c}\) & \(\diamond_{\mathrm{Vir}_c}\) &
\(\diamond_{\mathrm{Vir}_c}\) & \(\diamond_{\mathrm{Vir}_c}\) &
\(\diamond_{\mathrm{Vir}_c}\) \\
\hline
$\mathbf B$ &
$\mathbf H_{\Delta_5}
=\Phi_3^{(\Sigma_2,C)}(K3\times E)$
 & $0$ & $0$ & $3$ & $5$ & $24$ \\
\hline
\end{tabular}}
\end{table}

The matrix contains twenty-five cells.  The five cells of row
\(\mathbf B\) are computed above; the twenty cells in the Open rows
are the realization problems \(\diamond_A\).  Intrinsic Open
invariants such as
\(\kappa_{\mathrm{mod}}(\cH_k)=k\) and
\(\kappa_{\mathrm{mod}}(V_k(\fg))
=\dim(\fg)(k+h^\vee)/(2h^\vee)\)
belong to the OPE table below.  A trace-compatible realization
\(\iota_A\) supplies the comparison with a CY coordinate.

Row \(\mathbf B\) has five coordinates and realizes the four-element
set of numerical values
\[
\{0,3,5,24\}.
\]
The two zeroes arise from the categorical Euler characteristic and
the Hodge supertrace.  The CY host, the transverse specialization,
and the Borcherds input distinguish row \(\mathbf B\) as a
mathematical object.  The structural Borcherds relation is
\[
\kappa_{\mathrm{BKM}}(\Phi_N)=\frac{c_N(0)}2
\qquad
\bigl(N\in\{1,2,3,4,6\}\bigr)
\]
\cite{Borcherds1998,Gritsenko1999}.  The arithmetic equality
\(5=3+2\) records the \(N=1\) values of three separately defined
functions.

\paragraph{Cross-volume Hodge-coordinate conflict.}
Volume~III defines
\[
\kappa_{\mathrm{ch}}^{\mathrm{Hodge}}
\bigl(\PhiFA_d(\Perf X)\bigr)=\Xi(X).
\]
Volume~I computes from an OPE and a genus-one trace a scalar which we
denote by \(\kappa_{\mathrm{mod}}(A)\).  The elliptic example separates
the two functions:
\[
\kappa_{\mathrm{ch}}^{\mathrm{Hodge}}
\bigl(\PhiFA_1(\Perf E)\bigr)=\Xi(E)=0,
\qquad
\kappa_{\mathrm{mod}}(\cH_1)=1.
\]
Their equality is a comparison theorem whose datum is a
trace-preserving equivalence
\[
\vartheta_{X,\Sigma,C}\colon
\SpCh_{\Sigma,C}\bigl(\PhiFA_d(\Perf X)\bigr)
\longrightarrow A
\]
together with equality of the two trace functionals.  This comparison
is open for the four Open witnesses in the matrix.  The notation
\(\kappa_{\mathrm{ch}}^{\mathrm{Hodge}}\) is therefore reserved for
\(\Xi(X)\), and \(\kappa_{\mathrm{mod}}\) denotes the OPE-derived
Volume~I scalar.  A bare \(\kappa(\cA)\) anywhere in this census
denotes the trace-datum scalar \(\kappa_{\mathrm{mod}}(\cA)\), never
the Hodge coordinate \(\kappa_{\mathrm{ch}}^{\mathrm{Hodge}}\).

\medskip

\noindent\textbf{Two-stage CY-derived constructions.}
The CY examples in Remark~\ref{rem:landscape-2} carry the functor
\[
 \Phi_d^{(\Sigma_{d-1},\,C)}
 \;=\;
 \SpCh_{\Sigma_{d-1},\,C}\;\circ\;\PhiFA_d
 \;\colon\;
 \mathrm{CY}_d\text{-}\mathrm{Cat}
 \;\xrightarrow{\;\PhiFA_d\;}\;
 \EdHolFA(X)
 \;\xrightarrow{\;\SpCh_{\Sigma_{d-1},\,C}\;}\;
 \EnHolFA(C)
\]
(cf.\ Vol.~III
Definitions~\ref{V3-def:phi-fa-and-sp}--\ref{V3-rem:multiple-e1-specialisations}).
The first functor assigns a holomorphic factorisation algebra on
\(X\); the second takes factorisation homology over the chosen
transverse manifold \(\Sigma_{d-1}\) and produces a chiral algebra on
\(C\).  Row \(\mathbf B\) uses
\[
(d,\Sigma_{d-1},C)=(3,\Sigma_2,C)
\]
and the target \(\mathbf H_{\Delta_5}\).  The complete input
\((\mathcal C_X,\Sigma_{d-1},C)\) determines the transverse manifold,
the curve, and the target chiral algebra.

\section*{Master table of computed invariants}
\label{sec:master-table}
\index{master table of invariants|textbf}

\paragraph{OPE modular characteristic.}
The fifth column of Table~\ref{tab:master-invariants} records
\(\kappa_{\mathrm{mod}}(\cA)\), the scalar extracted from the
genus-one trace of the OPE deformation complex.  The Hodge coordinate
\(\kappa_{\mathrm{ch}}^{\mathrm{Hodge}}\) remains the CY function
\(\Xi(X)\) defined above.  The genus-one identity is
\[
\operatorname{obs}^{\mathrm{sc}}_1(\cA)
=\kappa_{\mathrm{mod}}(\cA)\lambda_1.
\]
The uniform-weight continuation to genus \(g\) and the multi-weight
correction \(\delta F_g^{\mathrm{cross}}\) are governed by
Theorems~\ref{thm:genus-universality} and
\ref{thm:multi-weight-genus-expansion}.

For each row, the table records the specified dual or partner, the
central charge, the conformal complementarity sum on its domain, and
\(\kappa_{\mathrm{mod}}\).  A CJ or HF entry names a conjectural or
open construction.  Set
\[
\begin{aligned}
K^c(\cA,\cA^!)&:=c(\cA)+c(\cA^!),\\
K^{\kappa_{\mathrm{mod}}}(\cA,\cA^!)
&:=\kappa_{\mathrm{mod}}(\cA)
+\kappa_{\mathrm{mod}}(\cA^!).
\end{aligned}
\]
These are numerical trace functions on the displayed Verdier
parameter loci.  The averaging morphism
\(\Av_\cA\), the derived centre \(\DerZ(\cA)\), and a physical bulk
algebra are three additional typed objects.  The K3 Mukai lattice has
rank~$24$, signature~$(4,20)$, and hence the exact arithmetic value
\[
2c_+\bigl(\widetilde H(K3,\mathbb Z)\bigr)=8.
\]
The chiral identification
$K^{\kappa_{\mathrm{ch}}}_{\mathbf B}=8$ has status
\ClaimStatusConjectured under the package~$H_{\mathsf B}$ displayed
above.  Independently, the Borcherds construction gives
\(\kappa_{\mathrm{BKM}}(\Delta_5)=5\).
Throughout, \(d=\dim\fg\), \(r=\operatorname{rank}\fg\),
\(h^\vee\) is the dual Coxeter number, and \(t=k+h^\vee\).

\begin{remark}[Standard-family constants and collision conventions]
On the standard landscape the modular characteristic is normalized by
the genus-$1$ obstruction coefficient and by the trace-form collision
residue:
\[
\begin{aligned}
\kappa_{\mathrm{mod}}(\mathcal H_{\ell}^{\oplus d}) &= d\,\ell,
&
\kappa_{\mathrm{mod}}(V_k(\fg)) &= \frac{\dim(\fg)(k+h^\vee)}{2h^\vee},
&
\kappa_{\mathrm{mod}}(\mathrm{Vir}_c) &= \frac{c}{2}.
\end{aligned}
\]
On the principal modular-companion surface
\(H_{\mathrm{diag}}^{g=1}+H_{W_N}^{\mathrm{DS/bar}}\),
\[
 \kappa_{\mathrm{mod}}(\mathcal W_N)
 =c\sum_{j=2}^{N}\frac1j
 =c(H_N-1),
 \qquad
 H_N=\sum_{j=1}^{N}\frac1j.
\]
Thus the rank-one Heisenberg algebra satisfies
\(\kappa_{\mathrm{mod}}(\mathcal H_k)=k\), while the rank-\(d\)
level-one free-boson algebra satisfies
\(\kappa_{\mathrm{mod}}(\mathcal H_1^{\oplus d})=d\).
For a simple Lie algebra, \(h^\vee>0\), and the affine formula has the
boundary values
\[
\kappa_{\mathrm{mod}}(V_0(\fg))=\frac{\dim(\fg)}{2},\qquad
\kappa_{\mathrm{mod}}(V_{-h^\vee}(\fg))=0.
\]
The degree-two coefficient depends on a chosen one-dimensional
subspace of fields.  For
\[
L_T=\mathbb C T,\qquad L_J=\mathbb C J,
\]
the restrictions are
\[
S_2|_{L_T}=S_2^T=\frac c2,\qquad
S_2|_{L_J}=S_2^J=k.
\]
A multi-generator family carries a specified projection
\(p_L:\Def_{\mathrm{cyc}}(\cA)\to L\), while
\(\kappa_{\mathrm{mod}}(\cA)\) is the scalar genus-one trace after
the stated averaging map.

On the Virasoro parameter locus
\[
U_{\mathrm{Vir}}
=\{c\in\mathbb C\mid c(5c+22)\ne0\},
\]
\[
S_2=\frac c2,\qquad
S_3=2,\qquad
S_4^{\mathrm{Gram}}=\frac{10}{c(5c+22)},\qquad
S_5^{\mathrm{Ricc}}=-\frac{48}{c^2(5c+22)},
\]
and
\[
\Delta^{\mathrm{Ricc}}=8\kappa_{\mathrm{mod}} S_4^{\mathrm{Gram}}
=\frac{40}{5c+22},\qquad
\Lambda={:}TT{:}-\frac{3}{10}\partial^2T,\qquad
\langle\Lambda|\Lambda\rangle=\frac{c(5c+22)}{10}.
\]
The first three entries are presentation data and the fourth is the
weighted-Riccati output.  A scalar extracted from the connected Ward
correlator carries the additional datum
$\mathsf H_{\mathrm{res}}(\Vir_c;X)$ and its normalized ordered
residue map.
The collision residues used in this census are
\[
r_{\mathcal H}^{\mathrm{coll}}(z)=\frac{k}{z},\qquad
r_{V_k(\fg)}^{\mathrm{coll}}(z)=\frac{k\Omega_{\mathrm{tr}}}{z},
\qquad
r_{\mathrm{Vir}_c}^{\mathrm{coll}}(z)=\frac{c/2}{z^3}+\frac{2T}{z}.
\]
The affine KZ presentation
$r_{\mathrm{KZ}}(z)=\Omega/((k+h^\vee)z)$, with
$\Omega=2h^\vee\Omega_{\mathrm{tr}}$, and the collision presentation
form the ordered pair
\[
\left(
r^{\mathrm{coll}}_{V_k(\fg)}(z),
r^{\mathrm{KZ}}(z)
\right)
=
\left(
\frac{k\Omega_{\mathrm{tr}}}{z},
\frac{\Omega}{(k+h^\vee)z}
\right).
\]
At \(k=0\) their values are
\[
\left(0,\frac{\Omega}{h^\vee z}\right),
\]
which records the trace-form and KZ normalizations explicitly.
\smallskip

The Heisenberg identity follows from the OPE
$J^a(z)J^b(w)\sim \ell\,\delta^{ab}(z-w)^{-2}$. The affine identity is
the Sugawara two-channel computation: the trace-form double-pole
projection contributes $\dim(\fg)k/(2h^\vee)$ and the simple-pole Lie
bracket channel contributes $\dim(\fg)/2$. The Virasoro identity is
the leading coefficient of
$T(z)T(w)\sim (c/2)(z-w)^{-4}+2T(w)(z-w)^{-2}+\partial T(w)(z-w)^{-1}$.
For principal $\mathcal W_N$, the Drinfeld--Sokolov strong-generator
weights are $2,\ldots,N$.  On
$H_{\mathrm{diag}}^{g=1}+H_{W_N}^{\mathrm{DS/bar}}$, their diagonal
trace gives the harmonic sum $\sum_{j=2}^{N}1/j$; the specialization
$N=2$ recovers Virasoro, while $N=3$ gives
\(\kappa_{\mathrm{mod}}(\mathcal W_3)=5c/6\) on
$H_{\mathrm{diag}}^{g=1}+H_{W_3}^{\mathrm{DS/bar}}$.
The quartic coefficient and the norm of \(\Lambda\) are computed from
the level-four Virasoro Gram matrix in
Theorem~\ref{thm:nms-virasoro-quartic-explicit}.  The displayed
coefficients are collected in
Proposition~\ref{prop:virasoro-shadow-canonical}.  The trace-form/KZ
comparison is Theorem~\ref{thm:q-convention-bridge-main}.
\end{remark}

\begin{table}[ht]
\centering
\caption{Computed invariants and conjectural rows for chiral algebras}
\label{tab:master-invariants}
\index{Koszul duality!master table}
\renewcommand{\arraystretch}{1.5}
{\small
\resizebox{\textwidth}{!}{%
\begin{tabular}{|l|l|c|c|c|c|}
\hline
\textbf{Algebra $\cA$} & \textbf{Koszul Dual / partner datum}
 & $\boldsymbol{c(\cA)}$ & $\boldsymbol{c + c'}$
 & $\boldsymbol{\kappa_{\mathrm{mod}}(\cA)}$ & \textbf{Status} \\
\hline
\multicolumn{6}{|c|}{\textit{Free Fields}} \\
\hline
Free fermion $\psi$
 & $\mathrm{Sym}^{\mathrm{ch}}(\gamma)$
 & $\tfrac{1}{2}$ & $0$ & $\tfrac{1}{4}$ & PH \\
\hline
$bc$ ghosts (weight $\lambda$)
 & $\beta\gamma$ system (weight $\lambda$)
 & $1 - 3(2\lambda{-}1)^2$ & $0$
 & $c/2$ & PH \\
\hline
\multicolumn{6}{|c|}{\textit{Heisenberg \textup{(}curved abelian\textup{)}}} \\
\hline
Heisenberg $\mathcal{H}_\ell$
 & $\mathrm{Sym}^{\mathrm{ch}}(V^*)$ (curved)
 & $1$ & n/a\textsuperscript{$\dagger$} & $\ell$ & PH \\
\hline
\multicolumn{6}{|c|}{\textit{Affine Kac--Moody
 \textup{(}non-critical universal locus; Feigin--Frenkel level shift $k \mapsto -k - 2h^\vee$\textup{)}}} \\
\hline
$\widehat{\fg}_k$ (general)
 & $\widehat{\fg}_{-k-2h^\vee}$
 & $\dfrac{kd}{k + h^\vee}$ & $2d$
 & $\dfrac{td}{2h^\vee}$ & PH \\[4pt]
\hline
$\widehat{\mathfrak{sl}}_2$ at level $k$
 & $\widehat{\mathfrak{sl}}_2$ at $-k{-}4$
 & $\dfrac{3k}{k+2}$ & $6$ & $\dfrac{3(k+2)}{4}$ & PH \\[4pt]
\hline
$\widehat{\mathfrak{sl}}_3$ at level $k$
 & $\widehat{\mathfrak{sl}}_3$ at $-k{-}6$
 & $\dfrac{8k}{k+3}$ & $16$ & $\dfrac{4(k+3)}{3}$ & PH \\[4pt]
\hline
$\widehat{\mathfrak{so}}_{8}$ at level $k$ (type $D_4$; triality)
 & $\widehat{\mathfrak{so}}_{8}$ at $-k{-}12$
 & $\dfrac{28k}{k+6}$ & $56$ & $\dfrac{7(k+6)}{3}$ & PH \\[4pt]
\hline
$\widehat{E}_6$ at level $k$
 & $\widehat{E}_6$ at $-k{-}24$
 & $\dfrac{78k}{k+12}$ & $156$ & $\dfrac{13(k+12)}{4}$ & PH \\[4pt]
\hline
$\widehat{E}_7$ at level $k$
 & $\widehat{E}_7$ at $-k{-}36$
 & $\dfrac{133k}{k+18}$ & $266$ & $\dfrac{133(k+18)}{36}$ & PH \\[4pt]
\hline
$\widehat{E}_8$ at level $k$
 & $\widehat{E}_8$ at $-k{-}60$
 & $\dfrac{248k}{k+30}$ & $496$ & $\dfrac{62(k+30)}{15}$ & PH \\[4pt]
\hline
\multicolumn{6}{|c|}{\textit{Affine Kac--Moody \textup{(}non-critical non-simply-laced, $h \neq h^\vee$\textup{)}}} \\
\hline
$\widehat{\mathfrak{so}}_{5,k}$ (type $B_2$)
 & $\widehat{\mathfrak{so}}_{5,-k-6}$
 & $\dfrac{10k}{k+3}$ & $20$ & $\dfrac{5(k+3)}{3}$ & PH \\[4pt]
\hline
$\widehat{\mathfrak{sp}}_{4,k}$ (type $C_2$)
 & $\widehat{\mathfrak{sp}}_{4,-k-6}$
 & $\dfrac{10k}{k+3}$ & $20$ & $\dfrac{5(k+3)}{3}$ & PH \\[4pt]
\hline
$\widehat{G}_{2,k}$
 & $\widehat{G}_{2,-k-8}$
 & $\dfrac{14k}{k+4}$ & $28$ & $\dfrac{7(k+4)}{4}$ & PH \\[4pt]
\hline
$\widehat{F}_{4,k}$
 & $\widehat{F}_{4,-k-18}$
 & $\dfrac{52k}{k+9}$ & $104$ & $\dfrac{26(k+9)}{9}$ & PH \\[4pt]
\hline
\multicolumn{6}{|c|}{\textit{$\mathcal{W}$-Algebras
 \textup{(}Drinfeld--Sokolov reduction, $k' = -k - 2h^\vee$\textup{)}}} \\
\hline
$\mathrm{Vir}_c = \mathcal{W}^k(\mathfrak{sl}_2)$
 & $\mathrm{Vir}_{26-c}$\textsuperscript{$\ddagger$}
 & $1 - \dfrac{6(k{+}1)^2}{k{+}2}$ & $26$ & $c/2$ & PH \\[4pt]
\hline
$\mathcal{W}_3^k(\mathfrak{sl}_3)$
 & $\mathcal{W}_3^{k'}(\mathfrak{sl}_3)$ on
 $H_{W_3}^{\mathrm{DS/bar}}$
 & $2 - \dfrac{24(k{+}2)^2}{k{+}3}$ & $100$\textsuperscript{$c$}
 & $\tfrac{5c}{6}$\textsuperscript{$\mathrm{mod/comp}$}
 & PH$^c$/CD$^{W_3}$ \\[4pt]
\hline
$\mathcal{W}_N^k(\mathfrak{sl}_N)$ (general)
 & $\mathcal{W}_N^{k'}(\mathfrak{sl}_N)$ on
 $H_{W_N}^{\mathrm{DS/bar}}$
 & $c(k)$ & $2(N{-}1)(2N^2{+}2N{+}1)$
 \textsuperscript{$c$}
 & $c \displaystyle\sum_{j=2}^{N}\tfrac{1}{j}$
 \textsuperscript{$\mathrm{mod/comp}$}
 & PH$^c$/CD$^{W_N}$ \\[6pt]
\hline
\multicolumn{6}{|c|}{\textit{Non-principal $\mathcal{W}$-algebras
 \textup{(}subregular DS reduction\textup{)}}} \\
\hline
$\mathrm{BP}_k = \mathcal{W}_3^{(2)}(k)$
 & $\mathrm{BP}_{-k-6}$ on $H_{\mathrm{BP}}^{\mathrm{DS/bar}}$
 & $-\dfrac{(2k{+}3)(3k{+}1)}{k{+}3}$ & $50$
 & open: direct genus-$1$ curvature
 & PH$^{\mathrm{BP},c}$/OP$^{\mathrm{BP},g=1}$/CD$^{\mathrm{BP,DS}}$ \\[4pt]
\hline
\multicolumn{6}{|c|}{\textit{Lattice VOAs and Exceptional}} \\
\hline
Leech lattice $V_{\Lambda_{24}}$
 & $V_{\Lambda_{24}}$ (self-dual\textsuperscript{$\|$})
 & $24$ & n/a\textsuperscript{$\|$}
 & $24$ & PH \\
\hline
Monster $V^\natural$ (FLM)
 & $\mathrm{Vir}_2$\textsuperscript{$\S$}
 & $24$ & $26$\textsuperscript{$\S$}
 & $12$ & PH \\
\hline
\multicolumn{6}{|c|}{\textit{$\Eone$-Chiral Algebras}} \\
\hline
Yangian $Y(\fg)$
 & $Y_{R^{-1}}(\fg)$
 & n/a & n/a & n/a & PH \\
\hline
Quantum lattice $\Vlat_\Lambda^{N,q}$
 & $(\Vlat_\Lambda^{N,-q})^c$
 & $d = \operatorname{rank}(\Lambda)$ & n/a
 & $\operatorname{rank}(\Lambda)$ & PH \\
\hline
\multicolumn{6}{|c|}{\textit{Honest frontier rows
 \textup{(}see Remark~\textup{\ref{rem:census-silent-non-coverage}}\textup{)}}} \\
\hline
$\mathcal{N}{=}2$ $\mathrm{SCA}_c$
 \textup{(}$c = 3k/(k{+}2)$, $k \ne -2$\textup{)}
 & $\mathrm{SCA}_{6-c}$ \textup{(}Ito--Tanaka involution\textup{)}
 & $\dfrac{3k}{k+2}$ & $6$
 & $\dfrac{k+4}{4}$ & PH$^{\mathrm{N2}}$ \\[4pt]
\hline
Triplet $\mathcal{W}(p)$ \textup{(}$p \ge 2$\textup{)}
 & $T$-channel self-pair; full logarithmic partner pending
 & $1 - \dfrac{6(p{-}1)^2}{p}$ & pending \textup{(}$T$-channel $2c$\textup{)}
 & $c/2$ on the $T$-channel & CJ$^{\mathrm{log}}$ \\[4pt]
\hline
Coset $\cA/\cB$
 \textup{(}Kazama--Suzuki, GKO, parafermion\textup{)}
 & Koszul partner pending
 & $c(\cA) - c(\cB)$ & pending
 & pending & HF \\[4pt]
\hline
Non-rational / indefinite lattice $V_L$
 \textup{(}$L$ not positive-definite\textup{)}
 & specified lattice partner pending
 & $\operatorname{sig}(L)=(p,q)$; positive cone and duality data required
 & pending
 & pending & HF \\[4pt]
\hline
Universal $V_k(\fg)$ restricted to $k + h^\vee = p/q$
 \textup{(}$\gcd(p,qh^\vee) = 1$, $p,q \ge 1$\textup{)}
 & $V_{-k-2h^\vee}(\fg)$ on the reflected universal locus
 & $\dfrac{kd}{k + h^\vee}$ & $2d$
 & $\dfrac{td}{2h^\vee}$ \textup{(}restriction\textup{)}
 & CJ$^{\mathrm{adm}}$ \\[4pt]
\hline
\end{tabular}%
}% end resizebox
}% end small
\end{table}

\noindent\textbf{Status column}: PH = proved in this monograph
(\ClaimStatusProvedHere); CD = conditional on the displayed hypothesis
package; CJ = conjectured (\ClaimStatusConjectured);
OP = open (\ClaimStatusOpen); HF = frontier (invariant not yet closed).
\begin{itemize}[leftmargin=*]
 \item PH$^c$/CD$^{W_N}$: the central conductor
 $c_N(k)+c_N(-k-2N)=4N^3-2N-2$ is an exact rational-function
 identity.  The displayed principal modular characteristic and its
 companion sum are statements on
 $H_{\mathrm{diag}}^{g=1}+H_{W_N}^{\mathrm{DS/bar}}$; at $N=3$ this
 package is $H_{\mathrm{diag}}^{g=1}+H_{W_3}^{\mathrm{DS/bar}}$, the
 exact central value is $100$, and the conditional modular sum is
 $250/3$.
 \item PH$^{\mathrm{BP},c}$/OP$^{\mathrm{BP},g=1}$/CD$^{\mathrm{BP,DS}}$:
 the standard Feigin--Semikhatov central charge and its companion
 sum~$50$ are proved central-charge identities.  Fehily--Kawasetsu--Ridout
 define BP as an ordinary bosonic vertex algebra with even strong
 generators $J,G^+,G^-,T$~\cite[\S2.1]{FKR20}.  The reciprocal-weight
 diagnostic is therefore
 \[
  1+\frac23+\frac23+\frac12=\frac{17}{6}.
 \]
 The BP genus-$1$ curvature computation determines
 $\kappa_{\mathrm{BP}}$, $\varrho_{\mathrm{BP}}$, and
 $K^\kappa_{\mathrm{BP}}$; their ledger status is
 \texttt{open-genus-one-computation} (\ClaimStatusOpen).  The same-family
 Verdier--Koszul comparison is carried by
 $H_{\mathrm{BP}}^{\mathrm{DS/bar}}$: a regular level $k\ne-3$,
 a finite-type bar dual or a completed continuous nondegenerate
 Verdier pairing, and subregular hook-type DS/bar transport.
 \item PH$^{\mathrm{N2}}$: the Kazama--Suzuki parameter
 $c = 3k/(k{+}2)$ satisfies the exact reflected identity
 $c(k)+c(-k-4)=6$.  Proposition~\ref{prop:n2-kappa} gives the
 conditional implication
 $\kappa=(k+4)/4$ under
 $H_{\mathcal N=2}^{\mathrm{mod}}$; Proposition~\ref{prop:n2-complementarity}
 gives $K^\kappa=1$ after adjoining
 $H_{\mathcal N=2}^{\mathrm{KD}}$.  The proved datum here is the
 OPE/coset parameter identity; the modular conductor remains an open
 comparison problem.
 \item CJ$^{\mathrm{log}}$: The $T$-channel is class~$M$ with
 $c = 1 - 6(p-1)^2/p$ \textup{(}$p = 2$: $c = -2$;
 $p = 3$: $c = -7$\textup{)} and $\kappa_T = c/2$ via
 Proposition~\ref{prop:wp-kappa}. The full logarithmic Koszul partner
 and conductor are not supplied by this $T$-channel scalar. The
 strong-generator count
 is $4 = 1\,(T) + 3\,(W^{\pm}, W^{0})$. Tempering of
 $\cW(p)$ is conjectural per
 Remark~\ref{rem:census-silent-non-coverage};
 Gurarie (1993) and Flohr (1996) construct unbounded
 Massey amplitudes, so the Adamovi\'c--Milas
 character-amplitude bound remains a frontier.
 \item HF (cosets): $c(\cA/\cB) = c(\cA) - c(\cB)$ is
 Goddard--Kent--Olive additivity; Kazama--Suzuki gives the
 $\mathcal{N}{=}2$ coset family. This central-charge subtraction does
 not determine the Koszul partner or the modular characteristic:
 both require family-by-family bar and commutant computations. See
 Remark~\ref{rem:census-silent-non-coverage}, item~3.
 \item HF (indefinite lattice): The Gaussian-sector formulas above
 have positive-definite lattice VOAs as their theorem domain.  A
 Lorentzian lattice row requires a chosen positive cone and a specified
 self-duality convention.  The Mukai row below records the proved
 arithmetic $c_+=4$ and $2c_+=8$; the formulas
 $\kappa=c_+$ and $K^\kappa=2c_+$ have status
 \ClaimStatusConjectured under~$H_{\mathsf B}$. The unimodular
 self-duality theorem (Theorem~\ref{thm:lattice:unimodular-self-dual})
 supplies the positive-definite input; $H_{\mathsf B}$ supplies the
 Lorentzian extension package.
 Canonical examples: $\mathrm{II}_{1,1}$ (Narain slab),
 $\mathrm{II}_{25,1}$ (fake-monster root lattice).
 \item CJ$^{\mathrm{adm}}$: The displayed central charge and
 $\kappa$ are the universal affine formulas restricted to an
 admissible parameter and paired with the reflected universal level.
 They are not a statement about the simple quotient $L_k(\fg)$.
 The periodic-CDG structure
 (Theorem~\ref{thm:periodic-cdg-is-koszul-compatible}) supplies the
 compatibility framework; finite-length $C_2$-cofinite admissible
 minimal models (Corollary~\ref{cor:class-M-admissible-minimal-model})
 require separate simple-quotient bar computations.
\end{itemize}

\begin{remark}[Bershadsky--Polyakov conformal-vector conventions]
\label{rem:bp-census-convention-separation}
The canonical census lane uses
\[
c_{\mathrm{BP}}(k)
=-\frac{(2k+3)(3k+1)}{k+3},
\qquad
c_{\mathrm{BP}}(k)+c_{\mathrm{BP}}(-k-6)=50.
\]
The explicitly shifted secondary lane is the rational function
\[
c_{\mathrm{BP}}^{\mathrm{shift}}(k)
=2-\frac{24(k+1)^2}{k+3},
\qquad
c_{\mathrm{BP}}^{\mathrm{shift}}(k)
+c_{\mathrm{BP}}^{\mathrm{shift}}(-k-6)=196.
\]
Both sums are identities in~$\mathbb Q(k)$, and $k=-3$ is their common
pole and the fixed point of $k\mapsto-k-6$.  The first lane supplies the
exact central-charge conductor $K^c_{\mathrm{BP}}=50$.  In the FKR
bosonic convention, $G^+$ and $G^-$ are even
\textup{(}\cite[\S2.1]{FKR20}\textup{)}, and the reciprocal-weight
diagnostic is $17/6$.  The former proposal
\[
 \kappa_{\mathrm{BP}}=c_{\mathrm{BP}}/6,
 \qquad
 \varrho_{\mathrm{BP}}=1/6,
 \qquad
 K^\kappa_{\mathrm{BP}}=25/3
\]
has status \ClaimStatusConjectured.  A direct non-separating genus-$1$
curvature computation of the BP modular trace is its proof obligation;
the active ledger tag is \texttt{open-genus-one-computation}.
The value~$196$ belongs to the computed-secondary shifted scalar lane.
Algebra-level same-family transport is the conditional comparison on
$H_{\mathrm{BP}}^{\mathrm{DS/bar}}$ stated above.
\end{remark}

\noindent The Schellekens $71$ rows at $c = 24$ form a closed
classification surface:
Theorem~\ref{thm:schellekens-structured-subset} identifies the
structured decomposition
\[
71 = 24_{\mathrm{Niemeier}} + 1_{V_1=0} + 46_{\mathrm{nonlat}}
\]
as a counting corollary of an explicit partition of Schellekens's
numbered list. When a concrete representative of the $V_1=0$ row is
needed we use $V^\natural$, but the row statement itself does not
assume moonshine uniqueness.
This master table still records only the \emph{coarse} $\kappa/r_{\max}$
rows represented by the Niemeier and Monster entries.  The remaining
$46$ non-lattice cases occupy distinct positions at this granularity;
they are organized by their Schellekens slot and
the weight-one threshold of
Proposition~\ref{prop:schellekens-weight-two-threshold}.  These two
coordinates retain the distinctions among the non-lattice rows.

\medskip

All PH dualities in this table have complete proofs at the semantic
level recorded here. CJ and HF rows record the specified scalar or
frontier datum explicitly displayed in their row and no more.

\begin{remark}[Euler--Koszul class]\label{rem:census-euler-koszul}
The Euler--Koszul class $\operatorname{ek}(\cA)$
(Definition~\ref{def:euler-koszul-tier}) assigns to each family:
$\operatorname{ek} = 0$ for Heisenberg, affine $V_k(\fg)$,
$\beta\gamma$, and lattice~$V_\Lambda$
(exact Euler--Koszul); $\operatorname{ek} = 1$ for $\mathrm{Vir}_c$;
$\operatorname{ek} = N{-}1$ for $\cW_N$.
This is independent of shadow depth~$\kappa_d$
(Theorem~\ref{thm:shadow-euler-independence}).
\end{remark}

\medskip

\noindent\textbf{The shadow obstruction tower as organizing frame.}
Every completed vertex-algebra entry in Table~\ref{tab:master-invariants}
is a projection of the universal MC element
$\Theta_\cA \in \MC(\gAmod)$
(Theorem~\ref{thm:mc2-bar-intrinsic}). The scalar $\kappa(\cA)$ is
the degree-$2$ projection; the complementarity sum $c + c'$ is the
ambient genus-$0$ shadow when a conformal partner is defined; the
Koszul dual is obtained via Verdier duality on the bar coalgebra.
The $\Eone$ and frontier rows are specified separately and are not read as
conformal vertex-algebra complementarity assertions.
The seed datum for each family is its modular Koszul triple
$\mathfrak{T} = (\cA, \cA^!, r(z))$
(Definition~\ref{def:modular-koszul-triple}):
all scalar normalizations in this table use the conventions of
Appendix~\ref{sec:scalar-normalization-appendix}.
\begin{equation}\label{eq:census-triples}
\begin{aligned}
\mathfrak{T}_{\cH}
 &= \bigl(\cH_\kappa,\;
 (\operatorname{Sym}^{\mathrm{ch}}(V^*[1]),m_0=-\kappa\omega),\;
 \kappa/z\bigr), &
\mathfrak{T}_{\cF}
 &= \bigl(\cF,\;
 \operatorname{Sym}^{\mathrm{ch}}(\gamma),\;
 0\bigr), \\
% Convention: trace-form collision residue. The KZ presentation is
% Omega/((k+h^\vee)z), with Omega = 2h^\vee Omega_tr. These are
% related by the trace/KZ bridge, not by equality of rational
% functions; only the trace-form residue has the level prefix k.
\mathfrak{T}_{\widehat{\fg}_k}
 &= \bigl(\widehat{\fg}_k,\;
 \widehat{\fg}_{-k-2h^\vee},\;
 k\Omega_{\mathrm{tr}}/z\bigr), &
\mathfrak{T}_{\beta\gamma}
 &= \bigl(\beta\gamma_\lambda,\;
 bc_\lambda,\;
 0\bigr), \\
\mathfrak{T}_{\mathrm{Vir}}
 &= \bigl(\mathrm{Vir}_c,\;
 \mathrm{Vir}_{26-c},\;
 \tfrac{c}{2z^3} + \tfrac{2T}{z}\bigr), &
\mathfrak{T}_{\mathcal{W}_3}
 &= \bigl(\mathcal{W}_3^k,\;
 \mathcal{W}_3^{-k-6},\;
 r_{\mathcal{W}_3}(z)\bigr).
\end{aligned}
\end{equation}
Table~\ref{tab:shadow-tower-census} below records the
shadow obstruction tower data that completes the portrait: the archetype class,
shadow depth $r_{\max}$, and the first nontrivial higher shadow for
each family.

\begin{proposition}[Minimal-model class transport under residue descent]
\label{prop:minimal-model-class-transport}
\ClaimStatusConditional
The simple quotient $\mathrm{Vir}_{c_{p,q}}$ at minimal-model central
charge $c_{p,q} = 1 - 6(p-q)^{2}/(pq)$ carries the universal stress
tensor and its exact level-four Gram form.  Assume that the quotient map
extends to the ordered residue complexes, intertwines their contractions,
and transports the quadratic recurrence with the same normalization.
For $c_{p,q}(5c_{p,q}+22)\ne0$, the transported recurrence has
$\Delta=40/(5c_{p,q}+22)$ and hence the class-M value
$r_{\max}=\infty$.  The residue-descent hypothesis is the load-bearing
statement about the null-vector ideal.
\end{proposition}

\begin{proof}
The assumed chain map transports the seeds
$\kappa=c_{p,q}/2$, $S_3=2$, and
$S_4=10/[c_{p,q}(5c_{p,q}+22)]$, together with the transferred
quadratic brackets.  Thus
$\Delta=8\kappa S_4=40/(5c_{p,q}+22)$.  The stated regularity gives
$\Delta\ne0$, and the square-root recurrence has two simple branch
points.  The class-M branch of the recurrence trichotomy gives
$r_{\max}=\infty$.
\end{proof}

\begin{proposition}[Virasoro Gram seed and weighted-Riccati value]
\label{prop:virasoro-shadow-canonical}
\ClaimStatusConditional
On the Gram-regular locus $c(5c+22)\ne0$, the stress-tensor
presentation supplies
\begin{align*}
S_2 &= \kappa(\mathrm{Vir}_c) = c/2, \\
S_3 &= 2, \\
S_4^{\mathrm{Gram}} &= \frac{10}{c(5c+22)}, \\
S_5^{\mathrm{Ricc}} &= \frac{-48}{c^2(5c+22)}, \\
\Delta^{\mathrm{Ricc}} &= 8\kappa S_4^{\mathrm{Gram}}
= \frac{40}{5c+22}.
\end{align*}
The Zamolodchikov norm of the first composite primary is
$\langle\Lambda | \Lambda\rangle = c(5c+22)/10$ with
$\Lambda = {:}TT{:} - (3/10)\partial^2 T$.  The first two lines are
Virasoro structure constants, the third is the inverse Gram seed, and
the fourth is the output of the weighted-Riccati recurrence.  An ordered
residue scalar $S_5(\Vir_c;\mathsf H_{\mathrm{res}})$ is obtained after
constructing $\mathsf H_{\mathrm{res}}(\Vir_c;X)$ and its normalized
projection; agreement with $S_5^{\mathrm{Ricc}}$ is the corresponding
comparison theorem.
\end{proposition}

\begin{proof}
In the vacuum basis
$\{L_{-4}\mathbf1,L_{-2}^2\mathbf1\}$ the level-four Gram matrix is
\[
 \begin{pmatrix}5c&3c\\3c&c(c+8)/2\end{pmatrix}.
\]
The quasi-primary vector
$\Lambda=L_{-2}^2\mathbf1-(3/5)L_{-4}\mathbf1$ therefore has norm
$c(5c+22)/10$, giving $S_4^{\mathrm{Gram}}$.  At arity five the
weighted recurrence reads
\[
 S_5^{\mathrm{Ricc}}
 =-\frac{6}{5\kappa}S_3S_4^{\mathrm{Gram}},
\]
which gives the displayed rational function after substituting
$\kappa=c/2$ and $S_3=2$.  The formula for
$\Delta^{\mathrm{Ricc}}$ is the same substitution in
$8\kappa S_4^{\mathrm{Gram}}$.
\end{proof}

\begin{table}[ht]
\centering
\caption{Shadow obstruction tower census: each family as a projection of $\Theta_\cA$}
\label{tab:shadow-tower-census}
\index{shadow obstruction tower!census table}
\renewcommand{\arraystretch}{1.5}
{\small
\resizebox{\textwidth}{!}{%
\begin{tabular}{|l|c|c|c|l|}
\hline
\textbf{Algebra $\cA$}
 & \textbf{Archetype}
 & $\boldsymbol{r_{\max}}$
 & \textbf{First higher shadow}
 & \textbf{Reference} \\
\hline
\multicolumn{5}{|c|}{\textit{Gaussian: tower terminates at scalar level}} \\
\hline
Free fermion $\psi$ & G & 2 & none & \S\ref{chap:free-fields} \\
\hline
Heisenberg $\mathcal{H}_\kappa$ & G & 2 & none & \S\ref{ch:heisenberg-frame} \\
\hline
Lattice VOA $V_\Lambda$ & G & 2 & none & \S\ref{sec:lattice-genera} \\
\hline
Niemeier $V_\Lambda$ (all 24; $\kappa = 24$) & G & 2 & none: blind to root system & Thm~\ref{thm:lattice:niemeier-shadow-universality}; Thm~\ref{thm:schellekens-structured-subset} \\
\hline
$Y_{1,1,1}[\Psi]$ (corner VOA, $c = 0$) & G & 2 & none: cubic vanishes ($J$ abelian) & Rem~\ref{rem:y111-class-g} \\
\hline
\multicolumn{5}{|c|}{\textit{Lie/tree: cubic shadow nontrivial}} \\
\hline
$\widehat{\fg}_k$ (simply-laced: $A_n, D_n, E_{6,7,8}$) & L & 3 & $C$: Lie bracket & \S\ref{chap:kac-moody-koszul} \\
\hline
$\widehat{\fg}_k$ (non-simply-laced: $B_N, C_N, G_2, F_4$) & L & 3 & $C$: Lie bracket & \S\ref{sec:B2-details}, \S\ref{sec:G2-details} \\
\hline
\multicolumn{5}{|c|}{\textit{Contact/quartic: quartic contact shadow nontrivial}} \\
\hline
$\beta\gamma$ system & C & 4 & $Q^{\mathrm{contact}}$; $\mu_{\beta\gamma}=0$ & \S\ref{chap:beta-gamma} \\
\hline
$bc$ ghosts & C & 4 & $Q^{\mathrm{contact}}$; standard conformal-weight family & \S\ref{chap:free-fields} \\
\hline
Symplectic boson ($\beta\gamma$ at $\lambda=1/2$, $c=-1$) & C & 4 & $Q^{\mathrm{contact}}$ in $\mathbb{Z}_2$-graded bar & \S\ref{sec:symplectic-boson-bar} \\
\hline
Symplectic fermion $\mathcal{SF}=bc_1$ ($c=-2$, $\kappa=-1$) & C & 4 & $Q^{\mathrm{contact}}$ on the $bc_\lambda$ family; logarithmic modules do not obstruct parent bar & \S\ref{chap:free-fields} \\
\hline
\multicolumn{5}{|c|}{\textit{Mixed/infinite: tower does not terminate}} \\
\hline
$\mathrm{Vir}_c$ & M & $\infty$ & $C$: cubic + $Q$: quartic & \S\ref{chap:w-algebra-koszul} \\
\hline
$\mathcal{W}_N$ (finite $N$) & M & $\infty$ & $C$+$Q$; quintic forced & \S\ref{chap:w-algebra-koszul} \\
\hline
$\mathcal{W}_3^k(\mathfrak{sl}_3)$ (principal) & M (candidate) & open & $N_\Lambda=c(5c{+}22)/10$; $WW\to\Lambda$ coupling $32/(22{+}5c)$; scalar quartic under $H_{\mathrm{bar}}+H_{\mathrm{proj}}$ & Thm~\ref{thm:w3-w-line-s4-zamolodchikov} \\
\hline
Ising minimal model $\mathcal{M}(4,3)$ ($c=1/2$) & M & $\infty$ & $S_4 = 40/49$; $\Delta = 80/49 \neq 0$ & Prop~\ref{prop:ising-shadow-invariants} \\
\hline
$\mathcal{N}{=}2$ superconformal $\mathrm{SCA}_c$ & M & $\infty$ & $T$-channel infinite; $J$-channel terminates at degree $2$; $G$-channel at degree $3$ & Rem~\ref{rem:n2-shadow-tower} \\
\hline
Triplet $\mathcal{W}(p)$ (logarithmic; $p \geq 2$) & M & $\infty$ & $T$-channel forces $r_{\max}^T = \infty$ regardless of $c$ & \S\ref{sec:wp-shadow-depth} \\
\hline
$Y_{0,0,N}[\Psi]$ ($N \geq 2$; $\mathcal{W}_N$-type) & M & $\infty$ & contains $\mathcal{W}_N$ subalgebra with spin-$\geq 3$ generator & Thm~\ref{thm:y-shadow-depth} \\
\hline
$\mathcal{W}_\infty$ (limit) & M & $\infty$ & one-dimensional cyclic line at the inverse limit, the $M^*$ refinement ($\dim H^2_{\mathrm{cyc}}=1$) & Thm~\ref{thm:winfty-scalar} \\
\hline
Monster $V^\natural$ ($c = 24$, $\kappa = 12$) & M & $\infty$ & Griess algebra at weight~$2$; $\Delta \neq 0$ & Rem~\ref{rem:lattice:monster-shadow}; Thm~\ref{thm:schellekens-structured-subset} \\
\hline
\multicolumn{5}{|c|}{\textit{Subregular and non-principal $\mathcal{W}$-algebras}} \\
\hline
$\mathrm{BP}_k = \mathcal{W}_3^{(2)}(k)$ (Bershadsky--Polyakov) & M (candidate) & open & standard $c_{\mathrm{BP}}(k)$ and exact $T$-line inverse norm; full $J,G^+,G^-$ projection under $H_{\mathrm{bar}}+H_{\mathrm{proj}}$ & Thm~\ref{thm:bp-t-line-rational-k} \\
\hline
\multicolumn{5}{|c|}{\textit{$\Eone$-chiral}} \\
\hline
Yangian $Y(\fg)$ & $\Eone$-RTT & n/a & ordered RTT $R$-matrix tower; no Virasoro shadow depth & \S\ref{sec:modular-dg-shifted-yangian} \\
\hline
Quantum lattice $\Vlat_\Lambda^{N,q}$ (strict $\Eone$, $N \geq 3$) & ? & unknown & deformed braiding; same character as $V_\Lambda$; strict $\Eone$-chiral Koszul self-pairing under $q\mapsto -q$ & Thm~\ref{thm:quantum-lattice-structure} \\
\hline
\multicolumn{5}{|c|}{\textit{Wild: non-Koszul boundary}} \\
\hline
Kronecker $K_m$ ($m \geq 3$) & W & n/a & bar SS does not collapse; DT invariants replace shadows & Rem~\ref{rem:wild-quiver-class} \\
\hline
\end{tabular}%
}% end resizebox
}% end small
\end{table}

\noindent
\textbf{Archetype column}: G = Gaussian, L = Lie/tree, C = contact/quartic,
M = mixed/infinite, with the inverse-limit $M^*$ line recorded as a
refinement inside class~M; $\Eone$-RTT = ordered nonlocal RTT tower
data, W = wild (non-Koszul, shadow depth undefined).
The archetype classifies the shadow depth: at which
degree the tower of projections $\Theta_\cA^{\leq r}$ stabilizes (or
fails to stabilize). Each archetype class is realized by at least one
standard vertex-algebra family, and the G/L/C/M classification is
exhaustive for the completed chirally Koszul vertex-algebra families
treated in this monograph. The finite conductor-domain count does not
split off the inverse-limit $M^*$ refinement and does not count
$\Eone$-RTT, strict $\Eone$ quantum lattice, or wild rows. A separate
refined spectral component admits depths
$d \in \{2,3,4,5,\dotsc\} \cup \{\infty\}$; every finite value is
realized by lattice spectral packets, with
$d = 3 + \dim S_{r/2}$ for even unimodular lattices
(Theorem~\ref{thm:shadow-spectral-correspondence}). This is the
shadow-spectral correspondence after adjoining cusp-form data, not the
scalar master-table row of a positive-definite lattice VOA. On that
scalar row, Table~\ref{tab:shadow-taxonomy} keeps the positive-definite
lattice VOA in class~G with $r_{\max}=2$. The infinite scalar row is
characterized by $d = \infty$ iff the stress tensor is a strong
generator ($T \in T_{(1)}T$;
Corollary~\ref{cor:conformal-vector-infinite-depth}).

\begin{remark}[Pole depth, logarithmic depth, and shadow depth]
The invariants $p_{\max}$, $k_{\max}$, and $r_{\max}$ measure different
objects. For the standard $\beta\gamma_\lambda$ and $bc_\lambda$
families the generator OPE has $p_{\max}=1$, $d\log$ absorption gives
$k_{\max}=0$, and the conformal-weight family has $r_{\max}=4$. For a
positive-definite lattice VOA the scalar oscillator component has
$p_{\max}=2$, $k_{\max}=1$, and $r_{\max}=2$; root operators carry
lattice-graded simple-pole residues, but no Lie-bracket cubic scalar
survives the averaging map. For a non-critical affine Kac--Moody
algebra the current bracket supplies the cubic scalar shadow, so
$r_{\max}=3$.
The generator-pole values for the remaining standard rows are read
off the defining OPEs of Table~\ref{tab:standard-family-status}.
The current OPE
$J^a(z)J^b(w)\sim k\kappa_\fg^{ab}(z-w)^{-2}
+f^{ab}{}_cJ^c(w)(z-w)^{-1}$ has quadratic maximal pole,
\begin{equation}\label{eq:affine-KM-p-max}
p_{\max}\bigl(V_k(\fg)\bigr) = 2;
\end{equation}
the stress-tensor OPE has quartic central pole
$\tfrac{c/2}{(z-w)^4}$,
\begin{equation}\label{eq:vir-p-max}
p_{\max}(\mathrm{Vir}_c) = 4;
\end{equation}
and the top strong generator $W_N$ of the principal $\cW_N$-algebra,
of conformal weight~$N$, has self-OPE with maximal central pole of
order~$2N$,
\begin{equation}\label{eq:w-N-p-max}
p_{\max}(\cW_N) = 2N.
\end{equation}
\end{remark}

\begin{proposition}[Conductor-domain census count]
\label{prop:census-conductor-row-count}
\ClaimStatusProvedHere
The completed vertex-algebra families in the landscape census that
lie in the averaged-shadow conductor domain are partitioned as
\[
G:5,\qquad L:2,\qquad C:4,\qquad M:10.
\]
Thus the $G/L/C/M$ conductor domain has exactly $21$ rows. The Yangian
row is the separate $\Eone$-RTT row; it has an ordered RTT
$R$-matrix tower and no Virasoro shadow depth. The strict $\Eone$
quantum lattice row is unclassified in this four-archetype count, and
the wild Kronecker row is outside the Koszul domain.
\end{proposition}

\begin{proof}
The Gaussian rows are the free fermion, Heisenberg, lattice VOA,
Niemeier lattice VOA, and $Y_{1,1,1}[\Psi]$. The Lie rows are the
simply-laced and non-simply-laced affine Kac--Moody families. The
contact rows are $\beta\gamma$, $bc$, the symplectic boson
$\beta\gamma_{1/2}$, and the symplectic-fermion point
$\mathcal{SF}=bc_1$. The mixed rows are Virasoro, finite
$\mathcal W_N$, principal $\mathcal W_3$, Ising, $\mathcal N{=}2$
SCA, triplet $\mathcal W(p)$, $Y_{0,0,N}[\Psi]$,
$\mathcal W_\infty$, Monster, and Bershadsky--Polyakov. These are
precisely the rows counted by
Theorem~\ref{thm:uc-landscape-universality}; the remaining rows have
the non-vertex or non-Koszul status stated above.
\end{proof}

\begin{remark}[Wild quiver class W]
\label{rem:wild-quiver-class}
\index{shadow depth!wild class}%
\index{wild quiver!class W}%
\index{Kronecker quiver!shadow depth}%
The shadow depth classification G/L/C/M applies to chirally Koszul
algebras: those whose bar spectral sequence collapses at~$E_2$
(Theorem~\ref{thm:koszul-equivalences-meta}). The Kronecker quiver
$K_m$ (two vertices, $m$ parallel arrows) exhibits a sharp Koszul
boundary at $m = 3$: $K_1$ is finite type ($A_2$, class~L),
$K_2$ is the tame affine-Kronecker boundary, and $K_m$ for
$m \geq 3$ is wild. The label class~M belongs to non-terminating
chirally Koszul vertex-algebra shadows, not to a bare tame quiver
without a specified chiral algebra and bar comparison. In the wild regime, the symmetric Euler form
has $\det = 4 - m^2 < 0$ (indefinite), the bar spectral sequence
does not collapse at any finite page, and the signed Euler series
$\prod_{n \geq 1}(1 - t^n)^{m+2}$ acquires negative coefficients
starting at weight~$2$.

These failures obstruct the construction of a shadow obstruction
tower in the sense of Definition~\ref{def:shadow-depth-classification}:
the MC element $\Theta_\cA$ does not exist in the formal deformation
complex because the bar cohomology lacks concentration. The role of
the shadow tower is taken by the Donaldson--Thomas invariants
$\Omega(d_0, d_1)$ of the Kronecker quiver
(Kontsevich--Soibelman~\cite{KS10}, Reineke~\cite{Re03}),
which grow exponentially in $\max(d_0, d_1)$ for $m \geq 3$,
reflecting the wild representation type. The class~W designation
marks the boundary of modular Koszul duality: on classes G/L/C/M,
the shadow obstruction tower is the organizing structure; on
class~W, motivic DT theory supplies the algebraic substitute.

Computation:
\texttt{compute/lib/wild\_quiver\_chiral\_engine.py} (89 tests).
\end{remark}

%% ================================================================
%% Collision-residue r-matrix census
%% ================================================================

\begin{table}[ht]
\centering
\caption{Collision-residue $r$-matrix census for all standard families}
\label{tab:rmatrix-census}
\index{r-matrix@$r$-matrix!census table|textbf}
\index{collision residue!census table}
\index{pole absorption!landscape census}
\renewcommand{\arraystretch}{1.5}
{\small
\resizebox{\textwidth}{!}{%
\begin{tabular}{|l|c|c|c|l|c|}
\hline
\textbf{Algebra $\cA$ / channel}
 & $\boldsymbol{h}$
 & \textbf{OPE poles}
 & \textbf{$r$-matrix poles}
 & $\boldsymbol{r^{\mathrm{coll}}(z)}$
 & \textbf{CYBE} \\
\hline
\multicolumn{6}{|c|}{\textit{Class G (Gaussian): abelian simple-pole or regular}} \\
\hline
$\cH_k$ ($JJ$)
 & $1$ & $\{2\}$ & $\{1\}$
 & $k/z$
 & trivial \\
\hline
$V_\Lambda$ (Cartan $J^a J^b$)
 & $1$ & $\{2\}$ & $\{1\}$
 & $\langle\alpha^a,\alpha^b\rangle/z$
 & trivial \\
\hline
Free fermion $\psi$ ($\psi_i\psi_i$)
 & $\tfrac{1}{2}$ & $\{1\}$ & $\varnothing$
 & $0$
 & trivial \\
\hline
\multicolumn{6}{|c|}{\textit{Class C (contact): regular $r$-matrix, quartic leading interaction}} \\
\hline
$\beta\gamma$ ($\beta\gamma$, $\gamma\beta$)
 & $\lambda, 1{-}\lambda$ & $\{1\}$ & $\varnothing$
 & $0$
 & trivial \\
\hline
$bc$ ghosts ($bc$)
 & $\lambda, 1{-}\lambda$ & $\{1\}$ & $\varnothing$
 & $0$
 & trivial \\
\hline
$\beta\gamma$ (diagonal)
 & n/a & $\varnothing$ & $\varnothing$
 & $0$
 & trivial \\
\hline
\multicolumn{6}{|c|}{\textit{Class L (Lie/tree): collision residue $r^{\mathrm{coll}}(z) = k\Omega_{\mathrm{tr}}/z$; KZ normalization $r^{\mathrm{KZ}}(z)=\Omega_{\mathrm{KZ}}/((k{+}h^\vee)z)$}} \\
\hline
$\widehat{\fg}_k$ (diagonal $J^a J^a$)
 & $1$ & $\{2\}$ & $\{1\}$
 & $k\Omega_{\mathrm{tr}}^{aa}/z$
 & IBR \\
\hline
$\widehat{\fg}_k$ (off-diagonal $J^a J^b$)
 & $1$ & $\{1\}$ & $\varnothing$
 & $0$
 & trivial \\
\hline
\multicolumn{6}{|c|}{\textit{Class M (mixed/infinite): higher-pole $r$-matrices}} \\
\hline
$\mathrm{Vir}_c$ ($TT$)
 & $2$ & $\{4,2,1\}$ & $\{3,1\}$
 & $(c/2)/z^3 + 2T/z$
 & odd \\
\hline
$\cW_3$ ($TT$)
 & $2$ & $\{4,2,1\}$ & $\{3,1\}$
 & $(c/2)/z^3 + 2T/z$
 & odd \\
\hline
$\cW_3$ ($TW$)
 & $2,3$ & $\{2,1\}$ & $\{1\}$
 & $3W/z$
 & n/a \\
\hline
$\cW_3$ ($WW$)
 & $3$ & $\{6,4,3,2,1\}$ & $\{5,3,2,1\}$
 & $(c/3)/z^5 + 2T/z^3 + \partial T/z^2 + \cdots$
 & mixed \\
\hline
\end{tabular}%
}% end resizebox
}% end small
\end{table}

\noindent
\textbf{Extraction rule.}
The bar construction uses the $d\log$ kernel
$d\log(z_i - z_j)$, which absorbs one power of $(z{-}w)$:
an OPE pole at $z^{-n}$ produces an $r$-matrix pole at
$z^{-(n-1)}$. In particular, simple poles ($z^{-1}$) in the OPE
become regular and drop.
This is a collision-residue table. The affine entries are written in
the trace-form convention $k\Omega_{\mathrm{tr}}/z$; the KZ connection
uses $\Omega/((k+h^\vee)z)$ after Casimir and coupling rescaling. The
two conventions should not be interchanged when testing level-prefix
vanishing.

\textbf{Column~$h$}: conformal weight of the generators involved.
\textbf{CYBE column}: ``trivial'' = abelian $r$-matrix,
CYBE $[r_{12}, r_{13}] + \text{cyc.} = 0$ holds vacuously;
``IBR'' = the CYBE reduces to the infinitesimal braid relation
$[\Omega_{ij}, \Omega_{ik} + \Omega_{jk}] = 0$ on the Casimir
tensor (checked numerically for $\mathfrak{sl}_2$ and
$\mathfrak{sl}_3$ in
\texttt{test\_rmatrix\_landscape.py});
``odd'' = $r(z)$ has only odd-order poles and is therefore an odd
function of~$z$, so skew-symmetry
$r_{12}(z) + r_{21}(-z) = 0$ holds;
``mixed'' = even-order poles present (see
Remark~\ref{rem:ww-even-poles-census} below).
Computed across
\texttt{test\_rmatrix\_landscape.py} and
\texttt{test\_rmatrix\_poles\_comprehensive.py}
($205$ tests, all passing).
The detailed pole structure is in
Table~\ref{tab:rmatrix-pole-landscape} (the bar complex
tables chapter).

\begin{remark}[Bosonic parity constraint on collision residues]
\label{rem:bosonic-parity-census}
\index{bosonic parity!$r$-matrix constraint}
\index{r-matrix@$r$-matrix!bosonic parity}
For a single bosonic generator of integer conformal weight~$h$,
the same-generator OPE poles lie at even orders
$z^{-2h}, z^{-(2h-2)}, \ldots, z^{-2}$ (from the conformal
algebra) together with a descendant pole at~$z^{-1}$. After
$d\log$ absorption ($n \mapsto n{-}1$), the even-order poles
become $z^{-(2h-1)}, z^{-(2h-3)}, \ldots, z^{-1}$, all odd,
and the descendant pole drops. Consequence: the $r$-matrix of a
single bosonic generator whose OPE has only even-order poles
(above $z^{-1}$) has no even-order poles.

The maximal $r$-matrix pole order follows the formula
$\max\operatorname{pole}(r) = 2h - 1$: the Heisenberg
($h = 1$) gives $z^{-1}$; the Virasoro ($h = 2$) gives $z^{-3}$;
the $W$ generator of $\cW_3$ ($h = 3$) gives $z^{-5}$ in the $WW$
channel.

The constraint applies whenever all OPE poles above $z^{-1}$
have even order. For the stress tensor ($h = 2$), this holds:
the $TT$ OPE has poles at $\{4, 2, 1\}$, all even above $z^{-1}$,
and the $r$-matrix poles $\{3, 1\}$ are both odd.
For the $\cW_3$ $WW$ channel ($h = 3$), the OPE has an additional
odd-order pole at $z^{-3}$ from the composite field $\partial T$,
which shifts to the even-order pole $z^{-2}$ in the $r$-matrix.
This is not a violation: the bosonic parity constraint
(Proposition~\ref{prop:bosonic-parity-constraint}) excludes even
$r$-matrix poles only when all OPE poles above $z^{-1}$ are even,
a condition the $WW$ OPE does not satisfy.
\end{remark}

\begin{remark}[$r$-matrix vs the degree-$3$ MC equation]
\label{rem:rmatrix-vs-mc3-census}
\index{r-matrix@$r$-matrix!vs MC equation}
The $r$-matrix $r(z) = \Res^{\mathrm{coll}}_{0,2}(\Theta_\cA)$
is the genus-$0$ binary collision residue of the universal MC
element~$\Theta_\cA$. It carries only part of the degree-$3$
content of the MC equation: the CYBE on~$r(z)$ follows
from the degree-$3$ MC equation after collision extraction and the
Arnold relation on $\mathrm{Conf}_3(\mathbb{C})$, but the two live
in different spaces.

For class-L algebras (non-critical affine Kac--Moody,
$r_{\max} = 3$), the
CYBE exhausts the degree-$3$ content. For class-M algebras
(Virasoro, $\cW_N$, $r_{\max} = \infty$), the degree-$3$ MC
equation carries the homotopy Jacobiator, strictly more than
the CYBE alone.
See Remark~\ref{rem:degree3-vs-cybe} for the full discussion.
\end{remark}

\begin{remark}[Even-order poles in the $\cW_3$ $WW$ $r$-matrix]
\label{rem:ww-even-poles-census}
\index{W-algebra@$\cW$-algebra!$r$-matrix even poles}
The $WW$ channel of $\cW_3$ has four $r$-matrix poles:
$\{5, 3, 2, 1\}$. The even pole at $z^{-2}$ arises from the
$\partial T/(z{-}w)^3$ term in the $W(z)W(w)$ OPE (pole order~$3$
shifted to~$2$ by $d\log$ absorption). This does not contradict
Remark~\ref{rem:bosonic-parity-census}: the $WW$ OPE has poles at
orders $\{6, 4, 3, 2, 1\}$, which include the odd-order pole
$z^{-3}$. The bosonic parity constraint excludes even $r$-matrix
poles only when all OPE poles above $z^{-1}$ have even order;
the $WW$ channel fails this hypothesis. For $\cW_N$ with
$N \geq 3$, the higher-spin generators produce OPEs with
odd-order poles from composite fields, and even-order $r$-matrix
poles are generically present.
\end{remark}

\begin{table}[ht]
\centering
\caption{Scalar-component free energies $F_1(\cA)=\kappa(\cA)/24$ and
$F_2^{\mathrm{scalar}}(\cA)=7\kappa(\cA)/5760$ for $21$ standard families;
multi-weight rows require the scalar term plus
$\delta F_2^{\mathrm{cross}}$}
\label{tab:free-energy-landscape}
\index{free energy!landscape table}
\renewcommand{\arraystretch}{1.5}
{\small
\begin{tabular}{|l|c|c|c|c|}
\hline
\textbf{Algebra $\cA$}
 & $\boldsymbol{\kappa(\cA)}$
 & $\boldsymbol{F_1(\cA)=\kappa(\cA)/24}$
 & $\boldsymbol{F_2^{\mathrm{scalar}}(\cA)=7\kappa(\cA)/5760}$
 & \textbf{Lane} \\
\hline
\multicolumn{5}{|c|}{\textit{Free Fields}} \\
\hline
Free fermion $\psi$ & $\tfrac{1}{4}$ & $\tfrac{1}{96}$ & $\tfrac{7}{23040}$ & scalar \\
\hline
$bc$ ghosts ($\lambda{=}0$) & $-1$ & $-\tfrac{1}{24}$ & $-\tfrac{7}{5760}$ & scalar \\
\hline
$\beta\gamma$ ($\lambda{=}1$) & $1$ & $\tfrac{1}{24}$ & $\tfrac{7}{5760}$ & multi \\
\hline
$\beta\gamma$ ($\lambda{=}\tfrac{1}{2}$) & $-\tfrac{1}{2}$ & $-\tfrac{1}{48}$ & $-\tfrac{7}{11520}$ & scalar \\
\hline
\multicolumn{5}{|c|}{\textit{Heisenberg}} \\
\hline
$\mathcal{H}_1$ & $1$ & $\tfrac{1}{24}$ & $\tfrac{7}{5760}$ & scalar \\
\hline
\multicolumn{5}{|c|}{\textit{Affine Kac--Moody at $k = 1$}} \\
\hline
$\widehat{\mathfrak{sl}}_2$ & $\tfrac{9}{4}$ & $\tfrac{3}{32}$ & $\tfrac{7}{2560}$ & scalar \\
\hline
$\widehat{\mathfrak{sl}}_3$ & $\tfrac{16}{3}$ & $\tfrac{2}{9}$ & $\tfrac{7}{1080}$ & scalar \\
\hline
$\widehat{\mathfrak{so}}_{5}$ (type $B_2$) & $\tfrac{20}{3}$ & $\tfrac{5}{18}$ & $\tfrac{7}{864}$ & scalar \\
\hline
$\widehat{\mathfrak{sp}}_{4}$ (type $C_2$) & $\tfrac{20}{3}$ & $\tfrac{5}{18}$ & $\tfrac{7}{864}$ & scalar \\
\hline
$\widehat{\mathfrak{so}}_{8}$ (type $D_4$) & $\tfrac{49}{3}$ & $\tfrac{49}{72}$ & $\tfrac{343}{17280}$ & scalar \\
\hline
$\widehat{G}_2$ & $\tfrac{35}{4}$ & $\tfrac{35}{96}$ & $\tfrac{49}{4608}$ & scalar \\
\hline
$\widehat{F}_4$ & $\tfrac{260}{9}$ & $\tfrac{65}{54}$ & $\tfrac{91}{2592}$ & scalar \\
\hline
$\widehat{E}_6$ & $\tfrac{169}{4}$ & $\tfrac{169}{96}$ & $\tfrac{1183}{23040}$ & scalar \\
\hline
$\widehat{E}_7$ & $\tfrac{2527}{36}$ & $\tfrac{2527}{864}$ & $\tfrac{17689}{207360}$ & scalar \\
\hline
$\widehat{E}_8$ & $\tfrac{1922}{15}$ & $\tfrac{961}{180}$ & $\tfrac{6727}{43200}$ & scalar \\
\hline
\multicolumn{5}{|p{0.96\textwidth}|}{%
\footnotesize\textit{Caveat (universal vs.\ simple quotient at simply-laced $k=1$):} the affine
KM entries above compute $\kappa$ on the \emph{universal} vacuum module $V_k(\fg)$ at $k=1$
with $\dim\fg$ independent weight-$1$ currents; $\kappa = \dim\fg \cdot (k+h^\vee)/(2h^\vee)$
by the Sugawara formula. At simply-laced $k=1$ the universal module is not simple: the Frenkel--Kac--Segal
isomorphism $L(1,0) = V_k(\fg)/\cI = V_{\Lambda_\fg}$ identifies root currents with composites
of the $\operatorname{rank}(\fg)$ Heisenberg generators, reducing the effective generator count
and the modular characteristic to $\kappa(V_{\Lambda_\fg}) = \operatorname{rank}(\fg)$
(Proposition~\ref{V1-prop:level1-kappa-reduction}). The lattice-VOA rows below record
$\kappa(V_{\Lambda_\fg})$ for $\Lambda_\fg$ the root lattice of $\fg$ (simple quotient); the
affine KM rows record $\kappa(V_k(\fg))$ for the universal vacuum (before quotient). Both are
correct in their respective categories; the FKS collapse relates them.
}\\
\hline
\multicolumn{5}{|c|}{\textit{$\mathcal{W}$-Algebras}} \\
\hline
$\mathrm{Vir}_{c}$ & $c/2$ & $c/48$ & $7c/11520$ & scalar \\
\hline
$\mathcal{W}_{3,c}$\textsuperscript{$\mathrm{mod/comp}$}
 & $5c/6$ & $5c/144$ & $7c/6912$ & multi \\
\hline
\multicolumn{5}{|c|}{\textit{Lattice VOAs}} \\
\hline
$V_{D_4}$ & $4$ & $\tfrac{1}{6}$ & $\tfrac{7}{1440}$ & scalar \\
\hline
$V_{E_8}$ & $8$ & $\tfrac{1}{3}$ & $\tfrac{7}{720}$ & scalar \\
\hline
Leech $V_{\Lambda_{24}}$ & $24$ & $1$ & $\tfrac{7}{240}$ & scalar \\
\hline
\multicolumn{5}{|c|}{\textit{Exceptional}} \\
\hline
Monster $V^\natural$ & $12$ & $\tfrac{1}{2}$ & $\tfrac{7}{480}$ & scalar \\
\hline
\end{tabular}
}% end small
\end{table}

\noindent
The principal $\mathcal W_3$ row in
Table~\ref{tab:free-energy-landscape} is conditional on
$H_{\mathrm{diag}}^{g=1}+H_{W_3}^{\mathrm{DS/bar}}$.  Its first
entry is the normalized genus-one diagonal trace; the subsequent
scalar entries are its Faber--Pandharipande multiples.

\noindent
\textbf{Lane column}: ``scalar'' = uniform-weight algebra (single generator or
all generators of the same conformal weight), for which
$F_g^{\mathrm{sc}}(\cA)=\kappa(\cA) \cdot \lambda_g^{\mathrm{FP}}$
is proved at all genera; on these scalar/exact rows
$F_g(\cA)=F_g^{\mathrm{sc}}(\cA)$
(Theorem~\ref{thm:genus-universality});
``multi'' = multi-weight algebra, for which
$F_g(\cA)=F_g^{\mathrm{sc}}(\cA)+\delta F_g^{\mathrm{cross}}(\cA)$
where
$F_g^{\mathrm{sc}}(\cA)=\kappa(\cA) \cdot \lambda_g^{\mathrm{FP}}$,
and the cross-channel correction is determined by
Theorem~\ref{thm:multi-weight-genus-expansion}: the scalar coefficient
$F_g^{\mathrm{sc}}(\cA)=\kappa(\cA) \cdot \lambda_g^{\mathrm{FP}}$
($g{=}1$ unconditional; all-weight at $g{=}1$) is the full coefficient
at $g=1$ and is not the full coefficient for $g \ge 2$ without the
$\delta F_g^{\mathrm{cross}}$ term
(Remark~\ref{rem:propagator-weight-universality}). The universal ratio
$F_2^{\mathrm{sc}}/F_1^{\mathrm{sc}} = 7/240$ is independent of~$\cA$.

\smallskip\noindent
For the first nontrivial multi-weight family, the correction terms are
\[
\delta F_2^{\mathrm{cross}}(\cW_3)
\;=\;
\frac{c+204}{16c},
\qquad
\delta F_3^{\mathrm{cross}}(\cW_3)
\;=\;
\frac{5c^3 + 3792c^2 + 1149120c + 217071360}{138240\,c^2},
\]
by Theorem~\ref{thm:multi-weight-genus-expansion}\textup{(vi)} and
Proposition~\ref{prop:w3-genus3-cross-channel}. On the Virasoro
family, $\delta F_g^{\mathrm{cross}}(\Vir_c)=0$ for every $g \ge 1$
because $\Vir_c$ has a single strong generator and therefore lies on
the uniform-weight locus.

\begin{table}[ht]
\centering
\caption{Shadow obstruction tower invariants: discriminant, shadow coefficients, growth rate}
\label{tab:shadow-invariants-landscape}
\index{shadow tower!landscape invariants}
\index{critical discriminant!landscape table}
\renewcommand{\arraystretch}{1.5}
{\small
\begin{tabular}{|l|c|c|c|c|c|c|}
\hline
\textbf{Algebra $\cA$}
 & \textbf{Class}
 & $\boldsymbol{S_3}$
 & $\boldsymbol{S_4}$
 & $\boldsymbol{\Delta = 8\kappa S_4}$
 & $\boldsymbol{\rho}$
 & $\boldsymbol{\kappa{+}\kappa'}$ \\
\hline
$\mathcal{H}_k$ & G & $0$ & $0$ & $0$ & $0$ & $0$ \\
\hline
Free fermion & G & $0$ & $0$ & $0$ & $0$ & $0$ \\
\hline
$bc$ ($\lambda$) & C & $0$ & $\neq 0$ & $\neq 0^\dagger$ & finite$^\dagger$ & $0$ \\
\hline
$V_\Lambda$ (rank $d$) & G & $0$ & $0$ & $0$ & $0$ & $0$ \\
\hline
$\widehat{\fg}_k$ & L & $\neq 0$ & $0$ & $0$ & $0$ & $0$ \\
\hline
$\beta\gamma_\lambda$ & C & $0$ & $\neq 0$ & $\neq 0^\dagger$ & finite$^\dagger$ & $0$ \\
\hline
$\mathrm{Vir}_c$ & M & $2$ & $\dfrac{10}{c(5c{+}22)}$
 & $\dfrac{40}{5c{+}22}$ & $\sqrt{\dfrac{180c{+}872}{c^2(5c{+}22)}}$ & $13$ \\[8pt]
\hline
$\mathcal{W}_{3,c}$, $T$-line & M & $2$ & $\dfrac{10}{c(5c{+}22)}$
 & $\dfrac{40}{5c{+}22}$ & \text{as Vir}
 & \multirow{2}{*}{$250/3$\textsuperscript{$\mathrm{mod/comp}$}} \\[4pt]
\cline{1-6}
$\mathcal{W}_{3,c}$, $W$-line & \text{even} & $0$ & $\dfrac{2560}{c(5c{+}22)^3}$
 & $\dfrac{20480}{3(5c{+}22)^3}$ & $\neq 0$ & \\[4pt]
\hline
$\mathcal{W}_N$ & M & $\neq 0$ & $\neq 0$ & $\neq 0$ & $> 0$
 & $\varrho_N K_N^c$\textsuperscript{$\mathrm{mod/comp}$} \\
\hline
Monster $V^\natural$ ($c{=}24$) & M & $2$ & $\dfrac{5}{1704}$
 & $\dfrac{20}{71}$ & $> 0$ & $13$\textsuperscript{$\S$} \\[4pt]
\hline
\end{tabular}
}% end small
\end{table}

\noindent
The superscript ``mod/comp'' in
Table~\ref{tab:shadow-invariants-landscape} denotes the package
$H_{\mathrm{diag}}^{g=1}+H_{W_N}^{\mathrm{DS/bar}}$.  In the
$N=3$ row it denotes
$H_{\mathrm{diag}}^{g=1}+H_{W_3}^{\mathrm{DS/bar}}$ and yields the
conditional value $250/3$ from the exact central conductor~$100$.

\noindent
${}^\dagger$For $\beta\gamma_\lambda$ and $bc_\lambda$, the class-C
witness is the standard conformal-weight family:
$S_2 = 6\lambda^2 - 6\lambda + 1$, $S_3 = 0$,
$S_4 = -5/12$, and $S_r = 0$ for $r \geq 5$.
Thus $\Delta=8\kappa S_4=-(10/3)(6\lambda^2-6\lambda+1)$
generically; the class-C label records finite shadow depth, not
vanishing discriminant.
The internal weight-changing line has zero shadow tower, while the
T-line is class~M; neither one-dimensional slice is the class-C
witness by itself.
The Virasoro discriminant complementarity is
$\Delta(c) + \Delta(26{-}c) = 6960/[(5c{+}22)(152{-}5c)]$
(Corollary~\ref{cor:discriminant-atlas}(ii)).
The exact central conductor for $\mathcal{W}_N$ is
$K_N^c = 4N^3 - 2N - 2$; explicitly:
$K_2^c = 26$, $K_3^c = 100$, $K_4^c = 246$, $K_5^c = 488$.
The self-dual central charges are $c^* = K_N^c/2$: $c^* = 13$ for Virasoro,
$c^* = 50$ for $\mathcal{W}_3$, $c^* = 123$ for $\mathcal{W}_4$.
The unconditional scalar Verdier values on the four witness families
form the set
\[
\{0,\,13\}.
\]
On
$H_{\mathrm{diag}}^{g=1}+H_{W_3}^{\mathrm{DS/bar}}$, the principal
$\mathsf M$ extension contributes the conditional value $250/3$ and
gives the comparison set
\[
\{0,\,13,\,250/3\}.
\]
The BP $\mathsf M$-extension contributes the
exact central conductor~$50$ and an open genus-$1$ modular slot.  After
adjoining the Mukai-enhanced K3 $\mathsf B$ candidate under
$H_{\mathsf B}$ together with
$H_{\mathrm{diag}}^{g=1}+H_{W_3}^{\mathrm{DS/bar}}$, the conditional
numerical ledger is
\[
\{0,\,8,\,13,\,250/3\}.
\]
The former BP proposal~$25/3$ has status \ClaimStatusConjectured as
typed in Remark~\ref{rem:bp-census-convention-separation}. The zero
bucket contains the Heisenberg, affine Kac--Moody, and
$\beta\gamma_\lambda/bc_\lambda$ Verdier scalar pairs:
\[
c_{\beta\gamma}(\lambda)=2(6\lambda^2-6\lambda+1),\qquad
c_{bc}(\lambda)=-2(6\lambda^2-6\lambda+1),
\]
\[
\kappa(\beta\gamma_\lambda)=6\lambda^2-6\lambda+1,\qquad
\kappa(bc_\lambda)=-(6\lambda^2-6\lambda+1).
\]
Hence $K^c=K^\kappa=0$ for the $\beta\gamma/bc$ Verdier pair. The Yangian
row has no scalar conductor value in this sense.
${}^\S$For the Monster module $V^\natural$: the complementarity
sum $\kappa + \kappa' = 13$ is the Virasoro-sector value
($\kappa(V^\natural) = 12$, $\kappa(\mathrm{Vir}_2) = 1$; sum $= 13$),
since $\dim V_1^\natural = 0$ forces $\kappa = c/2 = 12$
.

In the Virasoro row, the displayed partner $\mathrm{Vir}_{26-c}$ is
the proved M/S-level same-family shadow used in the complementarity
and semi-infinite calculations. The stronger H-level
infinite-generator realization is delegated to the exact
$W_\infty$ MC4 data: construct the filtered H-level target, prove
the residue identities
$\mathsf{C}^{\mathrm{res}}_{s,t;u;m,n}(N)
=\mathsf{C}^{\mathrm{DS}}_{s,t;u;m,n}(N)$,
and discharge them on the finite-detection packet $\mathcal{I}_N$.
The first live higher-spin packets are already explicit: stage-$3$
consists of fifteen primary coefficients, reduced by the explicit
$W_3$ OPE to three nonzero entries and twelve forced zeros. Stage-$4$
is cut by skew-symmetry to six live coefficients with twenty-three
forced zeros, organized as three local OPE blocks. The live
stage-$4$ comparison is the exact six-entry identity packet
on~$\mathcal{I}_4$:
\[
\mathsf{C}^{\mathrm{res}}_{s,t;u;m,n}(4)
=
\mathsf{C}^{\mathrm{DS}}_{s,t;u;m,n}(4)
\]
in the channels
\[
\begin{gathered}
(3,3;4;0,2),\qquad
(4,4;4;0,4),\\
(3,4;3;0,4),\qquad
(3,4;4;0,3),\qquad
(4,4;2;0,6),\\
(3,4;2;0,5).
\end{gathered}
\]
The last two channels are fixed exactly by
$\mathsf{C}^{\mathrm{DS}}_{4,4;2;0,6}=2$ and
$\mathsf{C}^{\mathrm{DS}}_{3,4;2;0,5}=0$.
Under the stage-$4$ Ward-normalized H-level input
(Definition~\ref{def:winfty-stage4-ward-normalized}),
Corollary~\ref{cor:winfty-stage4-residue-four-channel} contracts this
live packet to the four higher-spin channels.
Proposition~\ref{prop:winfty-mc4-frontier-package} identifies this
boundary exactly, and
Proposition~\ref{prop:w4-ds-ope-explicit} makes the principal
Drinfeld--Sokolov side explicit, so the live stage-$4$ question is
equality of the bar-side residue channels with those explicit targets.

For the Yangian row, ``PH'' records the proved DK core through DK-2/3
on its natural loci: the chain-level bar-cobar / $R \mapsto R^{-1}$
duality
(Theorem~\ref{thm:yangian-koszul-dual}), the evaluation-locus
factorization comparison
(Theorem~\ref{thm:factorization-dk-eval}), and the
evaluation-generated factorization core at all simple types
(Corollary~\ref{cor:dk23-all-types}; for finite-dimensional
type~$A$, Conjecture~\ref{conj:dk-fd-typeA} gives only the corresponding
thick-generation extension criterion). The full
category-$\mathcal{O}$/KL completion and the RTT-complete H-level
comparison are not being claimed in this table.

The remaining Yangian open problem begins beyond that proved core: the
extension to the ordinary-derived/completed/coderived enlargement and
then DK-4/DK-5, detected by the identities
$\mathsf{K}^{\mathrm{line}}_{a,b}(N)
=\mathsf{K}^{\mathrm{RTT}}_{a,b}(N)$
on the boundary strip $\Delta_{a,0}(N)$. For the standard type-A RTT
tower, Proposition~\ref{prop:yangian-dk4-typea-frontier} reduces this
to the auxiliary-kernel identity, and under one-loop exactness,
Corollary~\ref{cor:yangian-typea-degree2-plus-generators} sharpens it
to the single degree-$2$ mixed-tensor coefficient on the ordered
fundamental line, while
Proposition~\ref{prop:yangian-dk5-compact-generators} reduces the
completed module-category enhancement to the compact generator packet.
The first faithful-evaluation boundary-strip packets are already
closed computationally: $\Delta_{a,0}(4)$ and $\Delta_{a,0}(6)$
vanish on generic tensor lengths $2$ and $3$ in ranks $2,3,4$, so the
live Yangian target is the uniform all-rank transport of this
packet, not the first visible low-stage defects.

Central charges, $c{+}c'$, and $\kappa$ are not defined for
$\Eone$-chiral algebras (Yangians) in the same sense as for
vertex algebras.

\begin{remark}[Same central charge, different modular characteristic:
$V^\natural$ vs $V_{\Lambda_{24}}$]%
\label{rem:census-moonshine-leech-discrimination}%
\index{Monster module!modular characteristic}%
\index{Leech lattice!modular characteristic}%
\index{modular characteristic!discrimination at same $c$}%
The Monster module $V^\natural$ and the Leech lattice VOA
$V_{\Lambda_{24}}$ both have $c = 24$, but their modular
characteristics differ: $\kappa(V^\natural) = 12$ and
$\kappa(V_{\Lambda_{24}}) = 24$. This is the first instance in
the landscape where two algebras at the same central charge
belong to different shadow depth classes: $V^\natural$ is class~M
(infinite tower, single strong generator $T$ of weight~$2$,
$\Delta \neq 0$), while $V_{\Lambda_{24}}$ is class~G (Gaussian,
tower terminates at degree~$2$, $24$ weight-$1$ generators).

The root cause is: the modular characteristic depends on
the full algebra, not the Virasoro subalgebra alone. For the
Leech lattice VOA, the $24$ weight-$1$ Heisenberg generators
each contribute $\kappa = 1$ to the genus-$1$ obstruction,
giving $\kappa(V_{\Lambda_{24}}) = \operatorname{rank}(\Lambda_{24}) = 24$.
For the Monster module, $\dim V_1^\natural = 0$ eliminates all
weight-$1$ contributions, and the Virasoro sector alone determines
$\kappa(V^\natural) = c/2 = 12$.

The scalar free energy distinguishes the two: $F_1(V^\natural) = 1/2$
while $F_1(V_{\Lambda_{24}}) = 1$. The shadow obstruction towers
diverge qualitatively: $V_{\Lambda_{24}}$ has $\Delta = 0$ (all
higher shadows vanish), while $V^\natural$ has $\Delta = 20/71 \neq 0$
(infinite tower forced by the single-line dichotomy).
See Remark~\ref{rem:lattice:monster-shadow} for the detailed
computation.
\end{remark}

\begin{proposition}[Paired standard-tower MC4 completion packets;
 \ClaimStatusConditional]
\label{prop:paired-standard-mc4-frontier}
For the standard principal
\texorpdfstring{$W_\infty$}{W_infty} tower and the standard RTT Yangian
tower:
\begin{enumerate}[label=\textup{(\alph*)}]
\item the completed M-level comparison data are proved,
 given by
 Corollary~\ref{cor:winfty-standard-mc4-package} and
 Corollary~\ref{cor:yangian-standard-mc4-package};
\item the remaining H-level comparison reduces to finite-detection
 packets, namely the seed packets~$\mathcal{I}_N$ on the
 \texorpdfstring{$W_\infty$}{W_infty} side and the boundary strips
 $\Delta_{a,0}(N)$ on the Yangian side; on the
 \texorpdfstring{$W_\infty$}{W_infty} side,
 Proposition~\ref{prop:winfty-ds-stage-growth-packet} linearizes the
 promotion $\mathcal{I}_4\to\mathcal{I}_N$ as repeated closure of the
 incremental packets~$\mathcal{J}_{M+1}$, further reduced by
 Corollary~\ref{cor:winfty-ds-stage-growth-top-parity} to their
 top-pole/parity packets~$\mathcal{J}_{M+1}^{\mathrm{red}}$.
 On the Yangian side,
 The boundary-strip identity family
 \[
 \mathsf{K}^{\mathrm{line}}_{a,b}(N)
 =
 \mathsf{K}^{\mathrm{RTT}}_{a,b}(N)
 \]
 is sharpened in the standard type-$A$ RTT tower by
 Proposition~\ref{prop:yangian-dk4-typea-frontier} to the
 auxiliary-kernel identity, and under one-loop exactness by
 Corollary~\ref{cor:yangian-typea-degree2-plus-generators} to one
 degree-$2$ mixed-tensor coefficient together with the compact
 generator packet. On the factorization side,
 Corollary~\ref{cor:yangian-typea-factorization-local-closure}
 shows that this local coefficient is already closed; the live
 local move is dg-shifted transport to that factorization coefficient.

 \emph{Yangian side: H-level target and realization.}
 The H-level target already exists canonically as the
 formal-moduli tangent Lie algebra
 (Proposition~\ref{prop:yangian-canonical-hlevel-target}), with
 canonical associative dg model
 $U^{\mathrm{comp}}(\mathfrak{g}_{\cA})$
 (Proposition~\ref{prop:yangian-canonical-envelope}).
 Once that dg model carries the RTT-adapted
 finite-quotient/shared-seed data,
 Proposition~\ref{prop:yangian-typea-realization-criterion} identifies
 its standard RTT realization, and
 Corollary~\ref{cor:yangian-canonical-realization-to-spectral-seed}
 reduces the remaining input to the standard spectral vector
 seed-and-shift datum. On the spectral vector-line locus,
 Corollary~\ref{thm:yangian-canonical-realization-plus-vector-line}
 closes DK-4/DK-5 on the canonical target; on the sharper
 equivariant multiplicative spectral realization locus,
 Corollary~\ref{thm:yangian-canonical-realization-plus-one-seed}
 contracts this further to the single canonical spectral seed
 $V^\omega(0)=J_q^\omega(V(0))$.
 Corollary~\ref{thm:yangian-formal-moduli-plus-core-realization}
 expresses the remaining categorical step as compact-core realization
 of a chosen finite-dimensional factorization DK pair.
\end{enumerate}
Thus the standard-tower MC4 target is already reduced, on both sides,
to exact finite-detection packets. What remains is packet closure, not
another generic completion-existence problem.
\end{proposition}

\begin{proof}
Part~\textup{(a)} is exactly
Corollaries~\ref{cor:winfty-standard-mc4-package}
and~\ref{cor:yangian-standard-mc4-package}. On the
\texorpdfstring{$W_\infty$}{W_infty} side,
Corollary~\ref{cor:winfty-hlevel-comparison-criterion} reduces the
H-level comparison to compatible finite quotients, and
Proposition~\ref{prop:winfty-mc4-frontier-package} identifies the
relevant finite-detection packets~$\mathcal{I}_N$, while
Corollary~\ref{cor:winfty-stage4-residue-four-channel} gives the
four-channel refinement on the stage-$4$ Ward-normalized H-level
locus of Definition~\ref{def:winfty-stage4-ward-normalized}, and
Theorem~\ref{thm:winfty-all-stages-rigidity-closure} closes the
stage-$4$ packet inside the strong completion tower of
Theorem~\ref{thm:completed-bar-cobar-strong}.
On the Yangian side,
Corollary~\ref{cor:yangian-hlevel-comparison-criterion} reduces the
comparison to a filtered RTT-level target,
Proposition~\ref{prop:dg-shifted-rtt-boundary-seed} reduces the
coefficient problem to the boundary strips~$\Delta_{a,0}(N)$,
Proposition~\ref{prop:yangian-dk4-typea-frontier} gives the standard
type-$A$ sharpening to the auxiliary-kernel/residue packet,
Corollary~\ref{cor:yangian-typea-degree2-plus-generators} gives the
single degree-$2$ coefficient plus compact-generator formulation under
one-loop exactness,
Corollary~\ref{cor:yangian-typea-factorization-local-closure} closes
that local coefficient on the factorization side, and
Proposition~\ref{prop:yangian-canonical-hlevel-target} supplies the
canonical H-level target itself, while
Proposition~\ref{prop:yangian-typea-realization-criterion} realizes the
shared-seed finite-quotient data to an RTT-adapted realization of that
target, and
Corollary~\ref{thm:yangian-formal-moduli-plus-core-realization}
expresses the remaining categorical step. The last sentence is a
reformulation of these proved reductions.
\end{proof}

\begin{remark}[Current first live standard-tower packets]
On the \texorpdfstring{$W_\infty$}{W_infty} side, the stage-$4$
packet is closed, so the first next reduced packet is
$\mathcal{J}_5^{\mathrm{red}}$
\textup{(}Corollary~\textup{\ref{cor:winfty-ds-stage5-reduced-packet}}\textup{)}.
On the visible-pairing locus this stage-$5$ frontier contracts to the
single coefficient comparison isolated by
Proposition~\ref{prop:winfty-stage5-one-coefficient-reduction} and
Corollary~\ref{cor:winfty-stage5-exact-remaining-input}. On the
Yangian side, the local factorization coefficient is already closed by
Corollary~\ref{cor:yangian-typea-factorization-local-closure}; the
remaining live step is dg-shifted transport to that coefficient
together with the compact-core realization supplied by
Corollary~\ref{thm:yangian-formal-moduli-plus-core-realization}.
\end{remark}

\begin{corollary}[Minimal closure conditions for the standard-tower MC4
 completion target; \ClaimStatusConditional]
\label{cor:paired-standard-mc4-closure}
Under the hypotheses of
Proposition~\ref{prop:paired-standard-mc4-frontier}:
\begin{enumerate}[label=\textup{(\alph*)}]
\item the standard \texorpdfstring{$W_\infty$}{W_infty} side closes once
 the stagewise packet identities on~$\mathcal{I}_N$ are proved; the
 first unconditional finite packet is the six-entry stage-$4$ block
 on~$\mathcal{I}_4$, and on the stage-$4$ Ward-normalized H-level
 locus of Definition~\ref{def:winfty-stage4-ward-normalized} its live
 content is the four higher-spin channels of
 Corollary~\ref{cor:winfty-stage4-residue-four-channel};
\item on the standard type-$A$ RTT Yangian tower, the H-level
 comparison closes once the boundary-strip packet
 $\Delta_{a,0}(N)=0$ is proved and the completed module-category
 enhancement matches on compact generators; under one-loop exactness
 and seed normalization, the local boundary-strip packet is itself
 reduced to the single degree-$2$ coefficient of
 Corollary~\ref{cor:yangian-typea-degree2-plus-generators}, and on
 the factorization side that coefficient is already closed by
 Corollary~\ref{cor:yangian-typea-factorization-local-closure}. Once
 proved finite RTT quotients are added, the canonical formal-moduli
 target of Proposition~\ref{prop:yangian-canonical-hlevel-target}
 acquires its RTT-adapted realization by
 Proposition~\ref{prop:yangian-typea-realization-criterion}, and
 Corollary~\ref{thm:yangian-formal-moduli-plus-core-realization}
 reduces the completed categorical step to compact-core realization of
 a chosen finite-dimensional factorization DK pair.
\end{enumerate}
\end{corollary}

\begin{proof}
Part~\textup{(a)} combines
Corollaries~\ref{cor:winfty-stage4-closure-criterion}
and~\ref{cor:winfty-stage4-residue-four-channel}. Part~\textup{(b)} is
exactly
Corollaries~\ref{cor:yangian-typea-mc4-closure-criterion},
\ref{cor:yangian-typea-factorization-local-closure},
\ref{thm:yangian-typea-realization-plus-compacts},
and~\ref{thm:yangian-typea-realization-plus-fundamental-packet},
together with the canonical-target sharpening of
Corollary~\ref{thm:yangian-formal-moduli-plus-core-realization}.
\end{proof}

\begin{remark}[Conditional MC5 dependence]
\label{rem:paired-standard-mc4-closure-mc5}
Corollary~\ref{cor:paired-standard-mc4-closure} records only the
unconditional MC4 packet closure. Once the standard
\texorpdfstring{$W_\infty$}{W_infty} packet identities hold and the
realized Yangian DK-4/DK-5 data are in place, the successive
standard-tower MC5 dependence package is the conjectural reduction
\textup{(}Conjecture~\textup{\ref{conj:standard-tower-mc5-reduction}}\textup{)},
and on the canonical Yangian formal-moduli locus this sharpens to the
corresponding conditional closure statement
\textup{(}Corollary~\textup{\ref{cor:standard-tower-mc5-closure}}\textup{)}.
\end{remark}

\smallskip
\noindent\textbf{Proof tiers.} The invariants in the table rest on
different levels of the proof architecture:

\begin{enumerate}
\item[\textbf{T1.}] \emph{Unconditional} (free fields, Heisenberg):
 Koszulness is elementary ($d_{\mathrm{bracket}} = 0$ or $\dim V \le 2$).
 All invariants follow by direct computation.
\item[\textbf{T1$\frac{1}{2}$.}] \emph{Unconditional via PBW}
 (non-critical Kac--Moody and Virasoro at all genera):
 Koszulness proved via the PBW spectral sequence
 (Theorems~\ref{thm:km-chiral-koszul}
 and~\ref{thm:virasoro-chiral-koszul}); higher-genus PBW concentration
 proved unconditionally for Kac--Moody and Virasoro
 (Theorems~\ref{thm:pbw-allgenera-km}
 and~\ref{thm:pbw-allgenera-virasoro}).
 All invariants follow from the main theorems.
 The critical affine level is the FF boundary: PBW degeneration remains,
 but diagonal Koszul concentration fails.
\item[\textbf{T2.}] \emph{Unconditional higher-spin finite type}
 (principal finite-type $\mathcal{W}$-algebras):
 Koszulness is proved via the PBW spectral sequence in the finite-type
 cases, and higher-genus PBW concentration is now unconditional at all
 genera by Theorem~\ref{thm:pbw-allgenera-principal-w}.
\item[\textbf{DK-core.}] \emph{$\Eone$-Yangian regime}:
 the proved content is the chain-level $\Eone$-bar-cobar / DK square
 and the evaluation-locus factorization comparison
 \textup{(}Theorems~\ref{thm:e1-chiral-koszul-duality},
 \ref{thm:derived-dk-yangian}, \ref{thm:factorization-dk-eval}\textup{)},
 together with DK-2/3 on the evaluation-generated factorization core
 at all simple types
 \textup{(}Corollary~\ref{cor:dk23-all-types}; for finite-dimensional
 type~$A$, Conjecture~\ref{conj:dk-fd-typeA} only records the matching
 thick-generation extension criterion\textup{)}.
 What remains conjectural is the extension beyond that natural core
 and the DK-4/DK-5 dg-shifted RTT-complete comparison.
 The Virasoro-based invariants ($c$, $\kappa$) are not defined.
\end{enumerate}

\smallskip
\noindent\textsuperscript{$\dagger$}The Heisenberg algebra is the
abelian current algebra with its own level parameter:
$\kappa(\mathcal H_k)=k$ in rank one and
$\kappa(\mathcal H_{\ell}^{\oplus d})=d\ell$ in rank~$d$.
It is not the $h^\vee=0$ specialization of the non-abelian simple
affine formula $\kappa=(k+h^\vee)d/(2h^\vee)$.
The dual $\mathrm{Sym}^{\mathrm{ch}}(V^*)$ is a curved (non-conformal) chiral algebra,
so $c + c'$ is not defined in the usual sense.
The duality $\kappa + \kappa' = 0$ still holds: $\kappa' = -\kappa$
(see Theorem~\ref{thm:genus-universality}(ii)).
\textsuperscript{$\ddagger$}For the Virasoro row, the involution
$c \mapsto 26-c$ supplies the proved same-family partner governing the
current M/S-level genus and semi-infinite calculations. Identifying the
corresponding H-level infinite-generator dual object is the remaining
identification problem; the MC4 structural theorem is
Theorem~\ref{thm:completed-bar-cobar-strong}.
\textsuperscript{$\|$}The Leech lattice VOA $V_{\Lambda_{24}}$ is self-dual
as a lattice VOA: $\Lambda_{24}$ is unimodular, so
$V_{\Lambda_{24}}^! \cong V_{\Lambda_{24}}$
(Theorem~\ref{thm:lattice:unimodular-self-dual}).
The complementarity sum $c + c'$ is not defined in the
Feigin--Frenkel sense (no level shift); self-duality gives
$\kappa' = \kappa = 24$, so $\kappa + \kappa' = 48$.
\textsuperscript{$\S$}The Monster module $V^\natural$ has a single strong
generator $T$ of weight~$2$ ($\dim V_1^\natural = 0$), so the Virasoro
subalgebra determines~$\kappa$. The Koszul dual listed is the Virasoro-sector
partner $\mathrm{Vir}_{26-c} = \mathrm{Vir}_2$;
the full Koszul dual of $V^\natural$ as a VOA involves the
weight-$2$ Griess algebra and is not a standard-family object.
The complementarity sum $c + c' = 26$ is the Virasoro-sector value.

\begin{table}[ht]
\centering
\caption{Koszulness landscape: 15 algebras classified
(see also the detailed Table~\textup{\ref{tab:koszulness-landscape}} in
Chapter~\textup{\ref{chap:koszul-pairs}})}
\label{tab:koszulness-landscape-census}
\index{Koszul property!classification landscape}
\renewcommand{\arraystretch}{1.3}
{\small
\begin{tabular}{|l|c|c|l|}
\hline
\textbf{Algebra}
 & \textbf{Status}
 & \textbf{Equivalence}
 & \textbf{Proof mechanism} \\
\hline
\multicolumn{4}{|c|}{\textit{Proved Koszul (11 algebras)}} \\
\hline
Heisenberg $\cH_\kappa$ & \checkmark & (ii),(ix) & PBW universality \\
\hline
$V_k(\mathfrak{sl}_2)$ (generic $k$) & \checkmark & (ii),(ix) & PBW universality \\
\hline
$\mathrm{Vir}_c$ (generic $c$) & \checkmark & (ii),(ix) & PBW universality \\
\hline
$\mathrm{Vir}_0$ ($c=0$) & \checkmark & (ii) & PBW universality \\
\hline
$\cW_3$ (universal) & \checkmark & (ii) & PBW universality \\
\hline
$\beta\gamma$ ($\lambda=1$) & \checkmark & (ii) & PBW universality \\
\hline
Free fermion $\psi$ & \checkmark & (ii),(ix) & PBW universality \\
\hline
$D_4$ lattice VOA & \checkmark & (ii) & Lattice filtration \\
\hline
Symplectic fermion $\mathcal{SF}=bc_1$ & \checkmark & (ii) & PBW universality \\
\hline
$V_k(\mathfrak{sl}_2)$ (admissible, universal) & \checkmark & (ii) & PBW universality \\
\hline
$\cW_{1+\infty}$ & \checkmark & (ii) & MC4 completion tower \\
\hline
\multicolumn{4}{|c|}{\textit{Not Koszul (2 algebras)}} \\
\hline
Minimal model $L(c_{3,4},0)$ & $\boldsymbol{\times}$ & (ix) fails & Null vector at $h=6$ in bar range \\
\hline
$V_{-2}(\mathfrak{sl}_2)$ (critical) & $\boldsymbol{\times}$ & diagonal criterion fails & Feigin--Frenkel center; PBW page not diagonal \\
\hline
\multicolumn{4}{|c|}{\textit{Open (2 algebras)}} \\
\hline
$L_1(\mathfrak{sl}_2)$ (simple quotient) & ? & &
Li--bar/PBW off-diagonal obstruction complex; null at $h=2$ blocks PBW \\
\hline
Triplet $\cW(2)$ at $c=-2$ & ? & &
fixed-point Li--bar/PBW off-diagonal obstruction complex for
$\mathcal{SF}^{\mathbb Z_2}$; $C_2$-cofinite but non-rational \\
\hline
\end{tabular}
}% end small
\end{table}

\noindent
\textbf{Status column}: \checkmark\ = proved Koszul;
$\boldsymbol{\times}$ = proved not Koszul; ? = open.
\textbf{Equivalence column}: which of the seven intrinsic
bidirectional equivalences
\textup{(}i--v, ix--x\textup{)} of
Theorem~\ref{thm:koszul-equivalences-meta} was used for proof;
(ii) = PBW $E_2$-collapse;
(ix) = Kac--Shapovalov non-degeneracy. Condition~(vi) is a proved
consequence of~(v); (vii) is conditional on the factorization-homology
comparison theorem; (viii) records chiral Hochschild concentration
as a one-way consequence on the Koszul locus.
For the open rows, the named obstruction complex is the
Li--bar/PBW off-diagonal bar-cohomology complex controlled by
Theorem~\ref{thm:kac-shapovalov-koszulness} and
Remark~\ref{rem:li-bar-vs-kac-shapovalov}; the table asserts no
Koszul status until that complex is computed on the listed quotient
or fixed-point surface.

\begin{remark}[Corrections to naive expectations]
\label{rem:koszulness-corrections}
\index{Koszul property!corrections}
Three corrections emerge from the systematic landscape classification.
\begin{enumerate}
\item \emph{Symplectic boson and symplectic fermion are distinct
Koszul free-field points.}
 The point $\beta\gamma_{1/2}$ is the symplectic boson: it has
 central charge $c=-1$ and $\kappa=-1/2$. The symplectic-fermion point
 recorded in this census is the fermionic point $\mathcal{SF}=bc_1$:
 it has central charge $c=-2$ and $\kappa=-1$. PBW universality
 (Proposition~\ref{prop:pbw-universality}) applies to these parent
 free-field algebras. Logarithmic phenomena enter the module category
 and the $\mathbb{Z}_2$ orbifold/triplet construction
 (Remark~\ref{rem:symplectic-logarithmic}); they do not by themselves
 obstruct the parent free-field bar complex.
\item \emph{Quotient of Koszul $\not\Rightarrow$ Koszul.}
 The simple quotient $L_1(\mathfrak{sl}_2)$ has null vector at
 conformal weight $h = 2$, which is in the bar-relevant range
 $h \geq 2$. PBW fails; the question is open. More generally, for
 admissible levels $k = p/q - 2$, the first null vector appears at
 $h = (p{-}1)q$, which is in the bar-relevant range whenever
 $(p{-}1)q \geq 2$.
\item \emph{Critical level separates uncurving from Koszulness.}
 The universal algebra $V_{-h^\vee}(\fg)$ at the critical level
 is uncurved ($\kappa = 0$), and PBW degeneration still records the
 associated graded algebra. Full Koszulness requires diagonal
 concentration; it fails because the Feigin--Frenkel center contributes
 off-diagonal bar cohomology. Critical level breaks the Sugawara
 construction and the generic diagonal criterion at the same time.
\end{enumerate}
\end{remark}

\begin{remark}[Genus-\texorpdfstring{$1$}{1} free energy and the conformal anomaly]\label{rem:genus-1-verification}
\index{free energy!genus one}
The Faber--Pandharipande $\lambda_g$ formula (Theorem~\ref{thm:genus-universality}(iii))
gives the scalar genus-$g$ coefficient
$F_g^{\mathrm{sc}}(\cA)=\kappa(\cA)\cdot
\frac{2^{2g-1}-1}{2^{2g-1}}\cdot \frac{|B_{2g}|}{(2g)!}$.
At $g = 1$, $|B_2| = 1/6$, yielding
\[
F_1(\cA)=F_1^{\mathrm{sc}}(\cA)=\frac{\kappa(\cA)}{24}.
\]
For the Heisenberg algebra at level $\kappa = 1$ (single free boson, $c = 1$),
$F_1(\mathcal{H}_1) = 1/24 = c/24$, precisely the conformal anomaly coefficient
$q^{-c/24}$ in the torus partition function $Z(\tau) = q^{-c/24}\sum_h d(h)\,q^h$.
For $\widehat{\mathfrak{sl}}_2$ at level~$k$:
$F_1(\widehat{\mathfrak{sl}}_{2,k}) = 3(k{+}2)/96 = (k{+}2)/32$.
In general, $F_1(\cA)=\kappa(\cA)/24$ differs from the conformal anomaly
$c(\cA)/24$ by the ratio $\kappa(\cA)/c(\cA)$, which is~$1$ for the
Heisenberg algebra but varies for other algebras.
The displayed equality is a genus-$1$ scalar evaluation. The formula
$F_1(\cA)=\kappa(\cA)/24$ is the genus-$1$ evaluation
of the universal MC class: $\theta_1(\cA)=\kappa(\cA) \cdot \mu
\otimes \lambda_1$
(Theorem~\ref{thm:explicit-theta}(d)), integrated
against the Faber--Pandharipande $\lambda_1$-formula.
The ratio $\kappa(\cA)/c(\cA) = \varrho(\mathfrak{g})$
(Corollary~\ref{cor:genus1-anomaly-ratio}) measures
how much of the Lie bracket's cyclic deformation
survives DS reduction.
\end{remark}

\begin{remark}[Anomaly ratio and the Polyakov formula]\label{rem:anomaly-ratio-polyakov}
\index{anomaly ratio!Polyakov formula}
On the scalar component, the anomaly ratio
$\varrho(\cA) := \kappa(\cA)/c(\cA)$ measures how much of the
conformal anomaly survives from the Lie bracket's cyclic deformation.
For the rank-$d$ Heisenberg algebra $\cH_1^{\oplus d}$ at level~$1$,
$\varrho = d/d = 1$: the Polyakov formula $F_1 = c/24$ and the shadow
formula $F_1(\cH_1^{\oplus d})=\kappa(\cH_1^{\oplus d})/24$ coincide
because $\kappa(\cH_1^{\oplus d})=c(\cH_1^{\oplus d})=d$. For all
other standard families, $\varrho \neq 1$:
\begin{center}
\begin{tabular}{lccc}
Family & $c(\cA)$ & $\kappa(\cA)$ & $\varrho(\cA)=\kappa(\cA)/c(\cA)$ \\ \hline
$\mathcal{H}_1^{\oplus d}$ \textup{($d$ bosons)} & $d$ & $d$ & $1$ \\
$\widehat{\mathfrak{sl}}_{2,k}$ & $3k/(k{+}2)$ & $3(k{+}2)/4$ & $(k{+}2)^2/(4k)$ \\
$\beta\gamma_{\lambda=1}$ & $2$ & $1$ & $1/2$ \\
$\mathrm{Vir}_c$ & $c$ & $c/2$ & $1/2$ \\
$\mathcal{W}_{3,c}$\textsuperscript{$\mathrm{mod/comp}$}
 & $c$ & $5c/6$ & $5/6$ \\
\end{tabular}
\end{center}
The principal $\mathcal W_3$ entry is a statement on
$H_{\mathrm{diag}}^{g=1}+H_{W_3}^{\mathrm{DS/bar}}$.
The ratio $\varrho(\mathrm{Vir}_c) = 1/2$ means that half the conformal anomaly is absorbed by the nonlinearity of the Virasoro OPE. For $\widehat{\mathfrak{sl}}_{2}$ at level~$k$, $\varrho = (k{+}2)^2/(4k) \sim k/4$ as $k \to \infty$ (linearly divergent) and also diverges at $k \to 0$. The scalar genus-$g$ coefficient
$F_g^{\mathrm{sc}}(\cA)=\kappa(\cA) \cdot \lambda_g^{\mathrm{FP}}$
carries the modular characteristic~$\kappa(\cA)$, not the central
charge~$c(\cA)$; at genus~$1$,
$F_1(\cA)=F_1^{\mathrm{sc}}(\cA)=\kappa(\cA)/24$.
\end{remark}

\begin{corollary}[Anomaly ratio and DS reduction; \ClaimStatusConditional]
\label{cor:anomaly-ratio-ds}
\index{anomaly ratio!DS reduction}
\textbf{Type signature.}
$(\mathrm{Open},\mathrm{principal\ DS/bar},\mathrm{level}=1\to5,
H_{\mathrm{diag}}^{g=1}+H_{W(\fg)}^{\mathrm{DS/bar}})$.
For $\fg=\mathfrak{sl}_N$ this package is
$H_{\mathrm{diag}}^{g=1}+H_{W_N}^{\mathrm{DS/bar}}$; at $N=3$ it is
$H_{\mathrm{diag}}^{g=1}+H_{W_3}^{\mathrm{DS/bar}}$.
On this hypothesis package, every principal $\mathcal{W}$-algebra
$\mathcal{W}^k(\fg)$ of a simple Lie algebra~$\fg$ with exponents
$m_1, \ldots, m_r$ has level-independent anomaly ratio
$\varrho(\mathcal{W}^k(\fg)) = \sum_{i=1}^r 1/(m_i + 1)$. In particular:
$\varrho(\mathrm{Vir}) = 1/(1+1) = 1/2$,
$\varrho(\mathcal{W}_3) = 1/2 + 1/3 = 5/6$,
$\varrho(\mathcal{W}_4) = 1/2 + 1/3 + 1/4 = 13/12$.
The ratio $\varrho > 1$ for $\mathcal{W}_N$ with $N \geq 4$: the modular characteristic exceeds the central charge. The ratio $\varrho(\mathcal{W}_\infty) = \sum_{m=1}^\infty 1/(m+1) = \infty$: the divergence is the harmonic series.
\end{corollary}

\begin{proof}
On
$H_{\mathrm{diag}}^{g=1}+H_{W(\fg)}^{\mathrm{DS/bar}}$,
Theorem~\ref{thm:genus-universality}(ii) gives
$\kappa(\mathcal{W}^k(\fg))
 = c(\mathcal{W}^k(\fg))\cdot \sum_{i=1}^r 1/(m_i+1)$.
Since
$\varrho(\mathcal{W}^k(\fg))
 = \kappa(\mathcal{W}^k(\fg))/c(\mathcal{W}^k(\fg))$,
the ratio depends only on the exponents of~$\fg$, not on~$k$.
The special values follow by direct substitution.
\end{proof}

\begin{corollary}[Genus-\texorpdfstring{$1$}{1} scalar coefficient and anomaly ratio; \ClaimStatusConditional]
\label{cor:genus1-anomaly-ratio}
\index{anomaly ratio!genus-1 free energy}
\textbf{Type signature.}
$(\mathrm{Open},\mathrm{principal\ DS/bar},\mathrm{level}=1\to5,
H_{\mathrm{diag}}^{g=1}+H_{W(\mathfrak g)}^{\mathrm{DS/bar}})$.
For $\mathfrak g=\mathfrak{sl}_N$ this package is
$H_{\mathrm{diag}}^{g=1}+H_{W_N}^{\mathrm{DS/bar}}$; at $N=3$ it is
$H_{\mathrm{diag}}^{g=1}+H_{W_3}^{\mathrm{DS/bar}}$.
For any principal $\mathcal{W}$-algebra
$\mathcal{W}^k(\mathfrak{g})$ with
central charge $c(\mathcal{W}^k(\mathfrak{g})) \neq 0$, the genus-$1$
scalar coefficient on this hypothesis package satisfies
\begin{equation}\label{eq:F1-over-c}
\frac{F_1(\mathcal{W}^k(\mathfrak{g}))}
 {c(\mathcal{W}^k(\mathfrak{g}))}
 = \frac{\varrho(\mathfrak{g})}{24}
\end{equation}
where $\varrho(\mathfrak{g}) = \sum_{i=1}^r \frac{1}{m_i+1}$ is the anomaly ratio (the sum over exponents $m_1, \ldots, m_r$ of~$\mathfrak{g}$). This is $1/2$ for $\mathfrak{sl}_2$ (Virasoro), $5/6$ for $\mathfrak{sl}_3$ ($\mathcal{W}_3$), and $H_N - 1$ for $\mathfrak{sl}_N$ ($\mathcal{W}_N$). The identity follows from
$F_1(\mathcal W^k(\mathfrak{g}))=\kappa(\mathcal W^k(\mathfrak{g}))/24$
(Remark~\ref{rem:genus-1-verification}) and
$\kappa(\mathcal W^k(\mathfrak{g}))
 = c(\mathcal W^k(\mathfrak{g})) \cdot \varrho(\mathfrak{g})$
(Theorem~\ref{thm:genus-universality}\textup{(ii)}).
\end{corollary}

\begin{proof}
On
$H_{\mathrm{diag}}^{g=1}+H_{W(\mathfrak g)}^{\mathrm{DS/bar}}$,
combine
$F_1(\mathcal W^k(\mathfrak{g}))=\kappa(\mathcal W^k(\mathfrak{g}))/24$
(Remark~\ref{rem:genus-1-verification}) with
$\kappa(\mathcal W^k(\mathfrak{g}))
 = c(\mathcal W^k(\mathfrak{g})) \cdot \varrho(\mathfrak{g})$
(Theorem~\ref{thm:genus-universality}(ii)).
\end{proof}

\begin{remark}[Reading the table]\label{rem:reading-master-table}
Three structural features are visible from the table:
\begin{enumerate}
\item \emph{Complementarity.}
 Every $c + c'$ entry is independent of the level~$k$: it depends
 only on the root datum of~$\fg$.
 This is the content of Theorem~\ref{thm:central-charge-complementarity}.
\item \emph{Critical level.}
 Setting $\kappa(\cA) = 0$ (i.e., $t = k + h^\vee = 0$)
 gives $\mathfrak O_g^K=0$ under~$H_D^K$ and
 $F_g^{\mathrm{sc}}=0$ under~$H_D^{\mathrm{graph}}$.
 The bar complex becomes \emph{uncurved}, and
 the Feigin--Frenkel center $\mathfrak{z}(\widehat{\fg}_{-h^\vee})$
 emerges as the degree-zero cohomology.
\item \emph{Free fields.}
 The exact Koszul dualities ($\psi \leftrightarrow \gamma$,
 $bc \leftrightarrow \beta\gamma$) satisfy
 $\kappa(\cA) + \kappa(\cA^!) = 0$
 (Theorem~\ref{thm:fermion-all-genera};
 not to be confused with bosonization,
 Remark~\ref{rem:bosonization-not-koszul}): the individual
 obstruction coefficients $\kappa(\psi) = \tfrac{1}{4}$,
 $\kappa(bc) = c/2$ are nonzero (they measure the
 Virasoro subalgebra curvature), but the tensor product
 $bc \otimes \beta\gamma$ is uncurved.
 The Heisenberg is the abelian curved case, with
 $\kappa = \kappa$ recording the central extension.
\item \emph{Gaps.}
 The Heisenberg entry $c + c' = \text{n/a}$ reflects the fact that
 $\mathcal{H}_\kappa^! = \mathrm{Sym}^{\mathrm{ch}}(V^*)$ is
 curved (not a vertex algebra in the traditional sense), so
 its central charge is not defined by the Sugawara construction.
 The Yangian entries are blank because $Y(\fg)$ is an $\Eone$-chiral
 algebra, not conformally graded: it has no Virasoro subalgebra
 and hence no central charge or obstruction coefficient in the
 sense of the other rows.
\end{enumerate}
\end{remark}

\begin{theorem}[Witness-family complementarity from Verdier duality;
\ClaimStatusConditional]
\label{thm:census-witness-complementarity}
\index{complementarity!witness-family derivation|textbf}
\index{Verdier duality!witness-family complementarity}
\textbf{Type signature.}
$(\mathrm{Open},\mathrm{bar/Verdier\ and\ principal\ DS},
\mathrm{level}=1\to5,
H_{\mathrm{diag}}^{g=1}+H_{W_3}^{\mathrm{DS/bar}}
\text{ for the principal }\mathcal W_3\text{ clause})$.
For the four witness families of the shadow classes
\[
 \cH_k,\qquad
 \widehat{\fg}_k,\qquad
 \beta\gamma_\lambda,\qquad
 \mathrm{Vir}_c,
\]
their Koszul dual partners determined by Verdier duality on the bar
coalgebra are represented on the scalar component by
\[
 \cH_k^!=(\mathrm{Sym}^{\mathrm{ch}}(V^*[1]),m_0=-k\omega),\qquad
 \kappa(\widehat{\fg}_k^!)
 = \kappa(\widehat{\fg}_{-k-2h^\vee}),\qquad
 (\beta\gamma_\lambda)^! \simeq bc_\lambda,\qquad
 \mathrm{Vir}_c^! \simeq \mathrm{Vir}_{26-c}.
\]
The corresponding dual obstruction coefficients are
\[
 \kappa(\cH_k^!) = -k,\qquad
 \kappa(\widehat{\fg}_k^!)
 = -\frac{\dim(\fg)(k+h^\vee)}{2h^\vee},\qquad
 \kappa((\beta\gamma_\lambda)^!)
 = -(6\lambda^2 - 6\lambda + 1),\qquad
 \kappa(\mathrm{Vir}_c^!) = \frac{26-c}{2}.
\]
Hence the witness-family complementarity sums are
\[
 \kappa(\cH_k) + \kappa(\cH_k^!) = 0,\qquad
 \kappa(\widehat{\fg}_k) + \kappa(\widehat{\fg}_k^!) = 0,\qquad
 \kappa(\beta\gamma_\lambda) + \kappa((\beta\gamma_\lambda)^!) = 0,\qquad
 \kappa(\mathrm{Vir}_c) + \kappa(\mathrm{Vir}_c^!) = 13.
\]
On
$H_{\mathrm{diag}}^{g=1}+H_{W_3}^{\mathrm{DS/bar}}$, the principal
$\mathcal{W}_3$ modular conductor is
\[
 \kappa(\mathcal{W}_3^k) + \kappa(\mathcal{W}_3^{k'})
 = \frac{250}{3}.
\]
For Bershadsky--Polyakov, the exact companion identity is the
central-charge statement
\[
 c_{\mathrm{BP}}(k)+c_{\mathrm{BP}}(-k-6)=50.
\]
The BP modular sum
\[
 K^\kappa_{\mathrm{BP}}(k)
 :=\kappa(\mathrm{BP}_k)+\kappa(\mathrm{BP}_{-k-6})
\]
has status \ClaimStatusOpen pending the direct genus-$1$
non-separating curvature computation.  On
$H_{\mathrm{BP}}^{\mathrm{DS/bar}}$, the companion algebra
$\mathrm{BP}_{-k-6}$ represents the Verdier--Koszul partner.  The
unconditional scalar complementarity set of the four witness families
is $\{0,13\}$.  The package
$H_{\mathrm{diag}}^{g=1}+H_{W_3}^{\mathrm{DS/bar}}$ supplies the
conditional principal extension $\{0,13,250/3\}$.  The BP modular
slot has status \ClaimStatusOpen within the coarse archetype
class~$\mathbf{G}/\mathbf{L}/\mathbf{C}/\mathbf{M}$.
\end{theorem}

\begin{proof}
For Heisenberg, Theorem~\ref{thm:heisenberg-koszul-dual-early}
identifies the Verdier dual partner with the curved symmetric chiral
algebra $\mathrm{Sym}^{\mathrm{ch}}(V^*)$ carrying curvature
$m_0 = -k \cdot \mathbf{1}$. Since $\kappa$ is the scalar coefficient
of the genus-$1$ curvature, this gives
$\kappa(\cH_k^!) = -k$ and hence
$\kappa(\cH_k) + \kappa(\cH_k^!) = k + (-k) = 0$.
The sign is checked three ways on the witness locus:
directly from the curvature term, at the uncurved point $k=0$, and by
additivity, which reproduces the lattice identity
$\kappa(V_\Lambda^!) = -\operatorname{rank}(\Lambda)$ for rank-$1$
summands.

For affine Kac--Moody, Verdier duality gives the Feigin--Frenkel level
involution $k' = -k - 2h^\vee$, so
\[
 \kappa(\widehat{\fg}_k^!)
 =
 \frac{\dim(\fg)(k' + h^\vee)}{2h^\vee}
 =
 \frac{\dim(\fg)(-k-h^\vee)}{2h^\vee}
 =
 -\frac{\dim(\fg)(k+h^\vee)}{2h^\vee}
 =
 -\kappa(\widehat{\fg}_k).
\]
This is confirmed by the critical check $k=-h^\vee$, where both sides
vanish, and by the concrete specialization
$\fg=\mathfrak{sl}_2$, where
$3(k+2)/4 + 3(-k-2)/4 = 0$.

For $\beta\gamma_\lambda$, Theorem~\ref{thm:betagamma-bc-koszul}
identifies the Verdier dual partner with $bc_\lambda$.
Proposition~\ref{prop:betagamma-bc-koszul-detailed} gives
\[
 c_{\beta\gamma} = 2(6\lambda^2 - 6\lambda + 1),
 \qquad
 c_{bc} = -2(6\lambda^2 - 6\lambda + 1),
\]
so
\[
 \kappa((\beta\gamma_\lambda)^!)
 =
 \kappa(bc_\lambda)
 =
 \frac{c_{bc}}{2}
 =
 -(6\lambda^2 - 6\lambda + 1)
 =
 -\kappa(\beta\gamma_\lambda).
\]
Three independent checks are immediate: the bosonic symmetry
$\lambda \leftrightarrow 1-\lambda$ leaves
$\kappa(\beta\gamma_\lambda)$ unchanged but does not change the dual
partner; at $\lambda=\tfrac{1}{2}$ one gets
$-\tfrac{1}{2} + \tfrac{1}{2} = 0$; and at $\lambda=1$ one gets
$1 + (-1) = 0$.

For Virasoro, Proposition~\ref{prop:virasoro-generic-koszul-dual}
gives the same-family Verdier partner
$\mathrm{Vir}_{26-c}$, hence
\[
 \kappa(\mathrm{Vir}_c^!)
 =
 \kappa(\mathrm{Vir}_{26-c})
 =
 \frac{26-c}{2},
 \qquad
 \kappa(\mathrm{Vir}_c) + \kappa(\mathrm{Vir}_c^!)
 =
 \frac{c}{2} + \frac{26-c}{2}
 =
 13.
\]
This is checked again at the self-dual point $c=13$, where the sum is
$2 \cdot 13/2 = 13$, and at the two uncurved endpoints
$c=0$ and $c=26$, where one member of the pair has vanishing scalar
curvature.

On
$H_{\mathrm{diag}}^{g=1}+H_{W_3}^{\mathrm{DS/bar}}$,
Corollary~\ref{cor:general-w-obstruction} gives
$\kappa(\mathcal{W}_3) = \tfrac{5}{6}c$ and
the exact rational-function arithmetic of
Theorem~\ref{thm:central-charge-complementarity} gives
$c+c' = 100$.  Therefore
\[
 \kappa(\mathcal{W}_3^k) + \kappa(\mathcal{W}_3^{k'})
 =
 \frac{5}{6}(c+c')
 =
 \frac{5}{6}\cdot 100
 =
 \frac{250}{3}.
\]
The same number is recovered from the anomaly-ratio route
$(H_3-1)K_3^c = \tfrac{5}{6}\cdot 100$ and from the self-dual check
$2\kappa(\mathcal{W}_{3,c_*}) = 2 \cdot \tfrac{5}{6}\cdot 50$ at
$c_* = 50$.

For Bershadsky--Polyakov, Fehily--Kawasetsu--Ridout define an ordinary
bosonic vertex algebra with even generators $J,G^+,G^-,T$
\textup{(}\cite[\S2.1]{FKR20}\textup{)}.  The source-correct
reciprocal-weight diagnostic is therefore $17/6$.  Theorem~\ref{thm:bp-koszul-conductor-polynomial}
gives the exact rational-function identity
\[
 c_{\mathrm{BP}}(k)+c_{\mathrm{BP}}(-k-6)=50.
\]
The genus-$1$ modular characteristic and its companion sum have status
\ClaimStatusOpen.  Their proof obligation is a direct curvature
calculation of the non-separating BP modular trace, including the
charged ghosts, neutral fields, improvement term, and mixed channels.
Under $H_{\mathrm{BP}}^{\mathrm{DS/bar}}$,
Proposition~\ref{prop:bp-self-duality} identifies the companion algebra
with the Verdier--Koszul partner.  The exact central-charge midpoint is
$25$; its regular fibre consists of $k=-3\pm2i$, exchanged by the
companion involution, while the fixed level $k=-3$ is the pole.

Therefore the four witness families give the unconditional values
$0,0,0,13$ and the set $\{0,13\}$.  On
$H_{\mathrm{diag}}^{g=1}+H_{W_3}^{\mathrm{DS/bar}}$, the principal
$\mathcal{W}_3$ family contributes $250/3$ and enlarges the set to
$\{0,13,250/3\}$.  The BP family contributes an open modular slot and
the exact central conductor~$50$.  The scalar complementarity constant
is family-dependent.
\end{proof}

\begin{proposition}[Archetype-by-archetype
$\kappa + \kappa^!$ arithmetic with anomaly-ratio bridge;
\ClaimStatusConditional]
\label{prop:archetype-complementarity-bridge}
\index{complementarity!archetype-by-archetype}
\index{Trinity conductor!anomaly-ratio bridge}
\index{kappa plus kappa dual@$\kappa{+}\kappa^!$!archetype table}
\textup{[Regime: curved-central, Koszul locus;
Convention~\textup{\ref{conv:regime-tags}};
principal $\mathcal W_3$ package
$H_{\mathrm{diag}}^{g=1}+H_{W_3}^{\mathrm{DS/bar}}$.]}
The scalar complementarity sum
$K^\kappa(\cA) := \kappa(\cA) + \kappa(\cA^!)$ and the Trinity
conductor $K(\cA) := c(\cA) + c(\cA^!)$ are related, on the witness
landscape and those $\mathcal{W}$-family extensions with an established
level-independent anomaly ratio, together with the $\mathcal B$-family
candidate under~$H_{\mathsf B}$, by the
\emph{anomaly-ratio bridge}
\begin{equation}\label{eq:kappa-sum-rho-K-bridge}
 K^\kappa(\cA) \;=\; \varrho(\cA) \cdot K(\cA),
 \qquad
 \varrho(\cA) \;:=\; \frac{\kappa(\cA)}{c(\cA)},
\end{equation}
where $\varrho(\cA)$ is the level-independent anomaly ratio of
Corollary~\ref{cor:anomaly-ratio-ds}.  The four shadow-archetype
witnesses give the unconditional scalar set $\{0,13\}$.  On
$H_{\mathrm{diag}}^{g=1}+H_{W_3}^{\mathrm{DS/bar}}$, the principal
$\mathsf{M}$-family extension adds $250/3$.  Under $H_{\mathsf B}$,
the $\mathcal B$-family candidate adds~$8$.  Thus the joint conditional
list is $\{0,8,13,250/3\}$.  The BP row
records its exact central conductor and open modular entries in the table
below.  Each displayed scalar has the status stated in its row.
\begin{center}
\renewcommand{\arraystretch}{1.35}
\begin{tabular}{|c|c|c|c|c|c|}
\hline
Archetype & Witness $\cA$ & $\varrho(\cA)$ & $K(\cA) = c + c^!$
 & $K^\kappa(\cA) = \kappa + \kappa^!$
 & scalar midpoint $\kappa^{\mathrm{mid}}=K^\kappa/2$ \\
\hline
$\mathsf{G}$ & $\cH_k$ (Heisenberg) & $1$ & $0$ & $0$ & $0$ \\
\hline
$\mathsf{L}$ & $\widehat{\fg}_k$ (affine KM) & $0$ & $2\dim\fg$ & $0$ & $0$ \\
\hline
$\mathsf{C}$ & $\beta\gamma_\lambda$ & $1/2$ & $0$ & $0$ & $0$\\
\hline
$\mathsf{M}$ & $\mathrm{Vir}_c$ & $1/2$ & $26$ & $13$ & $13/2$ \\
\hline
\hline
$\mathsf{M}$-ext & $\mathcal{W}_3^k$ (principal)\textsuperscript{$\mathrm{CD}_{W_3}$}
 & $5/6$ & $100$\textsuperscript{exact} & $250/3$ & $125/3$ \\
\hline
$\mathsf{M}$-ext & $\mathrm{BP}_k$ (minimal) & open & $50$ & open & open \\
\hline
\hline
$\mathcal{B}$-fam & $\mathcal{H}_{\mathrm{Muk}}(K3)$ (Mukai-enhanced K3)
 & $1/6$\textsuperscript{$\mathrm{CJ}_{\mathsf B}$}
 & $48$\textsuperscript{$\mathrm{CJ}_{\mathsf B}$}
 & $8$\textsuperscript{$\mathrm{CJ}_{\mathsf B}$}
 & $4$\textsuperscript{$\mathrm{CJ}_{\mathsf B}$} \\
\hline
\end{tabular}
\end{center}
Here $\mathrm{CJ}_{\mathsf B}$ means \ClaimStatusConjectured under
$H_{\mathsf B}$.
The marker $\mathrm{CD}_{W_3}$ means that the central entry~$100$ is
exact and the entries $5/6$, $250/3$, and $125/3$ are conditional on
$H_{\mathrm{diag}}^{g=1}+H_{W_3}^{\mathrm{DS/bar}}$.
In particular, the bridge $K^\kappa=\varrho K$ gives the
unconditional set $\{0,13\}$ on the four Vol~I witnesses.  The
principal package
$H_{\mathrm{diag}}^{g=1}+H_{W_3}^{\mathrm{DS/bar}}$ gives the
conditional extension $\{0,13,250/3\}$.  Together with the
$\mathcal B$-family candidate under~$H_{\mathsf B}$, the conditional
ledger is $\{0,8,13,250/3\}$.  The four unconditional pairs use
$\varrho\in\{1,0,1/2,1/2\}$ and
$K\in\{0,2\dim\fg,0,26\}$; the principal conditional pair is
$(5/6,100)$ and the $\mathcal B$ conditional pair is $(1/6,48)$.
The seventh, BP, row has exact
$K^c_{\mathrm{BP}}=50$ and status \ClaimStatusOpen in its $\varrho$,
$K^\kappa$, and modular-midpoint entries.  The $\mathcal B$-family
arithmetic is proved: the Mukai lattice has rank~$24$, signature~$(4,20)$,
$c_+=4$, and $2c_+=8$.  The table interprets the candidate tuple
$(\varrho,K,K^\kappa,\kappa^{\mathrm{mid}})=(1/6,48,8,4)$ under
\[
H_{\mathsf B}
=(H_{\mathrm{chart}},H_{\mathrm{KD}},H_{\mathrm{scalar}},
H_{\mathrm{mod}},H_{\mathrm{quantum}}).
\]
Here $H_{\mathrm{mod}}$ constructs the Humbert--Borcherds comparison,
and $H_{\mathrm{quantum}}$ constructs the functor selecting root
order~$8$.  Bruinier's coefficient-field lemma and Lusztig's
root-of-unity construction supply the respective ambient arithmetic
and quantum settings; the two comparison hypotheses supply their link
to the Mukai chiral chart.
\end{proposition}

\begin{proof}
The bridge~\eqref{eq:kappa-sum-rho-K-bridge} is an identity of scalars
once the displayed chiral data and comparison package are fixed. Assume $\varrho$ is
level-independent on the family $\mathcal{F}$, so
$\varrho(\cA) = \kappa/c = \varrho(\cA^!) = \kappa^!/c^!$. Then
$\kappa + \kappa^! = \varrho(c + c^!)$; level-independence of
$\varrho$ on the standard families is
Corollary~\ref{cor:anomaly-ratio-ds} for
$\mathcal{W}^k(\fg)$ and is checked directly on the other families.

The four witness rows are computed individually.  The principal
$\mathcal W_3$ row is conditional on
$H_{\mathrm{diag}}^{g=1}+H_{W_3}^{\mathrm{DS/bar}}$.  The
$\mathcal B$ row records its conditional value under~$H_{\mathsf B}$,
and the BP row records its exact central conductor together with the
open modular obligation.

\emph{Row $\mathsf{G}$ (Heisenberg
$\cH_k \equiv \cH_1^{\oplus k}$, rank-$k$ free-boson
convention).} The OPE of the rank-$k$ Heisenberg
algebra $J^a(z) J^b(w) \sim \delta^{ab}/(z-w)^2$ gives
$c(\cH_k) = k$ and $\kappa(\cH_k) = k$
\textup{(}Remark~\ref{rem:anomaly-ratio-polyakov} with
$d = k$\textup{)}, hence
$\varrho(\cH_k) = 1$. The Verdier dual
$\cH_k^!=(\mathrm{Sym}^{\mathrm{ch}}(V^*[1]),m_0=-k\omega)$
\textup{(}Theorem~\ref{thm:heisenberg-koszul-dual-early}\textup{)}
carries $m_0 = -k \cdot \mathbf{1}$, so
$\kappa(\cH_k^!) = -k$ and $c(\cH_k^!) = -k$ at the level of the
scalar central-charge projection of the curved symmetric chiral
algebra. The Trinity conductor is $K = k + (-k) = 0$, and the bridge
gives $K^\kappa = 1 \cdot 0 = 0$; the sum
$\kappa + \kappa^! = k + (-k) = 0$ is the $\mathsf{G}$-archetype
witness of the free-field antisymmetric case of
Theorem~\ref{thm:fermion-all-genera}.

\emph{Row $\mathsf{L}$ (non-critical affine Kac--Moody
$\widehat{\fg}_k$).} The
critical-level reflection \(k'=-k-2h^\vee\) sends
$c(k) = k\dim\fg/(k+h^\vee)$ to $c(k') = (k+2h^\vee)\dim\fg/(k+h^\vee)$,
hence $c + c' = 2\dim\fg$
\textup{(}Theorem~\ref{thm:central-charge-complementarity}(a)\textup{)}.
The modular characteristic $\kappa(\widehat{\fg}_k) = \dim\fg\,(k+h^\vee)/(2h^\vee)$
has Koszul-dual image
$\kappa(\widehat{\fg}_{k'}) = -\kappa(\widehat{\fg}_k)$, so
$K^\kappa = 0$. The anomaly ratio
$\varrho(\widehat{\fg}_k) = [(k+h^\vee)/(2h^\vee)] \cdot [(k+h^\vee)/k]$
depends on~$k$, violating level-independence on the unreduced algebra;
the bridge applies in the \emph{level-symmetrised} form
$\varrho_{\mathrm{sym}}(\widehat{\fg}) = 0$ defined by
$(\kappa + \kappa')/(c + c') = 0/(2\dim\fg) = 0$, consistent with the
antisymmetric Feigin--Frenkel structure. Equivalently: on the
$\mathsf{L}$-archetype the sum $\kappa + \kappa^!$ vanishes not because
$\varrho = 0$ pointwise but because the level-pair symmetrisation sends
$\kappa \mapsto -\kappa$. We record this with
$\varrho = 0$ in the table and note the locus
caveat. At $k=-h^\vee$ the conformal central charge formula is not
defined and the row leaves the class-L Koszul locus.

\emph{Row $\mathsf{C}$ ($\beta\gamma_\lambda$).} The conformal-weight
family has $c(\beta\gamma_\lambda) = 2(6\lambda^2 - 6\lambda + 1)$
\textup{(}Sugawara-level computation,
Proposition~\ref{prop:betagamma-bc-koszul-detailed}\textup{)} and Verdier
dual $bc_\lambda$ with $c(bc_\lambda) = -c(\beta\gamma_\lambda)$
by Fock-space parity flip (bosonic versus fermionic reconstruction of
the same $\lambda$-twisted conformal structure). Hence $K = 0$,
$\kappa + \kappa^! = 0$, and $\varrho = 1/2$ by
$\kappa(\beta\gamma_\lambda) = c/2 = 6\lambda^2 - 6\lambda + 1$. The
bridge gives $K^\kappa = (1/2) \cdot 0 = 0$.

\emph{Row $\mathsf{M}$ (Virasoro $\mathrm{Vir}_c$).} The
same-family Verdier dual $\mathrm{Vir}_c^! \simeq \mathrm{Vir}_{26-c}$
\textup{(}Proposition~\ref{prop:virasoro-generic-koszul-dual}, at
$\mathsf{M}/\mathsf{S}$ locus\textup{)} gives $c + c' = 26$ and
$\kappa(\mathrm{Vir}_c) = c/2$, so
$\varrho(\mathrm{Vir}_c) = 1/2$ and the bridge gives
$K^\kappa = (1/2) \cdot 26 = 13$.

\emph{Row $\mathsf{M}$-ext (principal $\mathcal{W}_3^k$).}
The central-charge reflection gives the exact identity
$K^c(\mathcal{W}_3)=100$
\textup{(}Theorem~\ref{thm:central-charge-complementarity}(b) at
$\fg = \mathfrak{sl}_3$, $2r + 4h^\vee d = 4 + 96 = 100$\textup{)}.
On
$H_{\mathrm{diag}}^{g=1}+H_{W_3}^{\mathrm{DS/bar}}$, the normalized
diagonal trace and the principal bar companion give
$\kappa(\mathcal{W}_3^k) = (5/6) c(\mathcal{W}_3^k)$
\textup{(}Corollary~\ref{cor:anomaly-ratio-ds} with
$\varrho(\mathcal{W}_3) = 1/2 + 1/3 = 5/6$\textup{)}.
The bridge then gives the conditional modular conductor
$K^\kappa = (5/6) \cdot 100 = 500/6 = 250/3$.

\emph{Row $\mathsf{M}$-ext (Bershadsky--Polyakov $\mathrm{BP}_k$).}
The standard FKR central charge gives the exact conductor
$K(\mathrm{BP})=50$
\textup{(}Theorem~\ref{thm:bp-koszul-conductor-polynomial}\textup{)}.
The FKR strong generators $J,G^+,G^-,T$ are even, and their
reciprocal-weight diagnostic is $17/6$.  A direct genus-$1$ curvature
calculation determines $\kappa(\mathrm{BP}_k)$ and the anomaly ratio.
Thus the bridge entries $\varrho(\mathrm{BP})$ and
$K^\kappa(\mathrm{BP})$ have status \ClaimStatusOpen.  On
$H_{\mathrm{BP}}^{\mathrm{DS/bar}}$,
Proposition~\ref{prop:bp-self-duality} identifies the central companion
pair with the same-family Verdier--Koszul pair.

\emph{Scalar midpoints.} Define
$\kappa^{\mathrm{mid}}(\cA)=K^\kappa(\cA)/2$; halving gives the last
column.  At a regular same-family fixed point this midpoint is the
value of~$\kappa$.  For BP, the modular midpoint is open.  Its exact
central-charge midpoint is~$25$, attained on the regular companion
fibre $k=-3\pm2i$, while the fixed level $k=-3$ is the pole.

Consequently the four unconditional Vol~I pairs
$(\varrho,K)\in\{(1,0),(0,2\dim\fg),(1/2,0),(1/2,26)\}$
give the scalar set $\{0,13\}$.  This is the
\emph{Vol~I Theorem~C scalar ceiling} on the four shadow-archetype
classes $\mathsf{G}/\mathsf{L}/\mathsf{C}/\mathsf{M}$
\textup{(}Heisenberg, affine Kac--Moody, $\beta\gamma$, Virasoro\textup{)}.
On $H_{\mathrm{diag}}^{g=1}+H_{W_3}^{\mathrm{DS/bar}}$, the principal
pair $(5/6,100)$ enlarges the comparison set to
$\{0,13,250/3\}$.  The Mukai arithmetic gives
$c_+=4$, $\operatorname{rk}=24$, and $2c_+=8$ directly.  Under
$H_{\mathsf B}$ together with
$H_{\mathrm{diag}}^{g=1}+H_{W_3}^{\mathrm{DS/bar}}$, the candidate
pair $(\varrho,K)=(1/6,48)$ gives the conditional numerical ledger
$\{0,8,13,250/3\}$.  The five components of
$H_{\mathsf B}$ are precisely the remaining chiral-chart, Koszul,
scalar, modular, and quantum comparison obligations.  Bruinier's
coefficient-field modulus $\operatorname{lcm}(N,8)$ and Lusztig's
chosen root order are two separately typed ambient inputs for the last
two obligations; their comparison with the Mukai value~$8$ has status
\ClaimStatusConjectured.  The BP modular slot remains
\ClaimStatusOpen and is represented symbolically in both ledgers.
\end{proof}

\begin{proposition}[Mukai lattice arithmetic and the $\mathsf B$-family
anomaly-ratio proposal; \ClaimStatusConjectured]
\label{prop:G-B-heisenberg-rho-bifurcation}
\index{anomaly ratio!$\mathsf G$ vs $\mathsf B$ bifurcation}
\index{Heisenberg!two conventions (positive-definite vs Lorentzian)}
\textup{[Regime: curved-central, Koszul locus;
Convention~\textup{\ref{conv:regime-tags}};
hypothesis package $H_{\mathsf B}=(H_{\mathrm{chart}},H_{\mathrm{KD}},
H_{\mathrm{scalar}},H_{\mathrm{mod}},H_{\mathrm{quantum}})$.]}
The Mukai lattice
$\Lambda_{\mathrm{Muk}}=\widetilde H(K3,\mathbb Z)$ is even unimodular
of rank~$24$ and signature~$(4,20)$.  Hence
\[
 c_+(\Lambda_{\mathrm{Muk}})=4,
 \qquad
 \frac{c_+(\Lambda_{\mathrm{Muk}})}{\operatorname{rk}(\Lambda_{\mathrm{Muk}})}
 =\frac16,
 \qquad
 2c_+(\Lambda_{\mathrm{Muk}})=8.
\]
The positive-definite $\mathsf G$-Heisenberg row has the proved tuple
$(\varrho,K,K^\kappa)=(1,0,0)$.  Under~$H_{\mathsf B}$, the Mukai
arithmetic determines the proposed $\mathsf B$-family tuple
\[
 (\varrho,K,K^\kappa,\kappa^{\mathrm{mid}})
 =\left(\frac16,48,8,4\right),
 \qquad K^\kappa=\varrho K.
\]
The lattice identities in the first display have status
\ClaimStatusProvedHere; their chiral-conductor interpretation in the
second display has status \ClaimStatusConjectured under~$H_{\mathsf B}$.
\end{proposition}

\begin{proof}
The K3 intersection lattice is
$H^2(K3,\mathbb Z)\simeq U^3\oplus E_8(-1)^2$, of rank~$22$ and
signature~$(3,19)$.  The summand
$H^0(K3,\mathbb Z)\oplus H^4(K3,\mathbb Z)\simeq U$ adds one positive
and one negative direction.  Therefore
\[
 \widetilde H(K3,\mathbb Z)
 \simeq U^4\oplus E_8(-1)^2,
 \qquad
 \operatorname{sig}\widetilde H(K3,\mathbb Z)=(4,20),
\]
and the three arithmetic values follow.  This is the Mukai-lattice
calculation of~\cite{Mukai1987}.

For the $\mathsf G$ row,
Theorem~\ref{thm:heisenberg-koszul-dual-early} gives the curved
Verdier partner with scalar curvature~$-k$; hence
$(\varrho,K,K^\kappa)=(1,0,0)$.

For the $\mathsf B$ row, $H_{\mathrm{chart}}$ constructs an augmented
chiral chart from $\Lambda_{\mathrm{Muk}}$;
$H_{\mathrm{KD}}$ identifies its Koszul partner;
$H_{\mathrm{scalar}}$ identifies the companion sum with~$2c_+$;
$H_{\mathrm{mod}}$ constructs the Humbert--Borcherds comparison; and
$H_{\mathrm{quantum}}$ constructs the functor selecting root order~$8$.
These hypotheses turn the proved arithmetic
$(c_+/\operatorname{rk},2\operatorname{rk},2c_+)
=(1/6,48,8)$ into the displayed chiral tuple.  Thus the proposition
records a proved lattice calculation and its precisely typed
conjectural chiral realization.
\end{proof}

\begin{remark}[Three-path computation of the archetype table]
\label{rem:archetype-complementarity-three-path}
\index{complementarity!archetype table!three-path computation}
The table of
Proposition~\textup{\ref{prop:archetype-complementarity-bridge}} is
computed along three disjoint routes.
\textup{(V1)} \emph{Direct OPE computation.} For each archetype
witness, the scalar $\kappa(\cA)$ is computed from the OPE residue
$J(z)J(w) \sim \kappa/(z-w)^2 \otimes \Omega$
\textup{(}Heisenberg, KM via Sugawara\textup{)} or from the
conformal-weight parameter
\textup{(}$\beta\gamma_\lambda$: $\kappa = 6\lambda^2 - 6\lambda + 1$;
Virasoro: $\kappa = c/2$\textup{)}, then paired with the Verdier-dual
$\kappa^!$ evaluated at the dual parameter.
\textup{(V2)} \emph{Anomaly-ratio route.} Compute $\varrho = \kappa/c$
and $K = c + c^!$ separately, then multiply.
\textup{(V3)} \emph{Companion-pair midpoint.} Compute
$\kappa^{\mathrm{mid}}=K^\kappa/2$.  A regular fixed point of the
companion involution realizes this midpoint as an individual value.
The BP involution fixes the pole $k=-3$; its regular midpoint fibre
$c=25$ consists of $k=-3\pm2i$, exchanged by the involution.
The master computation
\texttt{test\_complementarity\_landscape.py} re-evaluates the four
unconditional witness formulas and the principal $\mathcal W_3$
formula.  The latter formula is read on
$H_{\mathrm{diag}}^{g=1}+H_{W_3}^{\mathrm{DS/bar}}$
\textup{(}see \texttt{compute/tests} for the stored
$(\varrho,K,K^\kappa,\kappa^*)$ tuples\textup{)}.  These evaluations
share the same companion formulas and constitute one computational
route.
The BP row carries the exact tuple component $K^c=50$ and open
genus-$1$ entries; the
$\mathsf B$ row carries the proved Mukai arithmetic
$(\operatorname{rk},\operatorname{sig},2c_+)=(24,(4,20),8)$ and the
conjectural chiral tuple under~$H_{\mathsf B}$.

\emph{Primary-source attribution for the $\mathcal{W}$-extension
values.} Feigin--Frenkel duality~\textup{\cite{Feigin-Frenkel}} and
Fateev--Lukyanov~\textup{\cite{FateevLukyanov88}} give the exact
principal central conductor $K_3^c=100$.  On
$H_{\mathrm{diag}}^{g=1}+H_{W_3}^{\mathrm{DS/bar}}$, the diagonal
multiplier $H_3-1=5/6$ carries this exact value to the conditional
modular conductor $250/3$.  The derivation uses the
Drinfeld--Sokolov central-charge formula at
$\fg = \mathfrak{sl}_3$; the number $25/3$
is the former BP proposal of
Remark~\textup{\ref{rem:bp-census-convention-separation}} and has
status \ClaimStatusConjectured.  Fehily--Kawasetsu--Ridout
\textup{\cite[\S2.1]{FKR20}} supply the bosonic/even parity and the
exact central-charge conductor~$50$; the direct genus-$1$ curvature
calculation remains the open modular input.  The same-family central
companion interpretation uses
Proposition~\textup{\ref{prop:bp-self-duality}} on
$H_{\mathrm{BP}}^{\mathrm{DS/bar}}$.
\end{remark}

\begin{remark}[Archetype and family data in the scalar sum]
\label{rem:archetype-vs-family-kappa-sum}
\index{quadrichotomy!and complementarity sum}
\index{archetype!does not determine $K^\kappa$}
The four shadow-archetype classes $\mathsf{G}/\mathsf{L}/\mathsf{C}/\mathsf{M}$
\textup{(}Theorem~\ref{thm:quadrichotomy}\textup{)} partition the
standard landscape by shadow-tower depth
$r_{\max} \in \{2, 3, 4, \infty\}$.  The scalar $K^\kappa$ is a
separate invariant.  The Virasoro witness contributes the
unconditional value~$13$.  On
$H_{\mathrm{diag}}^{g=1}+H_{W_N}^{\mathrm{DS/bar}}$, the principal
$\mathsf M$ families contribute the conditional values
$(H_N-1)K_N^c$, including $250/3$ at~$N=3$.  The BP family contributes
an open modular slot with exact central conductor~$50$. The
coarse archetype class controls the \emph{shadow tower} and the
\emph{$r$-matrix pole order} but not the scalar
complementarity constant. One letter $K$ carries two
referents -- the Trinity conductor $K = c + c^!$ and the scalar
sum $K^\kappa = \kappa + \kappa^!$ -- bridged by the anomaly ratio
$\varrho$. Only for Heisenberg, where $\varrho = 1$, do the two
referents coincide; for every other standard family they differ by
the factor $\varrho$, which is not archetype-invariant.
\end{remark}

\begin{remark}[Degree-three Hochschild and shadow witnesses across the
 archetype landscape]
\label{rem:chirhoch3-archetype-column}
\index{Hochschild cohomology!degree-three archetype table}
\index{Zamolodchikov W3 generator!as shadow/GF representative}
The archetype table of
Proposition~\textup{\ref{prop:archetype-complementarity-bridge}}
extends by one column, recording the degree-$3$ witness in its
native ambient. Theorem~\ref{thm:chirhoch3-five-archetype-portrait}
collapses through
Corollary~\ref{cor:chirhoch3-unified-formula} to
\[
 [\chi_3^{\cA}]\;=\;
 \begin{cases}
 0 & \cA\in\{\mathsf G,\mathsf L,\mathsf C,\mathsf M\}
      \text{ on the ambient curve}, \\
[W_3^{\mathrm{Zam}}]_{\mathrm{GF/shadow}} &
\cA=\mathrm{Vir}_c,\ c\text{ generic, GF/shadow or extension locus}, \\
 \mathbb C\cdot\bigl[{:}T\partial T{:}-\tfrac14\partial^3 T+
 \hbar\cdot\mathrm{qt}(J^{(3)})\bigr] &
 \cA=\mathbf H_{\Delta_5}=\Phi_3^{(\Sigma_2,C)}(\mathrm{K3}\times E)
      \text{ in the CY-$3$-enlarged product ambient},
 \end{cases}
\]
where the product ambient is $X\times(\mathrm{K3}\times E)$ for
the crown row. The weight-$5$ quasi-primary
$W_3^{\mathrm{Zam}}={:}T\partial T{:}-(1/4)\partial^3 T$ is the
Zamolodchikov $W_3$-extension generator
(Zamolodchikov~\textup{\cite{Zam1985}}~eq.~$(2.6)$). It is the
common GF/shadow representative in the $\mathsf M$ and $\mathbf B$
rows, not a non-zero ambient-curve $\ChirHoch^3$ class for
$\mathrm{Vir}_c$. The quantum tail
$\hbar\cdot\mathrm{qt}(J^{(3)})$ is the Schiffmann--Vasserot
degree-$3$ CoHA current; it produces the non-zero class only after
the CY-$3$ product ambient is included. The shadow-finite rows
$\mathsf G,\mathsf L,\mathsf C$ and the ambient-curve
$\mathsf M$ row admit only $\partial$-exact triples and carry
$[\chi_3^{\cA}]=0$ in curve-level chiral Hochschild cohomology.
\end{remark}

\begin{theorem}[Koszul self-dual locus across the witness landscape;
\ClaimStatusConditional]
\label{thm:census-self-dual-locus}
\index{Koszul self-duality!locus theorem|textbf}
\index{self-dual locus!witness families}
\index{Verdier duality!fixed points}
\textbf{Type signature.}
$(\mathrm{Open},\mathrm{bar/Verdier\ and\ principal\ DS},
\mathrm{level}=1\to5,
H_{\mathrm{diag}}^{g=1}+H_{W_3}^{\mathrm{DS/bar}}
\text{ for the principal modular midpoint})$.
For each witness family $\mathcal{F}$ in the list
$\{\cH_k,\,\widehat{\fg}_k,\,\beta\gamma_\lambda,\,\mathrm{Vir}_c,\,
\mathcal{W}_3^k,\,\mathrm{BP}_k\}$, let $\sigma_\mathcal{F}$ denote
the scalar companion involution displayed by
Theorem~\ref{thm:census-witness-complementarity}, and write
$\cA^\sigma$ for the family member at the companion parameter.
The corresponding comparison theorem identifies
$\cA^!\simeq\cA^\sigma$ on its named hypothesis package.  For BP this
package is $H_{\mathrm{BP}}^{\mathrm{DS/bar}}$. Write
$\mathrm{Fix}(\sigma_\mathcal{F})$ for the set of $\sigma_\mathcal{F}$-fixed
levels in the ambient parameter space of $\mathcal{F}$. Then the
Koszul self-dual locus $\{\cA \in \mathcal{F}\,:\,\cA \simeq \cA^!\}$
is tested, for the five families with an established genus-$1$
normalization, against the regular part of
$\mathrm{Fix}(\sigma_\mathcal{F})$ by the following scalar midpoint:
\begin{equation}\label{eq:self-dual-kappa-star}
 \kappa^*(\mathcal{F}) \;:=\;
 \tfrac{1}{2}\bigl(\kappa(\cA) + \kappa(\cA^\sigma)\bigr),
 \qquad
 \left.
 \begin{gathered}
 \cA^!\simeq\cA^\sigma,\\
 \cA\simeq\cA^!
 \end{gathered}
 \right\}
 \Longrightarrow \kappa(\cA) = \kappa^*(\mathcal{F}).
\end{equation}
For BP, the parameter involution and central-charge midpoint are exact,
while $\kappa^*(\mathrm{BP})$ has status \ClaimStatusOpen pending the
direct genus-$1$ non-separating curvature computation.
Concretely:
\begin{enumerate}
\item \emph{Heisenberg $\cH_k$ (class $\mathsf{G}$).}
 $\sigma_\cH(k) = -k$; $\mathrm{Fix}(\sigma_\cH) = \{0\}$, and
 $\kappa^*(\cH) = 0$. At the fixed level $k = 0$ the algebra is the
 uncurved free polynomial chiral algebra, \emph{but}
 $\cH_0 \not\simeq \cH_0^! = \mathrm{Sym}^{\mathrm{ch}}(V^*)$ as
 chiral algebras \textup{(}Theorem~\ref{thm:heisenberg-not-self-dual};
 Lie vs Com\textup{)}. Hence the self-dual locus is \emph{empty}, even
 though the fixed locus of $\sigma_\cH$ is non-empty. The scalar
 identity $\kappa = \kappa^* = 0$ is the numerical shadow of a
 categorical non-self-duality.

\item \emph{Affine Kac--Moody $\widehat{\fg}_k$ (class $\mathsf{L}$).}
 $\sigma_{\widehat{\fg}}(k) = -k - 2h^\vee$;
 $\mathrm{Fix}(\sigma_{\widehat{\fg}}) = \{-h^\vee\}$ (the critical
 level). $\kappa^*(\widehat{\fg}) = 0$ matches
 $\kappa(\widehat{\fg}_{-h^\vee}) = 0$ as scalars, but at $k = -h^\vee$
 the Sugawara construction degenerates, the Feigin--Frenkel center
 $\mathfrak{z}(\widehat{\fg})$ enlarges, and the Koszul dual is
 $(\widehat{\fg}_{-h^\vee})^! = \mathrm{CE}^{\mathrm{ch}}(\widehat{\fg}_{-h^\vee})$
 \textup{(}Theorem~\ref{thm:sl2-koszul-dual}, stated at $\fg =
 \mathfrak{sl}_2$; the proof is uniform in $\fg$\textup{)}. The
 critical algebra is \emph{not} isomorphic to its chiral
 Chevalley--Eilenberg complex; the self-dual locus is empty on the
 chiral-algebra side. The scalar fixed point $\kappa^* = 0$ is again
 a numerical-only self-duality.

\item \emph{$\beta\gamma_\lambda$ (class $\mathsf{C}$).}
 $\sigma_{\beta\gamma}(\lambda) = 1 - \lambda$ on the conformal-weight
 parameter, obtained from $(\beta\gamma_\lambda)^! \simeq bc_\lambda$
 composed with the parity exchange that takes the bosonic partner back
into the bosonic locus; equivalently, $\sigma_{\beta\gamma}$ is the
 unique involution of the $\lambda$-axis preserving
 $\kappa(\lambda) = 6\lambda^2 - 6\lambda + 1$.
 $\mathrm{Fix}(\sigma_{\beta\gamma}) = \{1/2\}$ on the
 conformal-weight axis, but this is not a fixed point of the actual
 Verdier-duality object
 $\beta\gamma_\lambda\mapsto bc_\lambda$; it is only the fixed point
 of the return map to the bosonic conformal-weight line. The
 complementarity constant of the actual pair is
 $K^\kappa_{\beta\gamma}=0$, hence
 $\kappa^*(\beta\gamma)=0$. Direct evaluation at $\lambda = 1/2$
 yields
 $\kappa(\beta\gamma_{1/2}) = -1/2$ and
 $\kappa((\beta\gamma_{1/2})^!) = \kappa(bc_{1/2}) = 1/2$, so the
 complementarity sum at the fixed point is $0$, consistent with
 Theorem~\ref{thm:census-witness-complementarity} \textup{(}which
 gives $\kappa + \kappa^! = 0$ at every $\lambda$\textup{)}. The
 partner $bc_{1/2}$ is the complex-fermion point
 $(c=1,\kappa=1/2)$, where
 $\beta\gamma_{1/2} \not\simeq bc_{1/2}$ as chiral algebras (bosonic
 versus fermionic parity), so the self-dual locus is empty. The
 parameter fixed-point scalar $\kappa = -1/2$ is not
 $\kappa^*(\beta\gamma)$; it is a numerical shadow of a categorical
 non-self-duality across the boson/fermion wall.

\item \emph{Virasoro $\mathrm{Vir}_c$ (class $\mathsf{M}$).}
 $\sigma_{\mathrm{Vir}}(c) = 26 - c$;
 $\mathrm{Fix}(\sigma_{\mathrm{Vir}}) = \{13\}$, and
 $\kappa^*(\mathrm{Vir}) = 13/2$. At $c = 13$ the Verdier-dual partner
 is $\mathrm{Vir}_{26-13} = \mathrm{Vir}_{13}$; the bar complex is
 same-family on the M/S-level locus
 \textup{(}Proposition~\ref{prop:virasoro-generic-koszul-dual}\textup{)},
 so the self-dual locus coincides with
 $\mathrm{Fix}(\sigma_{\mathrm{Vir}})$ on the scalar locus.

\item \emph{$\mathcal{W}_3^k$.} $\sigma_{\mathcal{W}_3}(k) = -k - 6$
 \textup{(}Feigin--Frenkel at $h^\vee = 3$\textup{)};
 $\mathrm{Fix}(\sigma_{\mathcal{W}_3}) = \{-3\}$.  The exact central
 conductor gives the formal midpoint $c_* = K_3^c/2 = 50$.  On
 $H_{\mathrm{diag}}^{g=1}+H_{W_3}^{\mathrm{DS/bar}}$, the conditional
 modular midpoint is $\kappa^*(\mathcal{W}_3) = 125/3$.
 The function $c(\mathcal{W}_3^{-3})$ approaches
 $\{-\infty,+\infty\}$ along the two
 approach directions.  The real-level image is
 $(-\infty,2]\cup[98,\infty)$, which is disjoint from the central
 midpoint~$50$.  Thus the midpoint fibre belongs to the complex
 parameter surface and its real-level intersection is~$\varnothing$.

\item \emph{Bershadsky--Polyakov $\mathrm{BP}_k$.}
 $\sigma_{\mathrm{BP}}(k) = -k - 6$;
 $\mathrm{Fix}(\sigma_{\mathrm{BP}}) = \{-3\}$, and this fixed level
 is the pole of the standard central-charge function.  The scalar
 central-charge companion midpoint is
 $c^{\mathrm{mid}}_{\mathrm{BP}}=50/2=25$.  The modular midpoint
 $\kappa^*(\mathrm{BP})$ is \ClaimStatusOpen.
 The BP central-charge image on real levels is
 $c(k) \leq 1$ for $k > -3$ and $c(k) \geq 49$ for $k < -3$, so
 the midpoint fibre $c(k)=25$ is the complex pair
 $k=-3\pm2i$, exchanged by~$\sigma_{\mathrm{BP}}$.
 Thus
 $\mathrm{Fix}(\sigma_{\mathrm{BP}})\cap U_{\mathrm{OPE}}=\varnothing$.
 On
 $H_{\mathrm{BP}}^{\mathrm{DS/bar}}$ the involution represents the
 same-family Verdier--Koszul comparison; its fixed point remains the
 singular level~$-3$.
\end{enumerate}
Consequently, among the regular fixed-parameter loci established for
the six witness families, Virasoro realizes the self-dual point
$c=13$.  The BP companion involution contributes the singular fixed
level~$k=-3$ and the exchanged midpoint fibre $k=-3\pm2i$. For the
$\mathcal{W}_N$-series at $N \geq 3$, the exact central conductor
gives the formal point $c^* = K_N^c/2 = 2N^3 - N - 1$ outside the
image of the level-to-central-charge map at every $N \geq 3$.
Equivalently, the real-level image and the formal fixed fibre are
disjoint.
\end{theorem}

\begin{proof}
On the comparison package for a family~$\mathcal F$ with established
genus-$1$ normalization, the equivalence
$\cA^!\simeq\cA^\sigma$ identifies
$K^\kappa_\mathcal{F}$ with
$\kappa(\cA)+\kappa(\cA^\sigma)$.
If $\cA\simeq\cA^!$, invariance of~$\kappa$ gives
$K^\kappa_\mathcal{F}=2\kappa(\cA)$ and therefore
$\kappa(\cA)=K^\kappa_\mathcal{F}/2=\kappa^*(\mathcal F)$.

\emph{Heisenberg.} $\sigma_\cH(k) = -k$ follows from
$\kappa(\cH_k^!) = -k$
\textup{(}Theorem~\ref{thm:census-witness-complementarity}\textup{)};
the fixed point is $k = 0$, where both $\cH_0$ and $\cH_0^!$ have
$\kappa = 0$. Theorem~\ref{thm:heisenberg-not-self-dual} shows that
the Koszul dual exchanges the Lie (non-commutative, centrally-extended
OPE) and Com (commutative, regular OPE) structures, hence
$\cH_0 \not\simeq \cH_0^!$ as chiral algebras. The vanishing
$\kappa^* = 0$ is a numerical identity and does not imply structural
self-duality.

\emph{Affine Kac--Moody.} \(\sigma_{\widehat{\fg}}(k)=-k-2h^\vee\)
is the critical-level reflection; the fixed point is
$k = -h^\vee$. Theorem~\ref{thm:sl2-koszul-dual} identifies the
Koszul dual with the chiral Chevalley--Eilenberg complex
$\mathrm{CE}^{\mathrm{ch}}(\widehat{\fg}_{-h^\vee})$, which is a
different chiral algebra (different underlying vector space:
$\widehat{\fg}$-valued chains versus $\widehat{\fg}^*$-valued
cochains). The proof in the $\mathfrak{sl}_2$ case is uniform in $\fg$
because the only input is the Sugawara-degeneration at the critical
level, which holds for every $\fg$ at $k = -h^\vee$ by the
Feigin--Frenkel center theorem \cite{Feigin-Frenkel}. Hence
$\widehat{\fg}_{-h^\vee} \not\simeq (\widehat{\fg}_{-h^\vee})^!$, and
the self-dual locus is empty.

\emph{$\beta\gamma_\lambda$.} The parameter-space involution
$\sigma_{\beta\gamma}$ on the conformal-weight $\lambda$ is determined
by the central-charge involution
$c(\lambda) = 2(6\lambda^2 - 6\lambda + 1)$ whose only non-trivial
automorphism is $\lambda \mapsto 1 - \lambda$ (fixing
$\kappa = c/2 = 6\lambda^2 - 6\lambda + 1$). The fixed point is
$\lambda = 1/2$, where $\kappa(\beta\gamma_{1/2}) = -1/2$. Its
Koszul dual is $bc_{1/2}$, the fermionic complex-fermion partner;
these have opposite statistics (bosonic versus fermionic) and differ
structurally (commutative-at-generator OPE residue for $\beta\gamma$
versus anti-commutative for $bc$; $\bar{B}^{\mathrm{ch}}$ is a
symmetric-coalgebra in one case and an anti-symmetric-coalgebra in
the other). Hence $\beta\gamma_{1/2} \not\simeq bc_{1/2}$, and the
self-dual locus is empty.

\emph{Virasoro.} Proposition~\ref{prop:virasoro-generic-koszul-dual}
identifies $\mathrm{Vir}_c^! \simeq \mathrm{Vir}_{26-c}$ on the
M/S-level locus. The fixed point of $\sigma_{\mathrm{Vir}}$ is
$c = 13$; the dual chiral algebra is $\mathrm{Vir}_{13}$, the same
family at the same central charge, so the self-dual isomorphism
holds on the M/S-level locus. The scalar
$\kappa^* = 13/2$ matches
$\kappa(\mathrm{Vir}_{13}) = 13/2$ by
$\kappa(\mathrm{Vir}_c) = c/2$
\textup{(}\texttt{landscape\_census.tex}\textup{)}. The locus is the
singleton $\{c = 13\}$ on the scalar locus.  The normalized
C1-to-Theorem~D trace gives
$\kappa+\kappa^!=13$ at this Virasoro self-dual point.

\emph{$\mathcal{W}_3^k$.} The central-charge formula
$c_{\mathcal{W}_3}(k) = 2 - 24(k+2)^2/(k+3)$ has image
$(-\infty, 2]$ for $k > -3$ and $[98, +\infty)$ for $k < -3$, with
a singularity at $k = -3$; the fixed point $c_* = 50$ lies in the
gap $(2, 98)$ \textup{(}exclusive on both endpoints, cf.\
Remark~\ref{rem:bp-conductor-comparison}\textup{)}.  Hence the
real-level image and the fixed central fibre are disjoint, and the
same-family fixed locus lies on the complex parameter surface.  The package
$H_{W_3}^{\mathrm{DS/bar}}$ promotes the scalar reflection to the
principal bar--Verdier companion.  The formal fixed level is
the complex-conjugate pair determined by
$(k+2)^2/(k+3) = -2$ on the $c = 50$ fibre, i.e.\
$k = -3 \pm i$.

\emph{Bershadsky--Polyakov.}
Theorem~\ref{thm:bp-koszul-conductor-polynomial} gives
$c(k)+c(-k-6)=50$ in~$\mathbb Q(k)$,
so the exact central-charge midpoint is~$25$.  The BP modular
quantities $\kappa_{\mathrm{BP}}$, $K^\kappa_{\mathrm{BP}}$, and
$\kappa^*(\mathrm{BP})$ have status \ClaimStatusOpen pending a direct
genus-$1$ curvature computation.  The BP central-charge function
$c_{\mathrm{BP}}(k) = -(2k+3)(3k+1)/(k+3)$
\textup{(}Kac--Roan--Wakimoto minimal reduction,
Proposition~\ref{prop:bp-central-charge}\textup{)}
has image
$c(k) \leq 1$ for $k > -3$ and $c(k) \geq 49$ for $k < -3$
\textup{(}Remark~\ref{rem:bp-self-duality-central-charge}\textup{)},
so the midpoint $c^{\mathrm{mid}}=25$ lies outside the real image.
Solving $c(k)=25$ gives $k=-3\pm2i$; these two regular levels are
exchanged by $k\mapsto-k-6$.  The unique fixed level is $k=-3$, the
pole.  Proposition~\ref{prop:bp-self-duality} promotes the companion
pair to the same-family Verdier--Koszul pair precisely on
$H_{\mathrm{BP}}^{\mathrm{DS/bar}}$.

\emph{$\mathcal{W}_N$ for $N \geq 3$.} The exact central conductor is
$K_N^c = 4N^3 - 2N - 2$
\textup{(}Remark~\ref{rem:koszul-conductor-explicit}\textup{)}, so
the formal self-dual central charge is $c^*_N = K_N^c/2 = 2N^3 - N - 1$:
$c^*_2 = 13$, $c^*_3 = 50$, $c^*_4 = 123$, $c^*_5 = 244$. At $N \geq 3$,
the principal central-charge image on real levels is
$c(k) \in (-\infty, N-1] \cup [c^+_N, \infty)$ where the upper lobe
$c^+_N$ is determined by the Kac--Roan--Wakimoto quadratic in
$(k + h^\vee - s_\lambda)^2$; in each case $c^+_N > c^*_N$ so the
fixed value $c^*_N$ lies in the excluded gap. Specifically at
$N = 3$, $c^+_3 = 98 > 50 = c^*_3$, confirmed by the explicit image
analysis above; analogous computations at $N = 4, 5$ give
$c^+_4 = 208 > 123$ and $c^+_5 = 400 > 244$, respectively.  Hence the
real-level image is disjoint from the same-family fixed fibre for
$N\geq3$.  The package $H_{W_N}^{\mathrm{DS/bar}}$
promotes the reflected scalar pair to the principal bar--Verdier
companion.
Only $N = 2$ (Virasoro) has its self-dual central charge
$c^*_2 = 13$ attained on the real locus.
\end{proof}

\begin{remark}[Three-path computation of
Theorem~\ref{thm:census-self-dual-locus}]
\label{rem:self-dual-locus-verification}
\index{Koszul self-duality!three-path computation}
The identification $\kappa^*(\mathcal{F}) = K^\kappa_\mathcal{F}/2$
is computed three independent ways on each family.
\textup{(V1)} \emph{Direct fixed-point computation of
$\sigma_\mathcal{F}$.} For Heisenberg:
$\sigma(k) = -k$, fixed point $k = 0$, $\kappa(0) = 0$. For affine KM:
$\sigma(k) = -k - 2h^\vee$, fixed point $k = -h^\vee$,
$\kappa(-h^\vee) = 0$. For Virasoro: $\sigma(c) = 26 - c$, fixed
point $c = 13$, $\kappa(13) = 13/2$. For
$\beta\gamma$: $\sigma(\lambda) = 1 - \lambda$, fixed
$\lambda = 1/2$, $\kappa(1/2) = -1/2$. For $\mathcal{W}_3$:
$\sigma(k) = -k - 6$, fixed $k = -3$, formal
$c$-midpoint $50$. For BP the same involution fixes the pole $k=-3$,
while its standard scalar midpoint is~$25$.
\textup{(V2)} \emph{Complementarity-sum halving.}
From $K^\kappa_\mathcal{F}$ in
Theorem~\ref{thm:census-witness-complementarity}:
$\{0,0,0,13\}/2=\{0,0,0,13/2\}$ for the four unconditional
Vol~I witness rows.  On
$H_{\mathrm{diag}}^{g=1}+H_{W_3}^{\mathrm{DS/bar}}$, the conditional
principal extension adds $(250/3)/2=125/3$.
The conditional $\mathsf B$ candidate gives $8/2=4$ under
$H_{\mathsf B}$.  The BP
modular midpoint is open; the former proposal is typed
\ClaimStatusConjectured in
Remark~\ref{rem:bp-census-convention-separation}. \textup{(V3)}
\emph{Standard scalar-conductor cross-check.}  The exact principal
$\mathcal W_3$ central conductor is~$100$, with midpoint~$50$; the standard
BP conductor is~$50$, with midpoint~$25$.  Their midpoint fibres are
$k=-3\pm i$ and $k=-3\pm2i$, respectively, and each pair is exchanged
by $k\mapsto-k-6$.  The involution itself fixes the common pole
$k=-3$.  The shifted BP scalar convention and its conductor~$196$
remain the separately typed lane of
Remark~\ref{rem:bp-census-convention-separation}.

The theorem's regular fixed-parameter conclusion---the realization of
$\{\mathrm{Vir}_{13}\}$ as the real-level same-family self-dual point
on the M/S-level locus of
Proposition~\ref{prop:virasoro-generic-koszul-dual}---is computed by
three paths:
\textup{(S1)} via the same-family partner
$\mathrm{Vir}_c^! \simeq \mathrm{Vir}_{26-c}$ at $c = 13$;
\textup{(S2)} via the central-charge involution $c \mapsto 26 - c$ and
its unique real fixed point $c = 13$;
\textup{(S3)} via the universal-conductor identity
$K_{\mathrm{Vir}} = 26$ of
Remark~\ref{rem:koszul-conductor-explicit} specialised to $N = 2$,
giving $c^*_2 = 13$ -- consistent with the universal-conductor
self-dual point predicted by the BRST construction of
Chapter~\ref{chap:universal-conductor}.
\end{remark}

\begin{remark}[Koszul conductor for \texorpdfstring{$\mathcal{W}_N$}{WN}]\label{rem:koszul-conductor-explicit}
\index{Koszul conductor!explicit formula}
\textbf{Type signature.}
$(\mathrm{Open},\mathrm{principal\ DS\ reflection},\mathrm{level}=1;
H_{W_N}^{\mathrm{DS/bar}}\text{ for the bar--Verdier companion})$.
The exact reflected central conductor admits a closed form for the
$\mathcal{W}_N$ series.  The Fateev--Lukyanov~\cite{FateevLukyanov88}
Drinfeld--Sokolov central charge on the principal orbit is
\[
c(\mathcal{W}_N^k) \;=\; (N-1) - \frac{N(N^2-1)(k+N-1)^2}{k+N}\,.
\]
The critical-level reflected companion level is
\(k'=-k-2N\).  It is the scalar reflection datum in the
Feigin--Frenkel/Drinfeld--Sokolov comparison lane.  On
$H_{W_N}^{\mathrm{DS/bar}}$ it represents the completed bar--Verdier
companion.  The exact rational-function sum evaluates (via the
identity \((k+N+1)^2-(k+N-1)^2=4(k+N)\)):
\[
K_N^c \;=\; c(\mathcal{W}_N^k) + c(\mathcal{W}_N^{k'})
\;=\; 2(N-1) + 4N(N^2-1)
\;=\; 2(N-1)(2N^2 + 2N + 1).
\]
The Freudenthal--de Vries identity gives the same total via
$c+c'=2r+4h^\vee d$ with $r=N-1$, $h^\vee=N$, and $d=N^2-1$.
Thus
\[
K_N^c=4N^3-2N-2,\qquad
(K_2^c,K_3^c,K_4^c,K_5^c)=(26,100,246,488).
\]
On $H_{\mathrm{diag}}^{g=1}+H_{W_N}^{\mathrm{DS/bar}}$, the modular
conductor is
\[
K_N^{\kappa_{\mathrm{mod}}}=(H_N-1)K_N^c.
\]
At $N=3$, the exact central conductor is $K_3^c=100$ and the
conditional modular conductor is
$K_3^{\kappa_{\mathrm{mod}}}=250/3$ on
$H_{\mathrm{diag}}^{g=1}+H_{W_3}^{\mathrm{DS/bar}}$.
\end{remark}

\begin{proposition}[Fateev--Lukyanov central-charge formula; canonical form; \ClaimStatusProvedHere]
\label{prop:fateev-lukyanov-canonical}
The principal $\cW_N$ central charge
\begin{equation}\label{eq:fateev-lukyanov-wn}
c(\mathcal{W}_N^k) \;=\; (N-1) - \frac{N(N^2-1)(k+N-1)^2}{k+N}
\end{equation}
is the canonical formula used in all three volumes. The alternative
parameterisation
$c(t)=(N-1)(1-N(N+1)(t-1)^2/t)$ with $t=k+N$
(Fateev--Lukyanov~\cite{FateevLukyanov88}, equation~(7)) is equivalent
by $(N-1)N(N+1)=N(N^2-1)$. Any displayed formula using
$(k+N)^{-1}(k'+N)^{-1}$ in place of $(k+N-1)^2/(k+N)$ is a
mixed-parameter identity valid only at the Koszul self-dual locus.
\end{proposition}

\begin{proof}
The displayed Fateev--Lukyanov form is obtained from
$c(t)=(N-1)(1-N(N+1)(t-1)^2/t)$ by substituting $t=k+N$ and using
$(N-1)N(N+1)=N(N^2-1)$. The critical-level reflection sends
$k+N$ to $-(k+N)$, so the conductor sum reduces to
$2(N-1)+4N(N^2-1)$. A formula with simultaneous factors
$(k+N)^{-1}(k'+N)^{-1}$ uses both parameters after imposing the
involution and therefore agrees with the canonical formula only at
the self-dual value.
\end{proof}

\begin{remark}[Numerical checks for the Koszul conductor]
Numerical checks:
\begin{itemize}
\item $N=2$: $c(\mathrm{Vir}_{k}) = 1 - 6(k+1)^2/(k+2)$; at $k=-3/2$ (central-charge zero locus) gives $c = 1 - 6 \cdot 1/4 \cdot 2 = 1 - 3 = -2$; at $k = 1$ gives $c = 1 - 24/3 = -7$ (minimal model $\mathcal{M}(2,5)$ via DS).
\item $N=3$: $c(\mathcal{W}_3^k) = 2 - 24(k+2)^2/(k+3)$; at $k = -3/2$ gives $c = 2 - 24 \cdot 1/4 \cdot 2/3 = 2 - 4 = -2$; at $k = -9/4$ gives $c = 2 - 24 \cdot 1/16 \cdot 4/3 = 2 - 2 = 0$ (degenerate $\mathcal{W}_3(4,3)$).
\item Complementarity at $N=2$: the partner $k' = -k - 4$ sends $k=1 \mapsto k' = -5$, giving $c' = 1 - 6 \cdot 16/(-3) = 1 + 32 = 33$, and $c + c' = -7 + 33 = 26 = K_2$.
\end{itemize}
\end{remark}

\begin{remark}[Exceptional and non-simply-laced complementarity]\label{rem:e8-complementarity}
\index{E8@$E_8$!complementarity}
\index{E6@$E_6$!complementarity}
\index{E7@$E_7$!complementarity}
\index{non-simply-laced!complementarity}

The following table records the complete data for exceptional and
non-simply-laced affine Kac--Moody algebras on the non-critical
universal locus, all proved instances of the Master Table. Every entry
satisfies $\kappa + \kappa' = 0$ and $c + c' = 2d$ (the clean
anti-symmetric case for Kac--Moody), and every non-critical family is
class L with shadow depth $r_{\max} = 3$.
\begin{center}
\renewcommand{\arraystretch}{1.3}
\begin{tabular}{lcccccl}
\toprule
$\mathfrak{g}$ & $d$ & $h$ & $h^\vee$ & $\kappa(\widehat{\mathfrak{g}}_k)$ & $c+c'$ & Koszul conductor $K$ \\
\midrule
\multicolumn{7}{c}{\textit{Classical simply-laced ($h = h^\vee$)}} \\
$D_4$ (triality) & $28$ & $6$ & $6$ & $7(k+6)/3$ & $56$ & $56 = 2\cdot 28$ \\
\midrule
\multicolumn{7}{c}{\textit{Exceptional (simply-laced, $h = h^\vee$)}} \\
$E_6$ & $78$ & $12$ & $12$ & $13(k+12)/4$ & $156$ & $156 = 2\cdot 78$ \\
$E_7$ & $133$ & $18$ & $18$ & $133(k+18)/36$ & $266$ & $266 = 2\cdot 133$ \\
$E_8$ & $248$ & $30$ & $30$ & $62(k+30)/15$ & $496$ & $496 = 2\cdot 248$ \\
\midrule
\multicolumn{7}{c}{\textit{Non-simply-laced ($h \neq h^\vee$; level shift uses $h^\vee$, not $h$)}} \\
$B_2 \cong \mathfrak{so}_5$ & $10$ & $4$ & $3$ & $5(k+3)/3$ & $20$ & $20 = 2\cdot 10$ \\
$C_2 \cong \mathfrak{sp}_4$ & $10$ & $4$ & $3$ & $5(k+3)/3$ & $20$ & $20 = 2\cdot 10$ \\
$G_2$ & $14$ & $6$ & $4$ & $7(k+4)/4$ & $28$ & $28 = 2\cdot 14$ \\
$F_4$ & $52$ & $12$ & $9$ & $26(k+9)/9$ & $104$ & $104 = 2\cdot 52$ \\
\midrule
\multicolumn{7}{c}{\textit{General BCD series ($N \geq 3$)}} \\
$B_N$ & $N(2N{+}1)$ & $2N$ & $2N{-}1$ & $N(2N{+}1)(k{+}2N{-}1)/(2(2N{-}1))$ & $2N(2N{+}1)$ & \\
$C_N$ & $N(2N{+}1)$ & $2N$ & $N{+}1$ & $N(2N{+}1)(k{+}N{+}1)/(2(N{+}1))$ & $2N(2N{+}1)$ & \\
$D_N$ & $N(2N{-}1)$ & $2(N{-}1)$ & $2(N{-}1)$ & $N(2N{-}1)(k{+}2N{-}2)/(4(N{-}1))$ & $2N(2N{-}1)$ & \\
\bottomrule
\end{tabular}
\end{center}

For $E_8$ the conductor is $K = 2\dim E_8 = 496$.  Its numerical coincidence with $\dim(E_8\times E_8) = 496$ is the identity $2\cdot248 = 496$, an instance of $K = 2\dim\fg$ valid on every row of the table; no derivation connects it to the heterotic gauge group.

The anomaly ratio $\varrho(\mathfrak{g}) = \sum_{i=1}^r 1/(m_i+1)$ for the associated principal $\mathcal{W}$-algebra (see Remark~\ref{rem:general-w-kappa-values}) satisfies $\varrho(E_8) = 121/126 < 1$: the $E_8$ shadow obstruction tower converges at all non-critical levels. More precisely:
\[
\varrho(E_8) = \frac{1}{2} + \frac{1}{8} + \frac{1}{12} + \frac{1}{14} + \frac{1}{18} + \frac{1}{20} + \frac{1}{24} + \frac{1}{30} = \frac{121}{126} \approx 0.960.
\]
The ordering $\varrho(E_6) > \varrho(E_7) > \varrho(E_8)$ (see Remark~\ref{rem:general-w-kappa-values}) reflects the increasing sparsity of the exponents: $E_8$ has large exponents $(29, 23, \ldots)$ giving small $1/(m_i + 1)$ contributions.

The type $D_4$ entry exhibits $S_3$-triality: the three
$8$-dimensional representations (vector, spinor, co-spinor) are
permuted by outer automorphisms, and all three give the same
$\kappa = 7(k+6)/3$. This is the unique simple Lie algebra with a
non-cyclic outer automorphism group.

For non-simply-laced types, the dual Coxeter number $h^\vee$ governs the kappa formula and critical-level reflection \(k \mapsto -k - 2h^\vee\). The Coxeter number $h$ appears in the Strange Formula $\lvert\rho\rvert^2 = dh/12$ and in the exponent sum $\sum m_i = rh/2$, but \emph{not} in this critical-level reflection or the curvature formula. The low-rank isomorphism $B_2 \cong C_2$ ($\mathfrak{so}_5 \cong \mathfrak{sp}_4$) produces identical entries: both have $d = 10$, $h^\vee = 3$, and the same kappa formula. For $N \geq 3$, types $B_N$ and $C_N$ are genuinely distinct: they share $d = N(2N+1)$ and $h = 2N$, but $h^\vee(B_N) = 2N-1$ while $h^\vee(C_N) = N+1$, so their kappa formulas and level-shifting dualities diverge.

\emph{Shadow depth universality away from the critical locus.} Every
non-critical affine Kac--Moody algebra $\widehat{\fg}_k$,
for \emph{every} simple Lie type ($A_N$, $B_N$, $C_N$, $D_N$, $G_2$,
$F_4$, $E_6$, $E_7$, $E_8$), has shadow depth class~L with
$r_{\max} = 3$. The cubic shadow $C$ is the Lie bracket, and the tower
terminates because the quartic shadow vanishes by the Jacobi identity.
At $k=-h^\vee$ the scalar curvature vanishes and the Feigin--Frenkel
center moves the bar object to the critical FF boundary; this is not a
class-L Koszul row. The exceptional and non-simply-laced non-critical
cases are checked by
\texttt{compute/lib/bar\_cohomology\_non\_simply\_laced\_engine.py}
(120+ tests).

Note: for Kac--Moody algebras the anomaly ratio $\varrho(\mathfrak{g})$ is defined for the associated $\mathcal{W}$-algebra $\mathcal{W}(\mathfrak{g}, f_{\mathrm{prin}})$, not for $\widehat{\mathfrak{g}}_k$ itself; the Kac--Moody obstruction coefficient is $\kappa = (k+h^\vee)\dim\mathfrak{g}/(2h^\vee)$ directly.
\end{remark}

\section*{Key results by category}

\emph{$\Einf$-Chiral Algebras (Vertex Algebras).}

\begin{enumerate}[label=(\roman*)]
\item \emph{Heisenberg}: Koszul dual is the commutative chiral algebra
$\mathrm{Sym}^{\mathrm{ch}}(V^*)$ with curvature $m_0 = -\kappa \cdot \mathbf{1}$;
the bar complex is the Chevalley--Eilenberg coalgebra
(Theorem~\ref{thm:heisenberg-koszul-dual-early}).
Complete genus-$1$ computation for $\widehat{\mathfrak{sl}}_2$ computed
in \S\ref{sec:sl2-genus-one-computation}.
\item \emph{$bc$--$\beta\gamma$}: Clifford/Sym exchange;
bar construction $\bar{B}(bc) = \Lambda^c(V^*)$, Koszul dual
$(bc)^! \cong \beta\gamma$
(\S\ref{sec:betagamma-koszul-dual}).
\item \emph{Affine Kac--Moody at non-critical level}: Koszul dual is
$\mathrm{CE}^{\mathrm{ch}}(\widehat{\fg}_{-k-2h^\vee})$,
the chiral CE algebra at the Feigin--Frenkel dual level
(same~$\kappa$ as $\widehat{\fg}_{-k-2h^\vee}$, but a
distinct object as a chiral algebra); explicit computation through degree~$3$
for $\mathfrak{sl}_2$ (Theorem~\ref{thm:sl2-koszul-dual}),
general case via Verdier duality on FM spaces
(Theorem~\ref{thm:universal-kac-moody-koszul}).
\item \emph{Virasoro and finite-type $\mathcal{W}$-algebras}:
for principal finite-type $\mathcal{W}_N$-algebras, the package
$H_{\mathrm{diag}}^{g=1}+H_{W_N}^{\mathrm{DS/bar}}$ gives
$\kappa(\mathcal{W}^k(\fg)^!) = \kappa(\mathcal{W}^{-k-2h^\vee}(\fg))$
\textup{(}Theorem~\ref{thm:w-algebra-koszul-main}\textup{)}.
For the Virasoro row, the same involution supplies the proved
same-family shadow partner $\mathrm{Vir}_{26-c}$ used by the
genus and semi-infinite calculations; the stronger H-level
infinite-generator realization is delegated to the exact
$W_\infty$ MC4 data of filtered H-level targets, residue identities,
and finite detection on $\mathcal{I}_N$, whose first live
higher-spin stages are already reduced on the theorem surface to the
stage-$3$ fifteen-coefficient packet and the stage-$4$
exact six-entry identity packet on~$\mathcal{I}_4$, organized as three
local OPE blocks and splitting into four higher-spin channels together
with two Virasoro-target channels whose principal values are fixed and
whose residue-side contraction requires the normalized two-point/Ward
identity of Corollary~\ref{cor:winfty-stage4-residue-four-channel}. The
first next reduced stage is the explicit eleven-entry packet
$\mathcal{J}_5^{\mathrm{red}}$ of
Corollary~\ref{cor:winfty-ds-stage5-reduced-packet}.
Complete genus-$1$ computations for $\mathrm{Vir}_c$
(\S\ref{sec:virasoro-genus-one-computation}, $c + c' = 26$, $\kappa = c/2$)
and the exact principal central identity
$c(\mathcal W_3^k)+c(\mathcal W_3^{-k-6})=100$.  On
$H_{\mathrm{diag}}^{g=1}+H_{W_3}^{\mathrm{DS/bar}}$, the principal
genus-$1$ modular value is $\kappa_{\mathrm{mod}}=5c/6$
\textup{(}\S\ref{sec:w3-genus-one-computation}\textup{)}.
For general principal $\mathcal{W}_N$, the formula
$\kappa_{\mathrm{mod}}=c(H_N-1)$ is conditional on
$H_{\mathrm{diag}}^{g=1}+H_{W_N}^{\mathrm{DS/bar}}$
\textup{(}Theorem~\ref{thm:wn-obstruction}\textup{)}.
\end{enumerate}

\emph{$\Eone$-Chiral Algebras (Nonlocal Vertex Algebras).}
\begin{enumerate}[label=(\roman*)]
\item Lattice algebras with non-symmetric cocycles: first strictly $\Eone$ examples;
Koszul dual uses inverse cocycle.
\item R-twisted vertex algebras: Yang--Baxter equation ensures associativity;
Yangians are fundamental examples.
\item Yangians: the proved content now includes the natural DK core
 through DK-2/3. At chain level
 $Y(\fg)^! \cong Y_{R^{-1}}(\fg)$ with braiding reversal
 \textup{(}Theorem~\ref{thm:yangian-koszul-dual}\textup{)}, and on the
 evaluation locus the ordered factorization categories are equivalent
 \textup{(}Theorem~\ref{thm:factorization-dk-eval}\textup{)}; the
 evaluation-generated factorization core is also proved at all simple
 types \textup{(}Corollary~\ref{cor:dk23-all-types}\textup{)}. The
 extension from that core to the
 ordinary-derived/completed/coderived enlargement, and then the
 dg-shifted Yangian comparison built on the standard RTT tower,
 remain conjectural
 \textup{(}Conjectures~\ref{conj:full-derived-dk}
 and~\ref{conj:full-dk-bridge}\textup{)}.
\item Toroidal and elliptic algebras: double affine structures with theta function OPEs.
\end{enumerate}

Deformation quantization.
\begin{enumerate}[label=(\roman*)]
\item Coisson algebras (PVA) quantize to $\Einf$-chiral algebras (vertex algebras)
via configuration space integrals (Theorem~\ref{thm:chiral-kontsevich}).
\item $\Pinf$-chiral algebras ($\Einf + L_\infty$) quantize to $\Eone$-chiral
algebras (nonlocal VAs; Construction~\ref{constr:quantization}).
\item Maurer--Cartan elements classify deformations at each level.
\end{enumerate}

\section*{The master table and the main theorems}

Every completed PH row of Table~\ref{tab:master-invariants}
instantiates the four main theorems on its stated locus:
Theorem~A (Theorem~\ref{thm:bar-cobar-isomorphism-main}) produces the
Koszul dual / partner column for $E_\infty$-chiral algebras (for the
Yangian, an $\Eone$-chiral algebra, the analogous role is played by
the $\Eone$-chiral Koszul duality theorem,
Theorem~\ref{thm:e1-chiral-koszul-duality});
Theorem~C (Theorem~\ref{thm:quantum-complementarity-main}) produces the
strict-flat centre-local-system eigensummands.  Their normalized
Theorem~D trace produces the $K^\kappa$ entries; the $c+c'$ column is
computed separately from the displayed family involutions and
central-charge formulae;
Theorem~\ref{thm:genus-universality} produces the $\kappa$ column.
The individual synthesis remarks (Remark~\ref{rem:sl2-three-theorems}
(Kac--Moody), Remark~\ref{rem:free-field-three-theorems} (free fields),
Remark~\ref{rem:lattice-three-theorems} (lattices),
Remark~\ref{rem:w-algebra-three-theorems} (W-algebras),
Remark~\ref{rem:deformation-three-theorems} (deformation quantization),
Remark~\ref{rem:yangian-three-theorems} (Yangians)) provide
detailed computations for each example class.
CJ and HF rows are boundary records: they display the known scalar or
structural input and leave the missing partner, conductor, or
bar-rank computation as an explicit obligation.

%% ================================================================
%% Shadow connection and propagator variance census
%% ================================================================

\section*{Shadow connection and propagator variance census}
\label{sec:shadow-connection-census}
\index{shadow connection!census}
\index{propagator variance!census}

The shadow connection
$\nabla^{\mathrm{sh}} = d - Q_L'/(2Q_L)\,dt$
(Theorem~\ref{thm:shadow-connection})
is a logarithmic connection with residue~$1/2$ at every
simple zero of the shadow metric $Q_L(t)$ and monodromy~$-1$
(the Koszul sign). The flat section is
$\Phi(t) = \sqrt{Q_L(t)/Q_L(0)}$, which is \emph{not}
the shadow generating function $H(t) = t^2\sqrt{Q_L(t)}$
($H$ has a double zero at $t=0$).

\begin{table}[ht]
\centering
\caption{Shadow connection data for Virasoro and $\cW_3$}
\label{tab:shadow-connection-census}
\renewcommand{\arraystretch}{1.3}
{\small
\begin{tabular}{|l|c|c|c|c|}
\hline
\textbf{Algebra / line}
 & $\boldsymbol{Q_L(t)}$
 & $\boldsymbol{\Delta}$
 & \textbf{Residues}
 & \textbf{Monodromy} \\
\hline
$\mathrm{Vir}_c$ (T-line)
 & $c^2 + 12ct + \frac{180c{+}872}{5c{+}22}\,t^2$
 & $\frac{40}{5c{+}22}$
 & $\frac{1}{2}$
 & $-1$ \\
\hline
$\cW_3$ (W-line)
 & $\frac{4c^2}{9} + \frac{40960}{3(5c{+}22)^3}\,t^2$
 & $\frac{20480}{3(5c{+}22)^3}$
 & $\frac{1}{2}$
 & $-1$ \\
\hline
$\cW_3$ (T-line)
 & \multicolumn{4}{c|}{Equals Virasoro data (autonomous)} \\
\hline
\end{tabular}
}% end small
\end{table}

\noindent
The W-line of $\cW_3$ has $\alpha_W = 0$ ($\mathbb{Z}_2$ parity),
so $Q_W$ is even in~$t$ with purely imaginary branch points.
The Gaussian decomposition is $Q_W = (2c/3)^2 + 2\Delta_W\,t^2$
with constant Gaussian envelope.
Checked by \texttt{test\_shadow\_connection.py}
(114 tests: residues, monodromy, flat sections, Picard--Fuchs,
Koszul duality, complementarity).

\begin{table}[ht]
\centering
\caption{Propagator variance $\delta_{\mathrm{mix}}$ for
 multi-generator families}
\label{tab:propagator-variance-census}
\index{propagator variance!census table}
\renewcommand{\arraystretch}{1.3}
{\small
\begin{tabular}{|l|c|c|c|l|}
\hline
\textbf{System}
 & \textbf{Rank}
 & $\boldsymbol{\delta_{\mathrm{mix}}}$
 & \textbf{Auto.}
 & \textbf{Enhanced symmetry} \\
\hline
$\cW_3$ ($T, W$)
 & $2$
 & $\frac{1280\,P^2}{c^3(5c{+}22)^6}$
 & No
 & $c = -2 \pm \tfrac{4\sqrt{33}}{5}$ \\[4pt]
\hline
$\cW_4$ ($T, W_3, W_4$)
 & $3$
 & $> 0$ for $c > 0$
 & No
 & None (complex only) \\
\hline
$\cW_5$ ($T, W_3, W_4, W_5$)
 & $4$
 & $> 0$ for $c > 0$
 & No
 & $c = -38/5$ \\
\hline
$\beta\gamma + bc$
 & $2$
 & $0$
 & Yes
 & Everywhere \\
\hline
$\cH + \mathrm{Vir}$
 & $2$
 & $\frac{1600}{c(c{+}2)(5c{+}22)^2}$
 & No
 & None ($P = \mathrm{const}$) \\[4pt]
\hline
$V(\mathfrak{sl}_2) + V(\mathfrak{sl}_3)$
 & $2$
 & $0$
 & Yes
 & Everywhere \\
\hline
\end{tabular}
}% end small
\end{table}

\noindent
Here $P = P(\cW_3) = 25c^2 + 100c - 428$ is the mixing polynomial.
Autonomy means $\delta_{\mathrm{mix}} = 0$: the multi-channel
shadow ODE decouples into independent single-channel equations
(Proposition~\ref{prop:propagator-variance}).
Enhanced symmetry points are the real zeros of~$P(c)$ where
the quartic gradients $f_i/\kappa_i$ become proportional.
Structural dichotomy: class~L + class~L is always autonomous
(both $f_i = 0$); class~G + class~M is never autonomous at
real $c > 0$ ($f_{\mathrm{G}} = 0$ but $f_{\mathrm{M}} \neq 0$,
so the ratios cannot coincide).
$\cW_4$ is \emph{never} autonomous at any real central charge.
Checked by
\texttt{test\_propagator\_variance\_landscape.py}
(143 tests: Cauchy--Schwarz, cross-engine consistency,
enhanced symmetry zeros, mixing polynomial structure).

%% ================================================================
%% Shadow fields and Epstein L-functions
%% ================================================================

\section*{Shadow fields and Epstein $L$-functions}
\label{sec:shadow-l-functions}
\index{shadow field!census|textbf}
\index{Epstein zeta function!standard families}
\index{Dirichlet L-function@Dirichlet $L$-function!from shadow metric}

The shadow metric $Q_L$ of a class~$\mathbf{M}$ algebra at
rational central charge~$c$ has discriminant
$\operatorname{disc}(Q_L) = -32\kappa^2\Delta$
(Proposition~\ref{prop:hankel-extraction}(iii)), which determines
an imaginary quadratic field $K_L = \bQ(\sqrt{D_0})$ and,
via the Dedekind factorization, a Dirichlet $L$-function
$L(s, \chi_d)$
(Theorem~\ref{thm:shadow-epstein-zeta},
Remark~\ref{rem:shadow-field}).
For class~$\mathbf{G}$ and~$\mathbf{L}$ algebras,
$\Delta = 0$ and the Epstein zeta degenerates.
Table~\ref{tab:shadow-l-functions} records the shadow field data
for all standard families at representative central charges.

\begin{table}[ht]
\centering
\caption{Shadow field and Epstein $L$-function data}
\label{tab:shadow-l-functions}
\index{shadow field!standard families}
\renewcommand{\arraystretch}{1.4}
{\small
\begin{tabular}{|l|c|c|c|c|c|}
\hline
\textbf{Family}
 & \textbf{Class}
 & $\boldsymbol{\operatorname{disc}(Q_L)}$
 & $\boldsymbol{K_L}$
 & $\boldsymbol{d}$
 & \textbf{$L$-function} \\
\hline
\multicolumn{6}{|c|}{%
 \textit{Class $\mathbf{G}$: Epstein zeta degenerate
 ($\Delta = 0$, $\operatorname{disc} = 0$)}} \\
\hline
Heisenberg $\cH_k$
 & $\mathbf{G}$ & $0$ & $\bQ$ & n/a & $\zeta(s)$ \\
\hline
Lattice $V_\Lambda$
 & $\mathbf{G}$ & $0$ & $\bQ$ & n/a & $\zeta(s)$ \\
\hline
Free fermion
 & $\mathbf{G}$ & $0$ & $\bQ$ & n/a & $\zeta(s)$ \\
\hline
\multicolumn{6}{|c|}{%
 \textit{Class $\mathbf{L}$: Epstein zeta degenerate
 ($\Delta = 0$, $\operatorname{disc} = 0$)}} \\
\hline
$\widehat{\fg}_k$ (generic)
 & $\mathbf{L}$ & $0$ & $\bQ$ & n/a & $\zeta(s)$ \\
\hline
$\widehat{\mathfrak{sl}}_2$ at $k = 1$
 & $\mathbf{L}$ & $0$ & $\bQ$ & n/a & $\zeta(s)$ \\
\hline
\multicolumn{6}{|c|}{%
 \textit{Class $\mathbf{C}$: standard conformal-weight family
 ($\beta\gamma_\lambda$ and $bc_\lambda$)}} \\
\hline
$\beta\gamma$
 & $\mathbf{C}$ & \textsuperscript{$\ast$}
 & n/a & n/a & n/a \\
\hline
\multicolumn{6}{|c|}{%
 \textit{Class $\mathbf{M}$: genuine imaginary quadratic field
 ($\Delta \neq 0$)}} \\
\hline
$\mathrm{Vir}_{1/2}$ (Ising)
 & $\mathbf{M}$ & $-160/49$ & $\bQ(\sqrt{-10})$
 & $-40$ & $L(s, \chi_{-40})$ \\[2pt]
\hline
$\mathrm{Vir}_1$
 & $\mathbf{M}$ & $-320/27$ & $\bQ(\sqrt{-15})$
 & $-15$ & $L(s, \chi_{-15})$ \\[2pt]
\hline
$\mathrm{Vir}_{13}$ (self-dual)
 & $\mathbf{M}$ & $-54080/87$ & $\bQ(\sqrt{-435})$
 & $-435$\textsuperscript{$\dagger$}
 & $L(s, \chi_{-435})$\textsuperscript{$\dagger$} \\[2pt]
\hline
$\mathrm{Vir}_{25}$
 & $\mathbf{M}$ & $-200000/147$
 & $\bQ(\sqrt{-15})$
 & $-15$
 & $L(s, \chi_{-15})$ \\[2pt]
\hline
$\cW_3$ (T-line)
 & $\mathbf{M}$
 & $-320c^2/(5c{+}22)$
 & \multicolumn{3}{c|}{Equals Virasoro data (autonomous)} \\
\hline
\end{tabular}
}% end small

\smallskip\noindent
{\small
${}^*$\,The class-C witness is the standard conformal-weight family:
$S_2 = 6\lambda^2 - 6\lambda + 1$, $S_3 = 0$,
$S_4 = -5/12$, $S_r = 0$ for $r \geq 5$.
The internal weight-changing line has $\kappa = 0$, so the shadow
metric (Definition~\ref{def:shadow-metric}) is not the witness that
detects class~$\mathbf{C}$.\\
${}^\dagger$\,$435 = 3 \cdot 5 \cdot 29$;
$d = -435$ is a fundamental discriminant ($-435 \equiv 1 \pmod{4}$,
squarefree) with $h(-435) = 4$
\textup{(}Computation~\textup{\ref{comp:fe-minimal-models}},
\texttt{shadow\_epstein\_zeta.py}\textup{)}.
}
\end{table}

\noindent
Each entry is computed from $\operatorname{disc}(Q_L)
= -320c^2/(5c{+}22)$ for the Virasoro T-line
(eq.~\eqref{eq:shadow-metric-discriminant}),
with the shadow field $K_L = \bQ(\sqrt{\delta})$
where $\delta$ is the squarefree part of the discriminant
in $\bQ^*/(\bQ^*)^2$.
The fundamental discriminant~$d$ is the discriminant
of the ring of integers $\cO_{K_L}$; when $\delta \equiv 1
\pmod 4$, $d = \delta$, otherwise $d = 4\delta$.
Computation:
\begin{itemize}[nosep]
\item $c = 1/2$:
 $\operatorname{disc} = -320/(4 \cdot 49/2) = -160/49$.
 $160 = 2^5 \cdot 5$, $49 = 7^2$.
 In $\bQ^*/(\bQ^*)^2$: $-160/49 \sim -2 \cdot 5 = -10$.
 $-10 \equiv 2 \pmod{4}$, so $d = -40$.
\item $c = 1$:
 $\operatorname{disc} = -320/27$.
 $320 = 2^6 \cdot 5$, $27 = 3^3$.
 In $\bQ^*/(\bQ^*)^2$: $-320/27 \sim -5 \cdot 3 = -15$.
 $-15 \equiv 1 \pmod{4}$, so $d = -15$.
\item $c = 25$:
 $\operatorname{disc} = -320 \cdot 625/147 = -200000/147$.
 $200000 = 2^6 \cdot 5^5$, $147 = 3 \cdot 7^2$.
 In $\bQ^*/(\bQ^*)^2$: $-2^6 \cdot 5^5/(3 \cdot 7^2)
 \sim -5 \cdot 3 = -15$.
 $-15 \equiv 1 \pmod{4}$, so $d = -15$.
\end{itemize}

%% ================================================================
%% NEW SECTIONS: bar dimensions, obstruction table, spectral sequences
%% ================================================================

\section*{Bar complex dimensions}
\label{sec:bar-dimensions}
\index{bar complex!dimension table|textbf}

Table~\ref{tab:bar-dimensions} records the bar-dual coalgebra dimensions
$\dim(\cA^{\mathrm i})_n$ for each algebra through degree~$10$.
By the $E_2$ collapse and finite/completed duality lane
(Corollary~\ref{cor:bar-cohomology-koszul-dual}), these same numbers are
the Koszul-dual Hilbert function $\dim(\cA^!)_n$ once the dual algebraic
branch exists; see
Remark~\ref{rem:table-vs-vanishing} for the distinction from the
total bar cohomology computed by the vanishing theorems.
(For Kac--Moody algebras, the bar \emph{chain-group} dimensions are much
larger: $\dim \bar{B}^n = (\dim\mathfrak{g})^n \cdot (n{-}1)!$,
the product of the generator factor and the Orlik--Solomon form
factor; see Lemma~\ref{lem:bar-dims-level-independent}.)

These dimensions are computed from the explicit
generators-and-relations
presentations of the bar complex given in Part~\ref{part:physics-bridges};
see Section~\ref{sec:cross-ref-guide} below for precise references.
Entries marked \texttt{unknown} name a missing finite computation;
entries marked \texttt{n/a} are not defined or not applicable.

Six of the nine algebras in Table~\ref{tab:bar-dimensions} now have
closed-form generating functions (Heisenberg, free fermion, $bc$,
$\beta\gamma$, $\widehat{\mathfrak{sl}}_2$, Virasoro).
For $\widehat{\mathfrak{sl}}_3$, the degree-$4$ bar cohomology
requires computing the rank of a $24576 \times 786432$ matrix
(the chain-group dimensions grow as $8^n \cdot (n{-}1)!$),
which exceeds practical hand computation. The $\mathcal{W}_3$ gap is
the missing rank computation for this matrix or an equivalent
closed-form generating function.

For the Yangian $Y(\mathfrak{sl}_2)$, the finite-character recurrence
datum
\[
a(n)=4a(n{-}1)-3a(n{-}2),\qquad a(1)=4,\quad a(2)=10,
\]
has closed form $a(n)=3^n+1$.  Its degree-$4$ coefficient is~$82$.
The bar-complex assertion at degree~$4$ is the finite rank identity
\[
\dim\ker(d^{4\to 3})-\dim\operatorname{im}(d^{5\to 4})=82.
\]
Appendix~\ref{subsec:yangian-h4-verification} records the consistency
checks that make this rank computation finite.

\begin{proposition}[Koszul dual Hilbert functions under finite-character data; \ClaimStatusConditional]
\label{prop:bar-dimensions-under-data}
Table~\textup{\ref{tab:bar-dimensions}} has the following exact reading.
Unmarked entries are the Koszul dual Hilbert dimensions computed from
the cited bar complexes. Entries marked with $\dagger$ or $*$ are exact
coefficients of the finite-character recurrence data written below; they
become bar-cohomology dimensions precisely when the corresponding
chain-level recurrence lift is supplied.
\end{proposition}

\begin{proof}
For unmarked rows, $\dim(\cA^{\mathrm i})_n$ is computed from the
$E_2$ collapse of the PBW spectral sequence
(Corollary~\ref{cor:bar-cohomology-koszul-dual}) applied to the explicit
generators-and-relations presentations established in the referenced
theorems.  The displayed $\cA^!$ Hilbert functions are the corresponding
finite/completed Verdier-dual dimensions.
The explicit formulas are:
\begin{enumerate}
\item \emph{Heisenberg} ($d = 1$, commutative):
 $\dim (\mathcal{H}_k^!)_1 = 1$; for $n \geq 2$,
 $\dim (\mathcal{H}_k^!)_n = p(n-2)$ where $p$ is the partition function.
 The underlying bar coalgebra structure is computed in
 Theorem~\ref{thm:heisenberg-bar}; the Hilbert function follows
 from Corollary~\ref{cor:bar-cohomology-koszul-dual}.
\item \emph{Free fermion} ($d = 1$, fermionic):
 $\dim (\psi^!)_n = p(n-1)$
 (Theorem~\ref{thm:fermion-bar-coalg},
 Corollary~\ref{cor:bar-cohomology-koszul-dual}).
\item \emph{$bc$ ghosts} ($d = 2$, fermionic):
 $\dim (bc^!)_n = 2^n - n + 1$
 (Theorem~\ref{thm:betagamma-bar-complex},
 Proposition~\ref{prop:bc-betagamma-orthogonality},
 Corollary~\ref{cor:bar-cohomology-koszul-dual}).
\item \emph{$\beta\gamma$ system} ($d = 2$, bosonic):
 at the lowest-weight truncation,
 $\dim (\beta\gamma^!)_n = [x^n]\sqrt{(1{+}x)/(1{-}3x)}$,
 with recurrence $n \cdot a(n) = 2n \cdot a(n{-}1) + 3(n{-}2) \cdot a(n{-}2)$
 (Theorem~\ref{thm:ds-bar-gf-discriminant},
 Corollary~\ref{cor:bar-cohomology-koszul-dual}).
 See Remark~\ref{rem:betagamma-conventions}
 for the relation between different truncation conventions.
\item \emph{$\widehat{\mathfrak{sl}}_2$}:
 The bar cohomology is computed by the $E_2$ collapse
 of the PBW spectral sequence
 (Corollary~\ref{cor:bar-cohomology-koszul-dual},
 Remark~\ref{rem:ce-vs-exterior}): the $d_1$
 differential is the CE differential of
 $\mathfrak{sl}_2 \otimes t^{-1}\mathbb{C}[t^{-1}]$,
 giving $\dim H^n = 2n+1$ for all $n \geq 1$
 (Computation~\ref{comp:sl2-ce-verification},
 Remark~\ref{rem:garland-lepowsky-sl2}).
 The earlier Riordan formula
 $\dim (\widehat{\mathfrak{sl}}_2^!)_n = R(n+3)$
 (Riordan numbers, OEIS~A005043) agrees at $n = 1$
 but gives $R(5) = 6 \neq 5$ at $n = 2$ and
 $R(6) = 15 \neq 7$ at $n = 3$; the Riordan generating
 function is superseded by the rational GF
 $x(3-x)/(1-x)^2 = \sum_{n \geq 1}(2n+1)x^n$
 (see Remark~\ref{rem:garland-lepowsky-sl2} and
 Appendix~\ref{subsec:riordan}).
\item \emph{$\widehat{\mathfrak{sl}}_3$}: computed through
 degree~$3$ from the PBW basis.  The marked degrees are generated by
 the finite-character datum~\eqref{eq:sl3-bar-recurrence}.
\item \emph{Virasoro}:
 $\dim (\mathrm{Vir}_c^!)_n = M(n{+}1) - M(n)$
 where $M(n)$ is the $n$-th Motzkin number
 (OEIS~A001006). This reflects the associahedron
 structure of $\overline{M}_{0,n+2}$, proved via
 chiral Koszulness of $\mathrm{Vir}_c$
 (Theorem~\ref{thm:virasoro-chiral-koszul}) and
 $E_2$ collapse
 (Corollary~\ref{cor:bar-cohomology-koszul-dual}).
\item \emph{$\mathcal{W}_3$}: computed through degree~$3$
 in \S\ref{sec:w3-bar-degree3}.  The marked degrees are generated by
 the recurrence $a(n) = 3a(n{-}1) + a(n{-}2) - 1$
 with $a(0)=0$, $a(1)=2$; hence the degree-$4$ and degree-$5$
 recurrence coefficients are $52$ and $171$.
 Appendix~\ref{subsec:w3-h5-verification} records the finite checks
 isolating the required degree-$5$ rank computation.\textsuperscript{$\dagger$}
\item \emph{Yangian $Y(\mathfrak{sl}_2)$}: computed through
 degree~$3$ from the RTT presentation
 (Theorem~\ref{thm:yangian-bar-rtt}).
 The recurrence datum $a(n)=4a(n{-}1)-3a(n{-}2)$ with
 $a(1)=4$, $a(2)=10$ has closed form $a(n)=3^n+1$ and rational
 generating function
 $(1-3x^2)/((1-x)(1-3x))$
 for the marked degrees.
\end{enumerate}
\end{proof}

\begin{table}[ht]
\centering
\caption{Bar-dual dimensions $\dim(\cA^{\mathrm i})_n$ through degree~$10$
\textup{(}equivalently $\dim(\cA^!)_n$ after the finite/completed dual
branch exists, and conformal-weight-graded bar cohomology;
see Remark~\ref{rem:table-vs-vanishing})}
\label{tab:bar-dimensions}
\index{bar complex!growth rate}
\renewcommand{\arraystretch}{1.4}
{\scriptsize
\begin{tabular}{|l|c|c|c|c|c|c|c|c|l|}
\hline
\textbf{Algebra $\cA$}
 & \textbf{1} & \textbf{2} & \textbf{3}
 & \textbf{4} & \textbf{5} & \textbf{6}
 & \textbf{7} & \textbf{8}
 & \textbf{Growth} \\
\hline
Heisenberg ($\operatorname{rk}1$)
 & 1 & 1 & 1 & 2 & 3 & 5 & 7 & 11 & sub-exp \\
\hline
Free fermion $\psi$
 & 1 & 1 & 2 & 3 & 5 & 7 & 11 & 15 & sub-exp \\
\hline
$bc$ ghosts
 & 2 & 3 & 6 & 13 & 28 & 59 & 122 & 249 & $2^n$ \\
\hline
$\beta\gamma$ system
 & 2 & 4 & 10 & 26 & 70 & 192 & 534 & 1500 & $3^n$ \\
\hline
$\widehat{\mathfrak{sl}}_2$ (level~$k$)
 & 3 & 5 & 7 & 9 & 11 & 13 & 15 & 17 & $2n{+}1$ \\
\hline
$\widehat{\mathfrak{sl}}_3$ (level~$k$)
 & 8 & 36 & 204 & 1352$^\dagger$ & 9892$^\dagger$ & 76084$^\dagger$ & 598592$^\dagger$ & 4755444$^\dagger$ & $8^n$ \\
\hline
$\mathrm{Vir}_c$
 & 1 & 2 & 5 & 12 & 30 & 76 & 196 & 512 & $3^n$ \\
\hline
$\mathcal{W}_3$
 & 2 & 5 & 16 & 52 & 171$^\dagger$ & 564$^\dagger$ & 1862$^\dagger$ & 6149$^\dagger$ & $3.3^n$ \\
\hline
Yangian $Y(\mathfrak{sl}_2)$
 & 4 & 10 & 28 & 82$^*$ & 244$^*$ & 730$^*$ & 2188$^*$ & 6562$^*$ & $3^n$ \\
\hline
\end{tabular}
}

\medskip
\noindent
{\footnotesize $^*$\,Exact coefficients of the Yangian recurrence datum
$a(n)=4a(n{-}1)-3a(n{-}2)$, $a(1)=4$, $a(2)=10$.\\
$^\dagger$\,Exact coefficients of the finite-character recurrence data:
$\widehat{\mathfrak{sl}}_3$ via recurrence~\eqref{eq:sl3-bar-recurrence},
$\mathcal{W}_3$ via recurrence~\eqref{eq:w3-bar-recurrence}.
Equality with bar cohomology from the first marked degree is the
corresponding rank computation for the chiral bar differential.}
\end{table}

\begin{remark}[Finite recurrence reading of \texorpdfstring{$H^4(\mathcal{W}_3) = 52$}{H4(W3) = 52}]
\label{rem:w3-h4-recurrence}
\index{W3-algebra@$\mathcal{W}_3$!degree-4 bar cohomology}
The value $52$ is the fourth coefficient of the unique depth-$2$
recurrence with constant $a(n) = 3a(n{-}1) + a(n{-}2) - 1$
fitted to the four computed values $\{0, 2, 5, 16\}$
$($with $a(0)=0)$.
Four data points determine this recurrence but do not determine the
ambient generating function among all algebraic alternatives; the
recurrence is computed, and no derivation from the chiral bar
differential is known.  The
bar-cohomology equality at degree~$4$ is the rank identity obtained by
direct bar differential computation, by a generators-and-relations
analysis of
$\barB(\mathcal{W}_3)_4$, or by extending the PBW spectral sequence
approach of Proposition~\ref{prop:sl3-pbw-ss} to the $\mathcal{W}_3$
setting.
\end{remark}

\begin{remark}[Table values vs.\ genus-\texorpdfstring{$0$}{0} vanishing]\label{rem:table-vs-vanishing}
\index{bar complex!table vs vanishing}
The dimensions in Table~\ref{tab:bar-dimensions} report the
Hilbert function of the Koszul dual algebra~$\cA^!$, which by the
$E_2$ collapse (Corollary~\ref{cor:bar-cohomology-koszul-dual})
equals the conformal-weight-graded bar cohomology
$\sum_h \dim H^{n,h}(\bar{B}(\cA))$.
These are nonzero because the conformal weight grading
prevents cancellation between different weight spaces.

The vanishing theorems for total bar cohomology at genus~$0$
(e.g., Theorem~\ref{thm:heisenberg-bar}: $H^n = 0$ for $n > 2$)
compute the \emph{total} cohomology $H^n = \bigoplus_h H^{n,h}$
after summing over all conformal weights with signs from the
differential; the weight-graded dimensions in the table are the
input to the spectral sequence whose output is the total
cohomology. In particular,
$\dim(\cA^{\mathrm i})_1=\dim(\cA^!)_1 = 1$ for the
Heisenberg algebra after the dual branch (the single
generator~$\alpha$), while the
total $H^1(\bar{B}_{\mathrm{geom}}(\mathcal{H}_k)) = 0$
on~$\mathbb{P}^1$ because $H^1(\mathbb{P}^1) = 0$.
\end{remark}

\begin{remark}[Growth rates and Koszulness]\label{rem:bar-growth-koszul}
\index{Koszul algebra!bar complex growth}
The growth rate of $\dim(\cA^{\mathrm i})_n$ as $n \to \infty$ is
an invariant of the bar-dual coalgebra with direct Koszul-theoretic
meaning; after the finite/completed dual branch it is the growth rate of
$\dim(\cA^!)_n$.
(Note: the bar \emph{chain-group} dimensions grow much faster;
for Kac--Moody algebras $\dim \bar{B}^n = (\dim\fg)^n \cdot (n{-}1)!$
is super-exponential. The growth rates below refer to
the Koszul dual Hilbert function.)
\begin{enumerate}
\item \emph{Sub-exponential growth.}
 The Heisenberg and free fermion Koszul duals grow
 sub-exponentially: $\dim (\cA^!)_n \sim
 e^{\pi\sqrt{2n/3}}/(4\sqrt{3}\,n)$ by the
 Hardy--Ramanujan asymptotics for the partition function.
\item \emph{Exponential growth.}
 All other algebras in the table have Koszul dual dimensions growing
 exponentially in~$n$, at rate~$\sim d^n$ where $d = \dim\fg$
 for Kac--Moody.
 The explicit formulas are:
 \begin{itemize}
 \item $bc$: $\dim (bc^!)_n = 2^n - n + 1$
 (growth rate~$2^n$);
 \item $\beta\gamma$: $\dim (\beta\gamma^!)_n = [x^n]\sqrt{(1{+}x)/(1{-}3x)}$
 (growth rate~$3^n/n^{1/2}$);
 \item $\widehat{\mathfrak{sl}}_2$:
 $\dim H^n = 2n+1$ for all $n \geq 1$
 (Garland--Lepowsky concentration;
 rational GF $x(3{-}x)/(1{-}x)^2$,
 growth rate linear;
 Remark~\ref{rem:garland-lepowsky-sl2});
 \item Virasoro: $\dim (\mathrm{Vir}_c^!)_n = M(n{+}1) - M(n)$,
 the first differences of Motzkin numbers
 (growth rate~$3^n/n^{3/2}$);
 \item Yangian $Y(\mathfrak{sl}_2)$: under the finite-character
 recurrence datum $a(n)=4a(n{-}1)-3a(n{-}2)$ with
 $a(1)=4$, $a(2)=10$, the marked coefficients are
 $a(n)=3^n+1$ (growth rate~$3^n$).
 \end{itemize}
\item In general, for a Koszul algebra, the exponential
 growth rate of $\dim (\cA^!)_n$ equals the
 exponential growth rate of the Koszul dual $\cA^!$.
 The Virasoro bar cohomology growth~$3^n$ is consistent
 with the infinite-dimensional same-family shadow partner
 $\mathrm{Vir}_{26-c}$ that controls the current M/S-level
 calculations. The MC4 structural theorem is
 Theorem~\ref{thm:completed-bar-cobar-strong}; the H-level
 infinite-generator identification remains open.
\end{enumerate}
\end{remark}

\begin{remark}[Garland--Lepowsky concentration for $\widehat{\mathfrak{sl}}_2$]
\label{rem:garland-lepowsky-sl2}
\index{Garland--Lepowsky concentration}
\index{bar complex!$\widehat{\mathfrak{sl}}_2$!linear growth}
The bar cohomology of the affine $\widehat{\mathfrak{sl}}_2$
vacuum module satisfies $\dim H^n = 2n+1$ for all $n \geq 1$,
with rational generating function
\[
 P_{\widehat{\mathfrak{sl}}_2}(x)
 \;=\; \frac{x(3-x)}{(1-x)^2}
 \;=\; 3x + 5x^2 + 7x^3 + 9x^4 + \cdots.
\]
This is a manifestation of the Garland--Lepowsky concentration
phenomenon: the CE differential on
$\Lambda^*(\mathfrak{sl}_2 \otimes t^{-1}\mathbb{C}[t^{-1}])^*$
kills the vast majority of chains at each bar degree,
leaving only $2n+1$ survivors at degree~$n$. The
mechanism is the semisimplicity of~$\mathfrak{sl}_2$:
the Whitehead lemma forces $H^1(\mathfrak{sl}_2; M) = 0$
for every finite-dimensional module, so the PBW
spectral sequence $E_1$ page is concentrated on the
invariant subalgebra
$S^*(\mathfrak{sl}_2[t^{-1}]t^{-1})^{\mathfrak{sl}_2}$,
which is generated by a single quadratic Casimir
(one generator per negative mode $t^{-m}$, $m \geq 1$).
The resulting bar cohomology at degree~$n$ and weight~$H$
counts the $\mathfrak{sl}_2$-invariant elements
in $\Lambda^n$ of the loop algebra restricted to
total weight~$H$; the net count is $2n+1$
(Computation~\ref{comp:sl2-ce-verification}).

The linear growth $\dim H^n = 2n+1$
supersedes the Riordan formula
$\dim H^n = R(n+3)$ (OEIS~A005043), which gives the
correct value at $n = 1$ ($R(4) = 3$) but
fails at $n = 2$ ($R(5) = 6 \neq 5$) and
at $n = 3$ ($R(6) = 15 \neq 7$).
A detailed treatment of the three concentration mechanisms
(Whitehead vanishing, Casimir generation, and
weight-by-weight cancellation) appears in the
companion paper~\cite{Lor-GL}.
\end{remark}

\begin{corollary}[Sub-exponential growth in the computed rows; \ClaimStatusProvedHere]
\label{cor:subexp-free-field}
\index{bar complex!sub-exponential growth}
Among the computed bar-complex rows in the landscape census, the
bar-dual dimensions
$\dim(\cA^{\mathrm i})_n$ \textup{(}equivalently the Koszul-dual Hilbert
function $\dim(\cA^!)_n$ after the dual branch\textup{)} grow
sub-exponentially exactly for the free fields (Heisenberg or free
fermion) and affine
$\widehat{\mathfrak{sl}}_2$ (linear growth $2n{+}1$).  The rows with
finite-character recurrence data have the growth rates displayed by
their recurrence polynomials.
\end{corollary}

\begin{proof}
The Heisenberg and free fermion dimensions are $p(n-2)$ and $p(n-1)$ respectively, both sub-exponential by Hardy--Ramanujan asymptotics ($p(n) \sim e^{\pi\sqrt{2n/3}}/(4\sqrt{3}\,n)$). The affine $\widehat{\mathfrak{sl}}_2$ has $\dim H^n = 2n+1$, linear in~$n$ by Computation~\ref{comp:sl2-ce-verification}. All other interacting entries in the computed list grow exponentially: $2^n$ for $bc$, $3^n$ for $\beta\gamma$ and Virasoro. The Heisenberg and free fermion are the unique free-field (non-interacting) algebras among the computed examples: they have no nontrivial OPE beyond the central term.
\end{proof}

\begin{corollary}[Closed forms for computed interacting rows; \ClaimStatusProvedHere]
\label{cor:algebraicity-koszul}
\index{bar complex!generating function!algebraicity}
Every computed interacting row with a closed form proved from the bar
complex has a rational or quadratic algebraic generating function
over~$\mathbb{Q}(x)$.  The marked
higher-rank and Yangian rows are rational functions only after the
finite-character recurrence datum is imposed; that datum is additional
bridge structure, not a formal consequence of Koszulness alone.
\end{corollary}

\begin{proof}
The interacting rows proved directly here have the following
generating functions:
$bc$ has $\sum (2^n-n+1)x^n$, hence a rational function;
$\beta\gamma$ has $\sqrt{(1+x)/(1-3x)}$, algebraic of degree~$2$;
$\widehat{\mathfrak{sl}}_2$ has $x(3-x)/(1-x)^2$, rational;
and Virasoro has the Motzkin first-difference series, algebraic of
degree~$2$. Hence every proved interacting row has a rational or
quadratic algebraic generating function.
For $\widehat{\mathfrak{sl}}_2$ the GF is rational, a special
quadratic-algebraic case reflecting the Garland--Lepowsky concentration to
linear growth verified by Computation~\ref{comp:sl2-ce-verification}.
\end{proof}

\begin{theorem}[Rank-one discriminant chart and finite-character denominator data; \ClaimStatusConditional]
\label{thm:ds-bar-gf-discriminant}
\index{Drinfeld--Sokolov reduction!bar cohomology generating function}
\index{bar complex!generating function!shared discriminant}
\index{Riordan numbers!relation to Motzkin numbers}
Let $\mathfrak{g}$ be a simple Lie algebra and
$\mathcal{W}^k(\mathfrak{g}) = H^0_{\mathrm{DS}}(\widehat{\mathfrak{g}}_k)$
the associated $\mathcal{W}$-algebra.
Define the Koszul dual Hilbert series
\[
P_{\widehat{\mathfrak{g}}}(x)
 = \sum_{n \geq 1} \dim (\widehat{\mathfrak{g}}_k^!)_n\, x^n,
\qquad
P_{\mathcal{W}}(x)
 = \sum_{n \geq 1} \dim (\mathcal{W}^k(\mathfrak{g})^!)_n\, x^n.
\]
The following statements hold in rank one and in the displayed
finite-character charts.

Explicitly, for $\mathfrak{g} = \mathfrak{sl}_2$,
the Virasoro and $\beta\gamma$ generating functions are:
\begin{align}
P_{\mathrm{Vir}}(x)
 &= \frac{4x}{(1 - x + \sqrt{1 - 2x - 3x^2})^2},
 \label{eq:gf-vir-motzkin} \\
P_{\beta\gamma}(x)
 &= \sqrt{\frac{1+x}{1-3x}}
 = \frac{1+x}{\sqrt{1-2x-3x^2}}.
 \label{eq:gf-betagamma}
\end{align}

Both are algebraic functions of the discriminant
$\Delta(x) = 1 - 2x - 3x^2 = (1-3x)(1+x)$.
The $\widehat{\mathfrak{sl}}_2$ generating function is
\emph{rational} and belongs to a different coefficient-recurrence
module (Remark~\ref{rem:garland-lepowsky-sl2}):
\begin{equation}\label{eq:gf-sl2-rational}
P_{\widehat{\mathfrak{sl}}_2}(x)
 \;=\; \frac{x(3-x)}{(1-x)^2}
 \;=\; \sum_{n \geq 1}(2n+1)\,x^n,
\end{equation}
whereas the Riordan algebraic function
\begin{equation}\label{eq:gf-sl2-riordan}
P_{\widehat{\mathfrak{sl}}_2}(x)
 \stackrel{\times}{=} \frac{1 + x - \sqrt{1 - 2x - 3x^2}}{2x(1+x)}
\end{equation}
is not the affine $\widehat{\mathfrak{sl}}_2$ bar series: it yields
$\dim(\widehat{\mathfrak{sl}}_2^!)_2 = 6$, conflicting
with the proved value $\dim H^2 = 5$
(Computation~\ref{comp:sl2-ce-verification},
Lemma~\ref{lem:bar-deg2-symmetric-square}),
and $\dim(\widehat{\mathfrak{sl}}_2^!)_3 = 15$, conflicting with
the proved $\dim H^3 = 7$.
The Virasoro and $\beta\gamma$ generating functions are ramified at
$x = 1/3$ and $x = -1$, sharing the
exponential growth rate~$3^n$
(with sub-exponential corrections $n^{-3/2}$ for
Virasoro, $n^{-1/2}$ for~$\beta\gamma$).
By contrast, $\widehat{\mathfrak{sl}}_2$ has linear growth
$\dim H^n = 2n + 1$, a consequence of the
Garland--Lepowsky concentration
(Remark~\ref{rem:garland-lepowsky-sl2}).

For the $\mathfrak{sl}_3/\mathcal{W}_3$ finite-character chart, the
shift polynomials
\[
p_{\widehat{\mathfrak{sl}}_3}(t)=(t-8)(t^2-3t-1),\qquad
p_{\mathcal{W}_3}(t)=(t-1)(t^2-3t-1)
\]
have common divisor $t^2-3t-1$.  The corresponding denominator
factor in the $x$-plane is $1-3x-x^2$.
\end{theorem}

\begin{proof}
The proved formulas~\eqref{eq:gf-vir-motzkin}
and~\eqref{eq:gf-betagamma} and the shared-discriminant
claim are established as follows.

\emph{$\widehat{\mathfrak{sl}}_2$.}
The corrected PBW analysis
(Corollary~\ref{cor:bar-cohomology-koszul-dual},
Remark~\ref{rem:ce-vs-exterior}) identifies the
bar cohomology with the CE cohomology of the loop algebra
$\mathfrak{sl}_2 \otimes t^{-1}\mathbb{C}[t^{-1}]$;
the Garland--Lepowsky concentration
(Remark~\ref{rem:garland-lepowsky-sl2}) gives
$\dim H^n = 2n+1$ for all $n \geq 1$, with rational
generating function~\eqref{eq:gf-sl2-rational}.
The original Riordan formula
$\dim (\widehat{\mathfrak{sl}}_{2,k}^!)_n = R(n+3)$
disagrees at $n = 2$ ($R(5) = 6 \neq 5$)
and at $n = 3$ ($R(6) = 15 \neq 7$);
the discriminant $\Delta(x) = (1-3x)(1+x)$
of the DS reduction still governs the Virasoro and
$\beta\gamma$ generating functions below.

\emph{Virasoro (Motzkin differences).}
The Koszul dual dimension $\dim (\mathrm{Vir}_c^!)_n
= M(n{+}1) - M(n)$ where $M(n)$ are the Motzkin numbers
(Theorem~\ref{thm:spectral-sequence-collapse}). The Motzkin
generating function is $M(x) = (1 - x - \sqrt{1-2x-3x^2})/(2x^2)$
(OEIS~A001006). Then
$\sum_{n \geq 0}(M(n{+}1) - M(n))x^n = (M(x) - 1)/x - M(x) + 1$;
simplification gives~\eqref{eq:gf-vir-motzkin}.

\emph{$\beta\gamma$ system.}
At the lowest-weight truncation
(Remark~\ref{rem:betagamma-conventions}), the bar cohomology
satisfies the recurrence
$n \cdot a(n) = 2n \cdot a(n{-}1) + 3(n{-}2) \cdot a(n{-}2)$
with $a(0) = 1$, $a(1) = 2$. This is the coefficient recurrence
for the ODE $(1{+}x)(1{-}3x)H'(x) = 2H(x)$, whose solution with
$H(0) = 1$ is $H(x) = \sqrt{(1{+}x)/(1{-}3x)}$,
giving~\eqref{eq:gf-betagamma}.
The first values $2, 4, 10, 26, 70, 192, 534, 1500$
are OEIS~A025565 (shifted by one).

\emph{Rank-one discriminant.}
The Virasoro and $\beta\gamma$ functions are algebraic of degree~$2$
over $\mathbb{Q}(x)$, with branch locus
$\Delta(x) = 1 - 2x - 3x^2 = 0$, i.e.\ $x = 1/3$ and $x = -1$.
The branch point at $x = 1/3$ determines the exponential growth rate
$3^n$ in these two sequences.  The sub-exponential correction
$n^{-3/2}$ for Virasoro follows from the square-root singularity
(Theorem~VI.1 of~\cite{FlajoletSedgewick}); the $\beta\gamma$
correction is $n^{-1/2}$, reflecting the $\Delta^{-1/2}$ singularity
type.  The affine $\widehat{\mathfrak{sl}}_2$ series is rational with
denominator $(1-x)^2$; the polynomial $\Delta$ is its DS branch divisor,
not its coefficient-recurrence denominator.

\emph{DS and Wakimoto bridge data.}
Since $\mathrm{Vir}_c = \mathcal{W}^k(\mathfrak{sl}_2)$ by
Drinfeld--Sokolov reduction, the functorial commutative
diagram~\eqref{eq:ds-kd-square}
(Theorem~\ref{thm:ds-koszul-intertwine}) gives
$\barBgeom(\mathrm{Vir}_c) \simeq
H^0_{\mathrm{DS}}(\barBgeom(\widehat{\mathfrak{sl}}_{2,k}))$.
This identifies the BRST bridge, while the displayed formulas identify
the branch divisor that survives on the Virasoro and $\beta\gamma$
side.  The affine coefficient sequence itself is governed by the
rational denominator $(1-x)^2$.

The $\beta\gamma$ system enters the same discriminant family
through the Wakimoto representation:
$\widehat{\mathfrak{sl}}_{2,k}$ embeds into
$\beta\gamma \otimes \mathcal{H}$ (Theorem~\ref{thm:wakimoto-bar}),
and the embedding preserves the algebraic structure of bar
generating functions. Since the Heisenberg factor~$\mathcal{H}$
contributes only sub-exponential growth
($\dim (\mathcal{H}^!)_n = p(n{-}2)$, the partition function),
the exponential growth rate~$3^n$ and the branch divisor~$\Delta(x)$
are carried by the Virasoro quotient (via DS) and the $\beta\gamma$
factor (via Wakimoto), not by the corrected affine
$\widehat{\mathfrak{sl}}_2$ coefficient recurrence.
\end{proof}

\begin{remark}[DS denominator chart]\label{rem:ds-discriminant-invariant}
\index{Drinfeld--Sokolov reduction!discriminant invariant}
A DS denominator chart consists of finite-character recurrence modules
for $\widehat{\mathfrak{g}}_k$ and $\mathcal{W}^k(\mathfrak{g})$
together with a BRST-compatible map of their divisor cores.  The
rank-one chart separates two facts: the corrected affine
$\widehat{\mathfrak{sl}}_2$ coefficient series has denominator
$(1-x)^2$, while the Virasoro and $\beta\gamma$ branch divisor is
$(1-3x)(1+x)$.  The $A_2$ chart records the common divisor
$t^2-3t-1$ of the two finite-character shift polynomials; in the
$x$-plane this is the shared denominator factor $1-3x-x^2$.
\end{remark}

\begin{theorem}[Divisor-core form of DS sub-discriminance; \ClaimStatusProvedHere]
\label{thm:ds-spectral-branch-preservation}
\index{Drinfeld--Sokolov reduction!spectral branch preservation}
\index{spectral branch object!DS sub-discriminant}
Let $\mathfrak{g}$ be a simple Lie algebra,
$\widehat{\mathfrak{g}}_k$ the associated Kac--Moody algebra, and
$\mathcal{W}^k(\mathfrak{g}) = H^0_{\mathrm{DS}}(\widehat{\mathfrak{g}}_k)$
its principal $\mathcal{W}$-algebra.  Assume finite-character
recurrence modules
\[
\mathcal{R}_{\widehat{\mathfrak{g}}}\cong k[t]/(p_{\widehat{\mathfrak{g}}}),
\qquad
\mathcal{R}_{\mathcal{W}}\cong k[t]/(p_{\mathcal{W}})
\]
and a DS-compatible map whose image on cyclic divisor cores is the
maximal common core.  If
$s(t)=\gcd(p_{\widehat{\mathfrak{g}}},p_{\mathcal{W}})$, then the
common sector is
\[
\mathcal{C}_s(p_{\widehat{\mathfrak{g}}})
\;\cong\; k[t]/(s)
\;\cong\; \mathcal{C}_s(p_{\mathcal{W}}).
\]
For the induced operator $KS_s$ on the dual sector,
\[
\det(1-x\,KS_s)=x^{\deg s}s(x^{-1}).
\]
\end{theorem}

\begin{proof}
The divisor-core theorem gives
$\mathcal{C}_s(p)\cong k[t]/(s)$ for every divisor $s\mid p$.
Applying it to $p_{\widehat{\mathfrak{g}}}$ and $p_{\mathcal{W}}$
with their maximal common divisor gives the displayed isomorphism.
The determinant formula is the characteristic polynomial calculation
for multiplication by~$t$ on $k[t]/(s)$, written in reciprocal
$x$-coordinates.
\end{proof}

\begin{remark}[Spectral decomposition of the branch transport]
\label{rem:ds-spectral-decomposition}
\index{spectral branch object!spectral decomposition}
Theorem~\ref{thm:ds-spectral-branch-preservation} refines the
denominator chart to a spectral decomposition: the finite-character
transport space decomposes as
$V_{\cA}^{\mathrm{br}} = V_{\mathrm{dom}} \oplus
V_{\mathrm{inv}}$, where $V_{\mathrm{dom}}$ carries the
algebra-specific dominant eigenvalue and $V_{\mathrm{inv}}$
carries the DS-invariant sub-discriminant. The DS reduction
acts as the identity on $V_{\mathrm{inv}}$ and replaces the
dominant eigenvalue on $V_{\mathrm{dom}}$.
\end{remark}

\begin{hypothesis}[Finite-character Kodaira--Spencer datum]
\label{conj:discriminant-ks-operator}
\label{hyp:discriminant-ks-operator}
\index{discriminant!Kodaira--Spencer operator}
\index{Kodaira--Spencer operator!discriminant}
For a simple Lie algebra $\mathfrak{g}$, the genus-$1$ component
$\Theta_{\cA}^{(1)}$ of the universal Maurer--Cartan class
is required to induce a finite-rank \emph{Kodaira--Spencer transport operator}
$\mathrm{KS}_{\mathfrak{g}}$ on a natural subquotient
$H^{\mathrm{red}}_1(\cA)$ of the first bar cohomology
$H^1(\bar{B}(\cA))$, of dimension
$\operatorname{rank}(\mathfrak{g}) + 1$, such that:
\begin{equation}\label{eq:discriminant-ks-operator}
\Delta_{\mathfrak{g}}(x)
= \det\bigl(1 - x \cdot \mathrm{KS}_{\mathfrak{g}}
 \mid H^{\mathrm{red}}_1(\cA)\bigr).
\end{equation}
For a lifted finite-character datum, the largest eigenvalue of
$\mathrm{KS}_{\mathfrak{g}}$ is $\dim\mathfrak{g}$
(Theorem~\ref{thm:dominant-branch-point}).

\noindent\emph{Rank-one model.}
For $\mathfrak{g} = \mathfrak{sl}_2$:
$\Delta = (1-3x)(1+x)$, so
$\mathrm{KS}_{\mathfrak{sl}_2}$
acts on the $2$-dimensional branch-divisor module with eigenvalues~$3$
and~$-1$.  The dominant eigenvalue satisfies
$3 = \dim \mathfrak{sl}_2$
(Theorem~\ref{thm:dominant-branch-point}); the
eigenvalue~$-1$ corresponds to the branch point of
$P_{\mathrm{Vir}}$ and $P_{\beta\gamma}$ at $x = -1$.
The corrected affine $\widehat{\mathfrak{sl}}_2$ coefficient series has
denominator $(1-x)^2$ and is not this branch-divisor module.

For $\mathfrak{g} = \mathfrak{sl}_3$ under the finite-character
datum~\eqref{eq:sl3-bar-recurrence}, the denominator
$\Delta_{\mathfrak{sl}_3} = (1 - 8x)(1 - 3x - x^2)$
gives a $3$-dimensional space with eigenvalues~$8 = \dim \mathfrak{sl}_3$,
$(3 {+} \sqrt{13})/2 \approx 3.303$, and
$(3 {-} \sqrt{13})/2 \approx -0.303$.
The sub-dominant eigenvalues encode the divisor core shared with
$\mathcal{W}_3$.

\emph{Why not $T_{\mathrm{KS}}$ on full $H^1$?}
A natural first attempt is to take $\mathrm{KS}_{\mathfrak{g}}$
acting on $H^1(\bar{B}(\cA)) \cong \mathfrak{g}$ directly.
But $T_{\mathrm{KS}}$ commutes with the $\mathfrak{g}$-action
(the vertex algebra structure is equivariant), and by Schur's lemma
on the irreducible adjoint representation, $T_{\mathrm{KS}}$
is scalar on~$H^1$, giving
$\det(1 - x\, T_{\mathrm{KS}}) = (1 - cx)^{\dim\mathfrak{g}}$,
which has degree $\dim\mathfrak{g}$ rather than
$\operatorname{rank}(\mathfrak{g}) + 1$.
The correct space $H^{\mathrm{red}}_1$ must
therefore be a \emph{subquotient} of $H^1$, not $H^1$ itself.
The required bridge data are the representation-theoretic construction
of this subquotient and the shift-compatible action of the genus-$1$
class on it.
This construction uses Theorem~\ref{thm:mc2-full-resolution}.
\end{hypothesis}

\begin{proposition}[A2 divisor-core calculation; \ClaimStatusProvedHere]
\label{prop:ds-invariant-discriminant}
Let
\[
p_{\widehat{\mathfrak{sl}}_3}(t)=(t-8)(t^2-3t-1),\qquad
p_{\mathcal{W}_3}(t)=(t-1)(t^2-3t-1)
\]
be the finite-character shift polynomials for the displayed
$\widehat{\mathfrak{sl}}_3$ and $\mathcal{W}_3$ denominator data.  Their
maximal common divisor is $s(t)=t^2-3t-1$, and the associated divisor
cores satisfy
\[
\mathcal{C}_s(p_{\widehat{\mathfrak{sl}}_3})
\cong k[t]/(s)
\cong \mathcal{C}_s(p_{\mathcal{W}_3}).
\]
On the dual core, multiplication by~$t$ has reciprocal characteristic
polynomial $1-3x-x^2$.
\end{proposition}

\begin{proof}
The polynomial identity
\[
\gcd\bigl((t-8)(t^2-3t-1),(t-1)(t^2-3t-1)\bigr)=t^2-3t-1
\]
is immediate.  The divisor-core theorem gives
$\mathcal{C}_s(p)\cong k[t]/(s)$ for each ambient polynomial~$p$.
The determinant identity is
$\det(1-xT_s)=x^2s(x^{-1})=1-3x-x^2$ for multiplication by~$t$ on
$k[t]/(s)$.
\end{proof}

\begin{remark}[Meaning of the \(A_2\) denominator core]
The preceding proposition is a theorem about finite-character
recurrence modules. It identifies the common DS denominator core
$t^2-3t-1$ and its reciprocal characteristic polynomial
$1-3x-x^2$. It does not yet identify either recurrence module with
the bar cohomology of
$\widehat{\mathfrak{sl}}_{3,k}$ or $\mathcal W_3$. That stronger
statement requires a shift-compatible DS lift from the chiral bar
complex to the recurrence module and the corresponding
bar-differential rank identities.
\end{remark}

\begin{remark}[The rank-plus-one pattern]\label{rem:rank-plus-one}
\index{discriminant!rank pattern}
The finite-character denominator degree in the two displayed Lie
families is $\operatorname{rank}(\mathfrak{g})+1$:
$\operatorname{rank}(\mathfrak{sl}_2) + 1 = 2$
(matching the rank-one branch-divisor module) and
$\operatorname{rank}(\mathfrak{sl}_3) + 1 = 3$
(matching the denominator degree of the finite-character datum
\eqref{eq:sl3-bar-recurrence}).
This is a computed observation on two families, one of which
($\widehat{\mathfrak{sl}}_3$) is itself the hypothesis
datum~\eqref{eq:sl3-bar-gf} matched to three computed values; no
derivation is known.  A heuristic mechanism, not a derivation: the
$\operatorname{rank}(\mathfrak{g})$ algebraically independent
Casimir operators $C_1, \ldots, C_r \in Z(U\mathfrak{g})$ act
on the PBW spectral sequence
(Theorem~\ref{thm:pbw-koszulness-criterion}), providing~$r$
independent constraints on the bar cohomology dimensions; together
with the central extension (one additional constraint from the
level~$k$), a finite-character lift would produce a recurrence
module of rank $r+1$, with transfer matrix
$\mathrm{KS}_{\mathfrak{g}}$.
\end{remark}

\begin{proposition}[Construction of \texorpdfstring{$H^{\mathrm{red}}_1$}{Hred_1} for
\texorpdfstring{$\mathfrak{sl}_2$}{sl_2}; \ClaimStatusConditional]
\label{prop:hred-sl2}
\index{bar complex!reduced bar cohomology!$\mathfrak{sl}_2$}
\index{Picard--Fuchs equation!bar cohomology}
For $\mathfrak{g} = \mathfrak{sl}_2$, the corrected affine bar
cohomology dimension sequence $a_n = 2n+1$ satisfies
the recurrence of order~$1$:
\begin{equation}\label{eq:sl2-bar-recurrence-corrected}
a_n \;=\; a_{n-1} + 2,
\qquad a_1 = 3.
\end{equation}
The rational generating function
$P(x) = x(3-x)/(1-x)^2$ is annihilated by the
first-order operator
$\widetilde{L} = (1-x)^2 \partial_x - (3-2x)$,
and its constant-coefficient shift module has polynomial $(t-1)^2$.
The separate \emph{branch divisor}
$\Delta(x) = (1-3x)(1+x)$ of the rank-one DS chart
(Theorem~\ref{thm:ds-bar-gf-discriminant})
governs the Virasoro and $\beta\gamma$ generating functions.  The
Picard--Fuchs $D$-module
$\mathcal{M}_{\mathfrak{sl}_2}
:= D_{\mathbb{A}^1} / D_{\mathbb{A}^1} \cdot L$,
where $L = (1 - 2x - 3x^2)\partial_x^2
+ (-1 + 3x)\partial_x + 1$
is the differential operator annihilating the
\emph{original} Riordan generating function
(retained here as it controls the DS-shared discriminant),
provides the canonical construction:
\[
H^{\mathrm{red}}_1(\widehat{\mathfrak{sl}}_2)
\;:=\; \mathcal{M}_{\mathfrak{sl}_2}\big|_{x=0}
\;\cong\; \mathbb{Q}^2.
\]
The monodromy representation of
$\pi_1(\mathbb{A}^1 \setminus \{1/3, -1\})$ acts on this
$2$-dimensional fiber, and the asymptotic companion matrix
\[
T_{\mathrm{rec}} =
\begin{pmatrix} 2 & 3 \\ 1 & 0 \end{pmatrix},
\qquad
\operatorname{Spec}(T_{\mathrm{rec}}) = \{3, -1\},
\]
recovers the discriminant:
$\det(1 - x\, T_{\mathrm{rec}}) = (1-3x)(1+x) = \Delta(x)$.
\end{proposition}

\begin{proof}
The identity $a_n=2n+1$ gives~\eqref{eq:sl2-bar-recurrence-corrected}
and the generating function $x(3-x)/(1-x)^2$ by summing the arithmetic
progression.  The branch-divisor recurrence
\eqref{eq:riordan-recurrence-explicit} follows by direct computation
from the Riordan algebraic function
(Theorem~\ref{thm:ds-bar-gf-discriminant}).
The Picard--Fuchs operator~$L$ is obtained by differentiating
the algebraic equation
$x(1{+}x)R^2 - (1{+}x)R + 1 = 0$
twice and eliminating~$R$: the result is a second-order Fuchsian
ODE with regular singular points at $x = 1/3$, $x = -1$, and
$x = \infty$. The fiber at $x = 0$ has dimension~$2$ (the
holonomic rank), and the
Poincar\'e--Perron theorem
(Theorem~\ref{thm:discriminant-spectral})
identifies the eigenvalues of~$T_{\mathrm{rec}}$ with the
reciprocals of the singular points.

Thus the rank-one sector has two branch modes, with eigenvalues
$3$ and $-1$, while the actual affine coefficient recurrence has the
double eigenvalue~$1$.
\end{proof}

\begin{remark}[Intrinsic vs.\ extrinsic: resolution]
\label{rem:hred-intrinsic}
Proposition~\ref{prop:hred-sl2} constructs $H^{\mathrm{red}}_1$
\emph{extrinsically}: the Picard--Fuchs D-module is computed from
the generating function, not directly from the bar complex.
The PBW rank-one discriminant theorem
(Theorem~\ref{thm:pbw-recurrence},
Chapter~\ref{chap:casimir-divisor}) supplies the scalar branch
divisor for $\mathfrak{sl}_2$ and the associated divisor-core
module. It does not convert the P-recursive coefficient sequence into
a global constant-coefficient shift recurrence outside the
finite-character recurrence hypothesis.

For higher rank, an intrinsic finite-character construction has two
structural inputs:
\begin{enumerate}[label=\textup{(\Roman*)}]
\item \emph{Invariant theory.}
The $r$~algebraically independent Casimir operators
$C_1,\ldots,C_r\in Z(U\mathfrak{g})$ act on the bar complex
$\barB(\widehat{\mathfrak{g}}_k)$ by functoriality, and this
action descends to bar cohomology. The PBW spectral sequence
couples the Casimir eigenspace structure to the Molien--Weyl
series~\eqref{V2-eq:molien-weyl-g}, which has $r$~independent
poles. Each pole contributes one holonomic constraint on
the bar cohomology generating function.
\item \emph{Arnold factorial.}
The Orlik--Solomon dimension
$\dim\Omega^{n-1}(\operatorname{Conf}_n)=(n{-}1)!$ contributes
one additional Euler-type constraint, raising the D-module
order by~$1$.
\end{enumerate}
Together, these inputs supply the shape of a rank $r{+}1$
finite-character module once the chain-level recurrence lift exists.
In rank~$1$, the growth-mode factorization
\textup{(}Theorem~\ref{thm:growth-mode-factorization}\textup{)}
decomposes the Casimir recurrence module into the dominant
chain-group growth mode and the DS-invariant spectral core. The same
decomposition in higher rank is exactly the missing lift datum.
\end{remark}

\begin{hypothesis}[\texorpdfstring{$\widehat{\mathfrak{sl}}_3$}{sl3-hat} finite-character recurrence]\label{conj:sl3-bar-gf}
\label{hyp:sl3-finite-character-recurrence}
\index{sl3@$\mathfrak{sl}_3$!bar generating function}
The $\widehat{\mathfrak{sl}}_3$ finite-character recurrence datum is the
rational series
\begin{equation}\label{eq:sl3-bar-gf}
P_{\widehat{\mathfrak{sl}}_3}(x)
= \frac{4x(2 - 13x - 2x^2)}{(1 - 8x)(1 - 3x - x^2)},
\end{equation}
with denominator $(1-8x)(1-3x-x^2)$ and recurrence
\begin{equation}\label{eq:sl3-bar-recurrence}
a(n) = 11\,a(n{-}1) - 23\,a(n{-}2) - 8\,a(n{-}3),
\qquad a(1) = 8,\; a(2) = 36,\; a(3) = 204.
\end{equation}

\emph{Finite checks.}

\begin{enumerate}[label=\textup{(\roman*)}]
\item The three known values $8, 36, 204$ are matched exactly.
\item The denominator factor $(1-3x-x^2)$ is shared with the
$\mathcal{W}_3$ recurrence datum
$P_{\mathcal{W}_3}(x) = x(2-3x)/\bigl((1-x)(1-3x-x^2)\bigr)$
(Remark~\ref{rem:ds-discriminant-invariant}):
$\widehat{\mathfrak{sl}}_3$ and $\mathcal{W}_3$ share the
factor $(1-3x-x^2)$ in the denominator, just as
$\widehat{\mathfrak{sl}}_2$ and $\mathrm{Vir}$ share the
discriminant $(1-3x)(1+x)$.
\item The pole at $x = 1/8$ gives exponential growth rate~$8^n = (\dim\mathfrak{sl}_3)^n$,
consistent with the chain-group bound
(Corollary~\ref{cor:growth-rate-dimg}).
\item The recurrence gives:
 $\dim H^4 = 1352$, $\dim H^5 = 9892$, $\dim H^6 = 76084$,
 $\dim H^7 = 598592$
 (internal consistency of the recurrence~\eqref{eq:sl3-bar-recurrence}
 checked through degree~$7$; these values are \emph{generated} by the
 recurrence from the three proved initial values, not independently
 computed bar cohomology dimensions).
\item The characteristic polynomial of the recurrence factors as
$(t - 8)(t^2 - 3t - 1)$, where the linear factor corresponds to the
dominant pole $x = 1/\dim\mathfrak{sl}_3$ and the quadratic factor
$t^2 - 3t - 1$ is the reciprocal polynomial of the DS discriminant
factor $1 - 3x - x^2$. This factorization mirrors the $\mathfrak{sl}_2$
case: the corrected bar GF $x(3{-}x)/(1{-}x)^2$ has
a double pole at $x = 1$ (growth rate linear, not exponential),
while the DS-shared discriminant
$\Delta = 1 - 2x - 3x^2 = (1-3x)(1+x)$
still governs the Virasoro and $\beta\gamma$ generating functions
(Remark~\ref{rem:garland-lepowsky-sl2}).
\end{enumerate}

\emph{Lift obligation.} Equality with the bar cohomology of
$\widehat{\mathfrak{sl}}_{3,k}$ is the degree-$4$ rank identity for the
chiral bar differential. The coefficient $1352$ is the first
nontrivial value forced by recurrence~\eqref{eq:sl3-bar-recurrence}.

The degree-$4$ chain group has dimension
\[
\dim \bar{B}^4
=(\dim\mathfrak{sl}_3)^4\cdot 3!
=24576,
\]
and the chiral bracket differential
$d_4\colon\bar B^4\to\bar B^3$ maps to the
$1024$-dimensional group $\bar B^3$. The coefficient
$1352$ is equivalent to the homology identity
\[
\dim H^4
=
24576-\operatorname{rk}(d_4)-\operatorname{rk}(d_5)
=1352,
\]
where $d_5\colon\bar B^5\to\bar B^4$ and
$\dim\bar B^5=8^5\cdot4!=786432$. Thus the required chain-level
verification is
\[
\operatorname{rk}(d_4)+\operatorname{rk}(d_5)=23224,
\]
together with the lower-degree compatibility giving the proved
values $8,36,204$. Computation~\ref{comp:sl3-modular-rank} gives the
weight-decomposed rank framework for these matrices; it does not by
itself prove the displayed degree-$4$ homology identity.
\end{hypothesis}

\begin{remark}[Rank computation needed for the lift]
The finite-character datum becomes a bar-cohomology theorem exactly
when the chiral bar matrices realize
\[
\operatorname{rk}(d_4)+\operatorname{rk}(d_5)=23224.
\]
Equivalently, after the lower-degree ranks have produced
$H^1,H^2,H^3=(8,36,204)$, the incoming rank from
$\bar B^5$ must leave a $1352$-dimensional quotient in degree~$4$.
This is a concrete finite linear-algebra problem over~$\mathbb Q$,
with modular reductions over suitable primes as independent checks.
\end{remark}

\begin{remark}[Two discriminant families]\label{rem:two-discriminant-families}
\index{bar complex!discriminant families}
\index{Drinfeld--Sokolov reduction!discriminant families}
The proved rank-one formulas and the finite-character $A_2$ data give
two denominator families:
\begin{center}
\renewcommand{\arraystretch}{1.3}
\begin{tabular}{lll}
\toprule
\textbf{Chart} & \textbf{Denominator/discriminant}
 & \textbf{Members} \\
\midrule
$\mathfrak{sl}_2$ branch chart &
 $\Delta = (1-3x)(1+x)$ &
 $\widehat{\mathfrak{sl}}_2$ branch core, $\mathrm{Vir}$, $\beta\gamma$ \\
$\mathfrak{sl}_3$ family &
 $(1-3x-x^2)$ &
 $\widehat{\mathfrak{sl}}_3$, $\mathcal{W}_3$ \\
\bottomrule
\end{tabular}
\end{center}
The Virasoro and $\beta\gamma$ functions have algebraic degree~$2$
and branch locus $\{\tfrac{1}{3}, -1\}$; the corrected affine
$\widehat{\mathfrak{sl}}_2$ coefficient series is rational with
denominator $(1-x)^2$ and is connected to the same chart through the
rank-one branch-divisor core.  The $A_2$ recurrence data have rational
series sharing the denominator factor $(1-3x-x^2)$ with roots at
$(-3 \pm \sqrt{13})/2$.
In each finite-character chart, DS reduction changes the growth pole
and preserves the displayed divisor-core factor.
For example, $\widehat{\mathfrak{sl}}_3$ has a
growth pole at $x = 1/8$ while $\mathcal{W}_3$ has a
growth pole at $x = (-3+\sqrt{13})/2 \approx 0.303$.
\end{remark}

\begin{theorem}[Bar cohomology generating-function chart; \ClaimStatusConditional]
\label{thm:bar-gf-classification}
\index{bar complex!generating function!classification|textbf}
\index{D-finiteness!bar cohomology dichotomy}
\index{generating function!bar cohomology!landscape}
The proved bar cohomology generating functions and the finite-character
recurrence data for
$P_\cA(x) = \sum_{n \geq 1} \dim H^n(\barBgeom(\cA))\, x^n$
are recorded in Table~\textup{\ref{tab:bar-gf-classification}}.
Three structural features emerge:
\begin{enumerate}[label=\textup{(\roman*)}]
\item \emph{D-finiteness dichotomy.}
 Every proved interacting row and every finite-character row
 \textup{(}nontrivial chiral bracket $a_{(0)}b \neq 0$ for some
 generators\textup{)} has a D-finite generating function
 \textup{(}rational or algebraic\textup{)}.
 The unique non-D-finite entry is the Heisenberg algebra,
 whose bar cohomology is governed by the partition function
 $\prod_{n \geq 1}(1 - q^n)^{-1}$, which is not D-finite
 \textup{(}its coefficients $p(n) \sim
 e^{\pi\sqrt{2n/3}}/(4\sqrt{3}\,n)$ violate the polynomial
 coefficient recurrence required of D-finite
 sequences\textup{)}.
\item \emph{Rational vs.\ algebraic.}
 The proved $bc$ and $\widehat{\mathfrak{sl}}_2$ rows are rational;
 the Virasoro and $\beta\gamma$ rows are algebraic of degree~$2$.
 The $\widehat{\mathfrak{sl}}_3$ and $\mathcal{W}_3$ rows are rational
 after imposing the displayed finite-character recurrences.
\item \emph{DS-invariant discriminant.}
 The denominator\slash branch-locus polynomials organize into
 DS orbits
 \textup{(}Remark~\textup{\ref{rem:two-discriminant-families})}\textup{:}
 the rank-one family has branch divisor $(1-3x)(1+x)$ on the
 Virasoro and $\beta\gamma$ side; the $A_2$ finite-character chart
 shares the irreducible divisor-core factor $1 - 3x - x^2$.
\end{enumerate}
\end{theorem}

\begin{table}[ht]
\centering
\caption{Bar cohomology generating function classification}
\label{tab:bar-gf-classification}
\index{bar complex!generating function!classification table}
\renewcommand{\arraystretch}{1.5}
{\small
\resizebox{\textwidth}{!}{%
\begin{tabular}{|l|c|c|c|c|c|}
\hline
\textbf{Algebra $\cA$}
 & \textbf{Generating function $P_\cA(x)$}
 & \textbf{Growth}
 & \textbf{Type}
 & \textbf{D-finite}
	 & \textbf{Source} \\
\hline
\multicolumn{6}{|c|}{\textit{$\mathfrak{sl}_2$ rank-one branch chart:
 $\Delta = (1-3x)(1+x)$}} \\
\hline
$\widehat{\mathfrak{sl}}_2$
 & $\dfrac{x(3-x)}{(1-x)^2}$
 & $2n{+}1$
 & rational
 & \checkmark
	 & bar complex \\[6pt]
\hline
$\mathrm{Vir}_c$
 & $x \cdot M(x)^2$\textsuperscript{$\dagger$}
 & $\sim 3^n / n^{3/2}$
 & alg.\ deg $2$
 & \checkmark
	 & bar complex \\[4pt]
\hline
$\beta\gamma$
 & $\sqrt{(1+x)/(1-3x)}$
 & $\sim 3^n / n^{1/2}$
 & alg.\ deg $2$
 & \checkmark
	 & bar complex \\[4pt]
\hline
\multicolumn{6}{|c|}{\textit{Fermionic: $bc$ ghosts}} \\
\hline
$bc$
 & $\dfrac{x(2-5x+4x^2)}{(1-2x)(1-x)^2}$
 & $\sim 2^n$
 & rational
 & \checkmark
	 & bar complex \\[6pt]
\hline
\multicolumn{6}{|c|}{\textit{$\mathfrak{sl}_3$ discriminant family:
 $\Delta^{\mathrm{inv}} = 1 - 3x - x^2$}} \\
\hline
$\widehat{\mathfrak{sl}}_3$
 & $\dfrac{4x(2-13x-2x^2)}{(1-8x)(1-3x-x^2)}$
 & $\sim 8^n$
 & rational
 & \checkmark
	 & recurrence datum\\[6pt]
\hline
$\mathcal{W}_3$
 & $\dfrac{x(2-3x)}{(1-x)(1-3x-x^2)}$
 & $\sim \bigl(\!\frac{3{+}\sqrt{13}}{2}\!\bigr)^n$
 & rational
 & \checkmark
	 & recurrence datum\\[6pt]
\hline
\multicolumn{6}{|c|}{\textit{Free fields: non-D-finite}} \\
\hline
Heisenberg $\cH_\kappa$
 & $\displaystyle\prod_{n \geq 1}(1 - q^n)^{-1}$
 & $e^{\pi\sqrt{2n/3}}$
 & not alg.
 & $\boldsymbol{\times}$
	 & bar complex \\[6pt]
\hline
\end{tabular}%
}% end resizebox
}% end small
\end{table}

\noindent
{\footnotesize $^{\dagger}$\,$M(x) = (1 - x - \sqrt{1-2x-3x^2})/(2x^2)$
is the Motzkin generating function (OEIS~A001006);
$P_{\mathrm{Vir}} = x \cdot M(x)^2 = 4x/(1-x+\sqrt{1-2x-3x^2})^2$.
The bar cohomology dimensions are the Motzkin first differences
$\dim H^n = M(n{+}1) - M(n)$:
$1, 2, 5, 12, 30, 76, 196, 512, \ldots$
(Theorem~\ref{thm:ds-bar-gf-discriminant}).}

\begin{proof}
The proved generating functions are established in
Theorem~\ref{thm:ds-bar-gf-discriminant}
(Virasoro, $\beta\gamma$, $\widehat{\mathfrak{sl}}_2$),
Remark~\ref{rem:garland-lepowsky-sl2}
($\widehat{\mathfrak{sl}}_2$ correcting Riordan),
and Proposition~\ref{prop:bar-bc-system}
with Corollary~\ref{cor:bar-cohomology-koszul-dual}
($bc$ ghosts: $\dim(bc^!)_n = 2^n - n + 1$, with rational
generating function
$P_{bc}(x) = x(2-5x+4x^2)/((1-2x)(1-x)^2)$ and dominant pole
at $x = 1/2$ giving growth rate~$2^n$).
The $\widehat{\mathfrak{sl}}_3$ and $\mathcal{W}_3$ rows are the
finite-character recurrence data
\eqref{eq:sl3-bar-recurrence} and~\eqref{eq:w3-bar-recurrence};
their equality with bar cohomology is the corresponding rank/lift
obligation.
The Heisenberg entry follows from
Theorem~\ref{thm:heisenberg-bar}: the bar cohomology of
$\cH_\kappa$ is the exterior algebra on $V^*$, giving
$\dim (\cH_\kappa^!)_n = p(n{-}2)$ (partitions);
the generating function $\prod(1-q^n)^{-1}$ is transcendental
(Nesterenko 1996: $e^{\pi}$ and $\Gamma(1/4)$ are
algebraically independent, forcing the $q$-series to be
non-algebraic) and not D-finite
(a theorem of Flajolet: the
Hardy--Ramanujan asymptotics $p(n) \sim
e^{\pi\sqrt{2n/3}}/(4\sqrt{3}\,n)$ violate the
Birkhoff--Trjitzinsky condition for D-finite sequences,
which requires at most polynomial-times-exponential growth).

The D-finiteness dichotomy follows:
rational functions are D-finite (trivially),
algebraic functions are D-finite
(by the closure theorem of Zeilberger),
and the partition function is the unique non-D-finite entry.
The discriminant family structure is
Remark~\ref{rem:two-discriminant-families} and
Theorem~\ref{thm:ds-spectral-branch-preservation}.
\end{proof}

\begin{remark}[The D-finiteness dichotomy and OPE structure]
\label{rem:d-finiteness-dichotomy}
\index{D-finiteness!OPE structure}
\index{free fields!non-D-finite bar cohomology}
The D-finiteness dichotomy in
Theorem~\ref{thm:bar-gf-classification} has a structural
explanation.
For an interacting algebra, the chiral bracket
$a_{(0)}b \neq 0$ forces the PBW spectral sequence
(Theorem~\ref{thm:pbw-koszulness-criterion}) to have a
nontrivial $d_1$ differential, which kills most of the
$E_1$~page and leaves a D-finite residue.
For the Heisenberg algebra, $a_{(0)}b = 0$ for all generators:
the bar differential acts only through the central extension
$a_{(1)}b = \kappa\,\delta_{a+b,0}$, and the surviving bar
cohomology is the full exterior algebra on infinitely many
generators (one per negative mode), whose Hilbert function
is the partition function.

The intermediate case of $\widehat{\mathfrak{sl}}_2$ is
instructive: the chiral bracket is nontrivial
($E_{(0)}F = H$, etc.), and the Garland--Lepowsky
concentration (Remark~\ref{rem:garland-lepowsky-sl2}) shows
that the semisimple $d_1$ differential is maximally
efficient, killing all but $2n{+}1$ cocycles at degree~$n$.
The resulting linear growth $\dim H^n = 2n+1$ corrects the
Riordan formula $\dim H^n = R(n{+}3)$, which incorrectly
treats the $E_1$~page as the bar cohomology.

The exact structural implication left by this chart is the construction
of a finite-character recurrence module from a nontrivial chiral bracket
among strong generators.
\end{remark}

\begin{proposition}[Discriminant as a finite characteristic invariant; \ClaimStatusConditional]
\label{prop:discriminant-characteristic}
\index{discriminant!characteristic invariant}
\index{modular characteristic hierarchy!discriminant}
For a modular Koszul chiral algebra $\cA$
\textup{(}Definition~\textup{\ref{def:modular-koszul-chiral})},
the Koszul dual discriminant $\Delta_{\cA}(x)$ is a finite
non-scalar invariant in the modular characteristic hierarchy whenever
the Hilbert series is algebraic or supplied by a finite-character
recurrence datum.
Specifically:
\begin{enumerate}[label=\textup{(\roman*)}]
\item If the bar-dual Hilbert series
 $P_{\cA}(x) = \sum_{n \geq 0} \dim(\cA^{\mathrm i})_n\, x^n$
 \textup{(}equivalently the $\cA^!$ Hilbert series after the
 finite/completed dual branch\textup{)}
 is algebraic of degree~$d$, then $\Delta_{\cA}(x)$ is a polynomial
 of degree at most $d(d-1)$ whose roots are the branch points
 of $P_{\cA}$.
\item The obstruction coefficient $\kappa(\cA)$ and
 the discriminant $\Delta_{\cA}(x)$ are related by: the growth
 rate $\rho = 1/\min\{|r| : \Delta_{\cA}(r) = 0\}$ satisfies
 $\log \rho \leq \log(\dim \cA_{[1]})$, with equality for
 Kac--Moody algebras.
\item Under a finite-character Drinfeld--Sokolov denominator chart
 $\cA = \widehat{\fg}_k \mapsto \mathcal{W}^k(\fg)$, the maximal
 common divisor core of the two shift polynomials is preserved
 \textup{(}Theorem~\textup{\ref{thm:ds-spectral-branch-preservation})}.
\item Under Koszul duality $\cA \mapsto \cA^!$, equality of
 discriminants is a statement about the induced finite-character
 recurrence module.  Verdier duality preserves the underlying
 divisor-core object once that module is constructed.
\end{enumerate}
\end{proposition}

\begin{proof}
Part (i) is the definition of the discriminant of an algebraic function.
Part (ii): the radius of convergence of $P_{\cA}(x)$ is the
reciprocal of the exponential growth rate $\rho$; for Kac--Moody,
$\dim (\widehat{\fg}_k^!)_n \leq (\dim \fg)^n \cdot (n-1)!$
(Corollary~\ref{cor:growth-rate-dimg}), so $\rho \leq \dim \fg$;
equality follows from the explicit chain-group count.
Part (iii) is Theorem~\ref{thm:ds-spectral-branch-preservation}.
Part (iv): the bar complex of $\cA^!$ is related to that of $\cA$
by Verdier duality
(Theorem~\ref{thm:bar-cobar-isomorphism-main}); the determinant claim
is therefore made on the finite cyclic divisor-core module, where
Verdier duality is an involution.  For the corrected affine
$\widehat{\mathfrak{sl}}_2$ coefficient series, the dual level has the
same rational denominator $(1-x)^2$; the separate rank-one branch
divisor is $(1-3x)(1+x)$.
\end{proof}

\begin{remark}[Discriminant and the universal class]\label{rem:discriminant-universal-class}
\index{discriminant!universal Maurer--Cartan class}
Proposition~\ref{prop:discriminant-characteristic} and
Hypothesis~\ref{hyp:discriminant-ks-operator} gives the finite-rank
interpretation of the discriminant $\Delta_{\cA}(x)$ as a shadow of the universal
Maurer--Cartan class $\Theta_{\cA}$
(\S\ref{sec:modular-koszul-theory}). The precise relationship
is: if $\Theta_{\cA}$ has a genus-$1$ component
$\Theta_{\cA}^{(1)} \in \operatorname{Def}_{\mathrm{cyc}}(\cA)
\otimes H^*(\overline{\mathcal{M}}_{1,1})$, then the reduced
genus-$1$ transport operator $\operatorname{KS}_{\cA}$ acts on
a finite-rank piece $H^{\mathrm{red}}_1(\cA)$, and the
discriminant is its characteristic polynomial:
$\Delta_{\cA}(x) = \det(1 - x \cdot \operatorname{KS}_{\cA}
\mid H^{\mathrm{red}}_1(\cA))$. For the $\mathfrak{sl}_2$ family
($\operatorname{rank} = 2$), the eigenvalues $3$ and $-1$ of
$\operatorname{KS}_{\cA}$ produce
$\Delta = (1 - 3x)(1 + x)$. This interpretation connects the
bar cohomology generating function to the genus-$1$ piece of
the proved full characteristic hierarchy, whose H-level apex is
the characteristic object~$\mathcal{C}_{\cA}$
(Definition~\ref{def:full-modular-package};
Remark~\ref{rem:modular-characteristic-package}).  In a DS denominator
chart, the preserved object is the common divisor-core sector of the
Kodaira--Spencer transport operator.
\end{remark}

\begin{remark}[DS transformation of bar generating functions]\label{rem:ds-gf-transformation}
\index{Drinfeld--Sokolov reduction!generating function transformation}
Theorem~\ref{thm:ds-bar-gf-discriminant} separates the corrected
rank-one affine coefficient series from the branch-divisor chart.  The
actual affine series is
\[
P_{\widehat{\mathfrak{sl}}_2}(x)=\frac{x(3-x)}{(1-x)^2}.
\]
The Virasoro series is
\[
P_{\mathrm{Vir}}(x) = \frac{4x}{(1 - x + \sqrt{\Delta})^2},
\qquad \Delta = 1 - 2x - 3x^2.
\]
Thus the rank-one DS bridge is a lift to the branch-divisor core
$\Delta=(1-3x)(1+x)$, not an equality of coefficient-recurrence
denominators.
\end{remark}

\begin{theorem}[Linear dependence in the rank-one branch family; \ClaimStatusProvedHere]
\label{thm:discriminant-linear-dependence}
\index{bar complex!generating function!linear dependence}
\index{Drinfeld--Sokolov reduction!linear dependence}
\index{Wakimoto representation!discriminant relation}
Let
\[
R_{\mathfrak{sl}_2}(x)
 := \frac{1 + x - \sqrt{1 - 2x - 3x^2}}{2x(1+x)}
\]
be the rank-one branch function.  It is the Riordan function, not the
corrected affine $\widehat{\mathfrak{sl}}_2$ bar series.  The three
rank-one branch functions satisfy the $\mathbb{Q}(x)$-linear relation
\begin{equation}\label{eq:discriminant-linear-relation}
x^2(1{+}x)\, P_{\mathrm{Vir}}(x)
= (1{-}2x{-}x^2)(1{+}x)\, R_{\mathfrak{sl}_2}(x)
 - (1{-}3x)\, P_{\beta\gamma}(x).
\end{equation}
Equivalently, with $\Delta = 1 - 2x - 3x^2$ and $u = \sqrt{\Delta}$:
\begin{equation}\label{eq:discriminant-family-decomposition}
R_{\mathfrak{sl}_2} = \frac{1}{2x} - \frac{u}{2x(1{+}x)},
\qquad
P_{\beta\gamma} = \frac{u}{1{-}3x},
\qquad
P_{\mathrm{Vir}} = \frac{1{-}2x{-}x^2}{2x^3} - \frac{(1{-}x)\,u}{2x^3}.
\end{equation}
\end{theorem}

\begin{proof}
The three displayed branch functions
lie in the quadratic extension
$\mathbb{Q}(x)[\sqrt{\Delta}]$ of $\mathbb{Q}(x)$,
which is a $2$-dimensional $\mathbb{Q}(x)$-vector space with basis
$\{1, u\}$ where $u = \sqrt{\Delta}$ and $u^2 = \Delta = (1{-}3x)(1{+}x)$.

\emph{Expansion in the basis $\{1, u\}$.}
From the displayed formula for $R_{\mathfrak{sl}_2}$:
$R_{\mathfrak{sl}_2} = (1{+}x{-}u)/(2x(1{+}x)) = 1/(2x) - u/(2x(1{+}x))$.
For $P_{\beta\gamma}$, rationalize $(1{+}x)/u = (1{+}x)u/u^2 = (1{+}x)u/\Delta = u/(1{-}3x)$.
For $P_{\mathrm{Vir}} = 4x/(1{-}x{+}u)^2$,
multiply numerator and denominator by the conjugate
$(1{-}2x{-}x^2) - (1{-}x)u$:
\[
(1{-}x{+}u)^2 = 2(1{-}2x{-}x^2) + 2(1{-}x)u,
\]
so $P_{\mathrm{Vir}} = 2x/[(1{-}2x{-}x^2)+(1{-}x)u]$.
Rationalizing gives denominator
$(1{-}2x{-}x^2)^2 - (1{-}x)^2\Delta = 4x^4$, hence
\[
P_{\mathrm{Vir}} = \frac{(1{-}2x{-}x^2)-(1{-}x)u}{2x^3}.
\]
This establishes~\eqref{eq:discriminant-family-decomposition}.

\emph{The linear relation.}
Since $\dim_{\mathbb{Q}(x)} \mathbb{Q}(x)[u] = 2$,
any three elements are $\mathbb{Q}(x)$-linearly dependent.
Solving $P_{\mathrm{Vir}} = \alpha\, R_{\mathfrak{sl}_2} + \beta\, P_{\beta\gamma}$
from the $\{1,u\}$-components:
\begin{align*}
\text{constant:} \quad
 \frac{1{-}2x{-}x^2}{2x^3} &= \frac{\alpha}{2x},
 &&\Rightarrow\ \alpha = \frac{1{-}2x{-}x^2}{x^2}, \\
\text{$u$-part:} \quad
 \frac{-(1{-}x)}{2x^3} &= \frac{-\alpha}{2x(1{+}x)} + \frac{\beta}{1{-}3x},
 &&\Rightarrow\ \beta = \frac{-(1{-}3x)}{x^2(1{+}x)}.
\end{align*}
Multiplying $P_{\mathrm{Vir}} = \alpha\, R_{\mathfrak{sl}_2} + \beta\, P_{\beta\gamma}$
through by $x^2(1{+}x)$ gives~\eqref{eq:discriminant-linear-relation}.
\end{proof}

\begin{remark}[Consequences of branch linear dependence]\label{rem:linear-dependence-consequences}
\index{bar complex!generating function!mutual determination}
Theorem~\ref{thm:discriminant-linear-dependence} has three consequences
inside the rank-one branch chart:
\begin{enumerate}
\item \emph{Mutual determination.}
Any two branch functions determine the third via~\eqref{eq:discriminant-linear-relation}. In particular, the $\beta\gamma$ generating function can be recovered from the Riordan branch function and Motzkin family:
\[
P_{\beta\gamma} = \frac{(1{-}2x{-}x^2)(1{+}x)\, R_{\mathfrak{sl}_2}
 - x^2(1{+}x)\, P_{\mathrm{Vir}}}{1{-}3x}.
\]
\item \emph{Algebraic structure.}
The rank-one branch family forms a $1$-dimensional
$\mathbb{Q}(x)$-affine subspace of $\mathbb{Q}(x)[u]$:
every member has the form $a(x) + b(x)u$ where
$a, b$ are constrained by the DS and Wakimoto functors.
The two ``canonical'' basis vectors are
$R_{\mathfrak{sl}_2}$ (with negative $u$-coefficient)
and $P_{\beta\gamma}$ (with positive $u$-coefficient, reflecting the
$\Delta^{-1/2}$ vs $\Delta^{1/2}$ singularity types).
\item \emph{Singularity exchange.}
The relation~\eqref{eq:discriminant-linear-relation} interchanges
the sub-exponential correction exponents:
the Riordan branch function and $P_{\mathrm{Vir}}$ have corrections
$n^{-3/2}$ (from $u = \Delta^{1/2}$),
while $P_{\beta\gamma}$ has correction $n^{-1/2}$
(from $u^{-1} = \Delta^{-1/2}$, cf.\ Corollary~\ref{cor:betagamma-inverse-discriminant}).
The linear relation~\eqref{eq:discriminant-linear-relation}
mixes these two singularity types through the
$\mathbb{Q}(x)$-coefficients.
\end{enumerate}
\end{remark}

\begin{proposition}[Obstruction to a coefficient-level relation; \ClaimStatusConditional]
\label{prop:linear-relation-functorial}
\index{Wakimoto representation!linear relation}
\index{Drinfeld--Sokolov reduction!linear relation}
The $\mathbb{Q}(x)$-linear
relation~\textup{\eqref{eq:discriminant-linear-relation}} is a relation
in the rank-one branch-divisor chart.  It is not a relation among the
three corrected coefficient series, because
$P_{\widehat{\mathfrak{sl}}_2}(x)=x(3-x)/(1-x)^2$ lies in
$\mathbb{Q}(x)$ and has shift polynomial $(t-1)^2$.  A functorial
coefficient-level statement requires a chain map from the affine
coefficient recurrence module to the branch-divisor module
$k[t]/((t-3)(t+1))$ compatible with the diagram
\begin{equation}\label{eq:wakimoto-ds-triangle}
\begin{tikzcd}
\bar{B}(\widehat{\mathfrak{sl}}_{2,k})
 \ar[r, "\mathrm{DS}"]
 \ar[d, hook, "\mathrm{Wak}"']
& \bar{B}(\mathrm{Vir}_c) \\
\bar{B}(\beta\gamma) \otimes \bar{B}(\mathcal{H})
 \ar[ur, dashed]
\end{tikzcd}
\end{equation}
arising from the Wakimoto embedding
$\widehat{\mathfrak{sl}}_{2,k} \hookrightarrow
\beta\gamma \otimes \mathcal{H}$
\textup{(}Theorem~\textup{\ref{thm:wakimoto-bar})} and the
Drinfeld--Sokolov reduction
$\mathrm{Vir}_c = H^0_{\mathrm{DS}}(
\widehat{\mathfrak{sl}}_{2,k})$
\textup{(}Theorem~\textup{\ref{thm:ds-koszul-intertwine})}.

The two structural maps therefore impose the following intertwining data:
\begin{enumerate}[label=\textup{(\roman*)}]
\item \emph{DS projection}: a BRST-compatible lift from the affine
 coefficient module to the Virasoro branch module.
\item \emph{Wakimoto embedding}: relates
 the affine coefficient module to the $\beta\gamma$ branch module via
 the factorization
 $B(\widehat{\mathfrak{sl}}_2)
 \hookrightarrow B(\beta\gamma) \otimes B(\mathcal{H})$.
 The Heisenberg factor contributes only sub-exponential growth
 ($\dim (\mathcal{H}^!)_n = p(n{-}2)$), so the exponential
 behavior and discriminant are determined by the
 $\beta\gamma$ factor.
\end{enumerate}
The $\mathbb{Q}(x)$-coefficients in~\textup{\eqref{eq:discriminant-linear-relation}}
encode the branch-divisor chart after this bridge is supplied.
\end{proposition}

\begin{proof}
The corrected affine series is rational with denominator $(1-x)^2$,
whereas the branch function $R_{\mathfrak{sl}_2}$ has branch divisor
$(1-3x)(1+x)$.  Hence the relation
\eqref{eq:discriminant-linear-relation} cannot be obtained by simply
substituting the corrected affine bar series into the rank-one branch
calculation.  The displayed diagram supplies the only possible source
of a coefficient-level theorem: a compatible lift between the two
finite cyclic modules.  Once such a lift is given, the linear relation
is exactly the $\{1,u\}$-decomposition in the proof of
Theorem~\ref{thm:discriminant-linear-dependence}.
\end{proof}

\begin{proposition}[Pole decomposition and singularity type; \ClaimStatusConditional]
\label{prop:pole-singularity-type}
\index{bar complex!generating function!singularity type}
\index{pole decomposition!sub-polynomial correction}
Within the rank-one branch family
(Theorem~\textup{\ref{thm:ds-bar-gf-discriminant}}), the sub-polynomial
correction to exponential growth is determined by the pole decomposition
of the bar differential
(Proposition~\textup{\ref{prop:pole-decomposition}}):
\begin{enumerate}
\item If $d_{\mathrm{bracket}} \neq 0$ (bracket channel active),
 the generating function has a \emph{square-root branch point}
 at $x = 1/3$ and the coefficients grow as
 $\dim (\cA^!)_n \sim C \cdot 3^n \cdot n^{-3/2}$.
 This applies to the Riordan branch function
 $R_{\mathfrak{sl}_2}$ and to Virasoro.  The corrected affine
 $\widehat{\mathfrak{sl}}_2$ coefficient series has a double pole at
 $x=1$ and linear growth.
\item If $d_{\mathrm{bracket}} = 0$ (curvature-only), the generating function
 has an \emph{inverse square-root singularity} at $x = 1/3$ and
 the coefficients grow as
 $\dim (\cA^!)_n \sim C' \cdot 3^n \cdot n^{-1/2}$.
 This applies to $\beta\gamma$.
\end{enumerate}
\end{proposition}

\begin{proof}
From the $\{1, u\}$-decomposition~\eqref{eq:discriminant-family-decomposition}
with $u = \sqrt{\Delta}$, $\Delta = (1{-}3x)(1{+}x)$:

For the Riordan branch function:
$R_{\mathfrak{sl}_2} = 1/(2x) - u/(2x(1{+}x))$.
The singular part near $x = 1/3$ is the $u$-term. Since
$u = \sqrt{(1{-}3x)(1{+}x)} \approx c_0\sqrt{1{-}3x}$ as $x \to 1/3^-$,
the transfer theorem (Flajolet--Sedgewick, Theorem~VI.3) gives
$[x^n] u/(2x(1{+}x)) \sim C \cdot 3^n \cdot n^{-3/2}$.
The $\beta\gamma$ OPE has a simple pole only
($\beta_{(0)}\gamma = 1 \in \mathbb{C}$, a scalar),
so $d_{\mathrm{bracket}} = 0$
(Proposition~\ref{prop:pole-decomposition}: the simple-pole residue is
the vacuum scalar rather than a field value).
Its generating function $P_{\beta\gamma} = u/(1{-}3x) = (1{+}x)/u$
has the \emph{inverse} singularity $u^{-1} \sim c_1 / \sqrt{1{-}3x}$
near $x = 1/3^-$, giving $[x^n] \sim C' \cdot 3^n \cdot n^{-1/2}$
by the same transfer theorem.

The Virasoro case: $P_{\mathrm{Vir}} = ((1{-}2x{-}x^2) - (1{-}x)u)/(2x^3)$
has the same $u$-type singularity as $R_{\mathfrak{sl}_2}$
(the coefficient of $u$ is a polynomial with $u \sim \sqrt{1{-}3x}$),
hence the same $n^{-3/2}$ correction.
The Virasoro OPE \emph{does} have a bracket channel
($T_{(0)}T = \partial T$, $T_{(1)}T = 2T$),
so $d_{\mathrm{bracket}} \neq 0$, consistent with case (1).
\end{proof}

\begin{remark}[Bar dimensions as Koszul dual Hilbert series]\label{rem:bar-hilbert-koszul-dual}
\index{bar complex!Koszul dual Hilbert series}
For a Koszul algebra $\cA$, the conformal-weight-graded bar cohomology
dimension at bar degree~$n$ is
$\dim H^\bullet(\barBgeom(\cA))_n$.  For a quadratic presentation
$\cA=T_X(V)/(R)$, the coalgebra
$\cA^{\mathrm i}=C_X(s^{-1}V,s^{-2}R)$ carries
$q_\cA\colon\cA^{\mathrm i}\to\barBgeom(\cA)$; on the Koszul locus
$H^\bullet(q_\cA)$ identifies its cohomology with bar cohomology.
Under the finite-type or completed Verdier-duality hypotheses of
Corollary~\ref{cor:bar-cohomology-koszul-dual}, the pair map
$\nu_\cA\colon K_X(\cA)\to\cA^!$ identifies this number with the
Koszul-dual Hilbert function $\dim(\cA^!)_n$.  Thus
$P_{\cA}(x)=\sum_n\dim H^\bullet(\barBgeom(\cA))_n x^n$ is also the
Hilbert series of $\cA^!$ on that comparison surface. In the Kac--Moody case,
$P_{\widehat{\mathfrak{g}}_k}(x)$ coincides with the Hilbert series of
$\mathrm{CE}^{\mathrm{ch}}(\widehat{\mathfrak{g}}_{-k-2h^\vee})$ as a
graded algebra; this is independent of~$k$ because the Koszul dual branch
has the same graded structure at all generic levels.
\end{remark}

\begin{lemma}[Degree-\texorpdfstring{$2$}{2} bar cohomology at lowest weight; \ClaimStatusProvedHere]\label{lem:bar-deg2-symmetric-square}
\index{bar complex!degree-2 cohomology}
At conformal weight $h = 2$, the bar complex of
$\widehat{\mathfrak{g}}_k$ (semisimple~$\mathfrak{g}$) satisfies
$H^2(\barBgeom(\widehat{\mathfrak{g}}_k))_{h=2} = 0$.
The chain groups are $\bar{B}^2_{h=2} \cong \mathfrak{g}^{\otimes 2}$
(dimension~$(\dim\mathfrak{g})^2$) and
$\bar{B}^1_{h=2} \cong \bar{A}_2 = \mathfrak{g}_2 \oplus S^2(\mathfrak{g}_1)$
(dimension $\dim\mathfrak{g} + \binom{\dim\mathfrak{g}+1}{2}$),
where $\mathfrak{g}_n$ denotes the mode-$n$ copy of~$\mathfrak{g}$
and $S^2(\mathfrak{g}_1)$ is spanned by the normally ordered products
$\normord{J^a J^b} = J^a_{-1}J^b_{-1}|0\rangle$. For
$\mathfrak{sl}_2$: both $\bar{B}^2_{h=2}$ and $\bar{B}^1_{h=2}$
have dimension~$9$.

The bar differential $d \colon \bar{B}^2_{h=2} \to \bar{B}^1$
has two components, corresponding to the two singular terms in the
OPE $J^a(z)J^b(w) \sim k\delta^{ab}/(z-w)^2 + f^{ab}{}_c
J^c(w)/(z-w)$:
\begin{itemize}
\item $d_{(-1)} \colon \bar{B}^2_{h=2} \to \bar{B}^1_{h=2}$
 via the $(-1)$-st product (normal ordering):
 $J^a \otimes J^b \mapsto J^a_{-1}J^b_{-1}|0\rangle$;

\item $d_{(0)} \colon \bar{B}^2_{h=2} \to \bar{B}^1_{h=1}$
 via the $(0)$-th product (Lie bracket):
 $J^a \otimes J^b \mapsto f^{ab}{}_c\, J^c$.
\end{itemize}

An element $\sum c_{ab}\, J^a \otimes J^b$ lies in $\ker(d)$ iff
both components vanish. The bracket condition
$\sum c_{ab}\, f^{ab}{}_c = 0$ for all~$c$ forces $c_{ab}$
to be \emph{symmetric} (since $f^{ab}{}_c$ is antisymmetric in
$a,b$). The normal-ordering condition then requires
$\sum_{a \leq b} c'_{ab}\, \normord{J^a J^b} = 0$
in~$\bar{A}_2$; but $\{\normord{J^a J^b}\}_{a \leq b}$ are
linearly independent PBW basis elements of
$S^2(\mathfrak{g}_1)$, so $c'_{ab} = 0$.
Hence $\ker(d) = 0$. Since $\bar{B}^3_{h=2} = 0$ (three
mode-$1$ generators have weight~$3$), one obtains $H^2_{h=2} = 0$.
For $\mathfrak{sl}_2$, the normal-ordering map alone is an
isomorphism ($9 \to 9$), but for $\mathfrak{g}$ of rank
${\geq}\, 2$ both components of~$d$ are needed.

There is therefore no discrepancy with the PBW spectral sequence
(Computation~\ref{comp:sl2-ce-verification}):
$H^2\bigl(\barBgeom(\widehat{\mathfrak{sl}}_2)\bigr) = 5$,
with all five classes at weight~${\geq}\, 3$.
\end{lemma}

\begin{corollary}[Exponential growth rate in a finite-character Kac--Moody chart; \ClaimStatusProvedHere]
\label{cor:growth-rate-dimg}
\index{bar complex!growth rate!Lie algebra dimension}
Let a finite-character recurrence datum for
$\widehat{\mathfrak{g}}_k$ have denominator
$\Delta_{\mathfrak{g}}(x)$ with a simple factor
$1-\dim(\mathfrak{g})x$ and no root of smaller modulus.  Its
coefficients $a_n$ satisfy
\[
\lim_{n \to \infty} \frac{\log |a_n|}{n} = \log(\dim \mathfrak{g}).
\]
The corrected affine $\widehat{\mathfrak{sl}}_2$ coefficient series is
outside this conclusion: it has denominator $(1-x)^2$ and linear
growth.  The rank-one branch chart has spectral radius~$3$, and the
$A_2$ finite-character datum has spectral radius~$8$.
\end{corollary}

\begin{proof}
For a rational or algebraic finite-character series, the exponential
growth rate is the reciprocal of the smallest-modulus singularity.  The
assumption makes $x=1/\dim(\mathfrak{g})$ the nearest singularity, hence
the displayed limit.
\end{proof}

\begin{theorem}[Dominant branch point in a finite-character Kac--Moody chart; \ClaimStatusProvedHere]
\label{thm:dominant-branch-point}
\index{bar complex!dominant branch point}
\index{discriminant!dominant eigenvalue}
Let $\mathfrak{g}$ be a simple Lie algebra and suppose a
finite-character denominator chart for $\widehat{\mathfrak{g}}_k$ has
algebraic or rational series $P(x)$ with denominator/discriminant
$\Delta(x)$ and nearest singularity $R=1/\dim\mathfrak{g}$. Then:
\begin{enumerate}[label=\textup{(\roman*)}]
\item The radius of convergence of $P(x)$ is
 $R = 1/\dim\mathfrak{g}$.
\item The point $x = R$ is a pole or branch point of~$P$, hence a root
 of the denominator/discriminant~$\Delta(x)$.
\item $\dim\mathfrak{g}$ is the spectral radius of the branch
 transfer matrix~$T_{\mathrm{br}}$ defined by
 $\Delta(x)/\Delta(0) = \det(1 - x\,T_{\mathrm{br}})$.
\end{enumerate}
\end{theorem}

\begin{proof}
\emph{(i).}
By Corollary~\ref{cor:growth-rate-dimg}, the finite-character
coefficient growth rate is $\dim\mathfrak{g}$, so
$R = 1/\dim\mathfrak{g}$ by the
Cauchy--Hadamard formula.

\emph{(ii).}
The denominator/discriminant chart records all finite singularities of
the finite-character series.  Since $R$ is the nearest singularity by
hypothesis, $R=1/\dim\mathfrak{g}$ is a root of~$\Delta(x)$.

\emph{(iii).}
The eigenvalues of $T_{\mathrm{br}}$ are the reciprocals
$\alpha_i = 1/x_i$ of the roots~$x_i$ of~$\Delta$.
The spectral radius $\max_i |\alpha_i| = 1/\min_i |x_i| = 1/R
= \dim\mathfrak{g}$, since the nearest branch point to the origin
determines the radius of convergence. Part~(ii) guarantees that
this nearest branch point is $x = R$.
\end{proof}

\begin{remark}[Computed cases]\label{rem:dominant-branch-verified}
For the rank-one branch chart: $R = 1/3$,
$\Delta = (1 - 3x)(1+x)$,
$T_{\mathrm{br}} = \operatorname{diag}(3, -1)$, spectral radius $3
= \dim\mathfrak{sl}_2$.
For the $\mathfrak{sl}_3$ finite-character generating function
(Hypothesis~\ref{hyp:sl3-finite-character-recurrence}): $R = 1/8$,
$\Delta = (1 - 8x)(1 - 3x - x^2)$, spectral radius $8
= \dim\mathfrak{sl}_3$. The sub-dominant eigenvalues
$(3 \pm \sqrt{13})/2$ are DS-invariant
(Remark~\ref{rem:ds-discriminant-invariant}).
\end{remark}

\begin{remark}[DS dimension inequality]\label{rem:ds-dimension-inequality}
\index{Drinfeld--Sokolov reduction!dimension inequality}
For any simple $\mathfrak{g}$, DS reduction cannot increase Koszul dual dimensions:
\[
\dim (\mathcal{W}^k(\mathfrak{g})^!)_n \leq \dim (\widehat{\mathfrak{g}}_k^!)_n
\]
for all $n \geq 1$. This follows from the BRST reduction: $(\mathcal{W}^!)_n = H^0_{\mathrm{DS}}((\widehat{\mathfrak{g}}^!)_n)$ is a subquotient. The inequality is strict for $\mathfrak{sl}_2$: for example at $n = 3$, $\dim (\mathrm{Vir}_c^!)_3 = 5 < 15 = \dim (\widehat{\mathfrak{sl}}_2^!)_3$.
\end{remark}

\begin{remark}[Conformal weight filtration on bar cohomology]\label{rem:conformal-weight-filtration}
\index{bar complex!conformal weight filtration}
Each Koszul dual component $(\cA^!)_n$ carries a natural filtration by conformal weight:
\[
F^{\leq h} (\cA^!)_n = \bigl\{[\alpha] \in (\cA^!)_n :
\text{all generators in $\alpha$ have weight $\leq h$}\bigr\}.
\]
The associated graded $\operatorname{gr}^h (\cA^!)_n$ refines the single number
$\dim (\cA^!)_n$ into a bigraded Hilbert series
$\sum_{n,h} \dim \operatorname{gr}^h (\cA^!)_n \cdot x^n q^h$.
For the Virasoro algebra, the conformal weight grading on
$(\mathrm{Vir}_c^!)_n$ comes from the weight of the stress tensor~$T$
(weight~$2$) and its descendants, giving a $q$-refinement of the
Motzkin difference sequence.
\end{remark}

\begin{theorem}[Motzkin path model for Virasoro bar cohomology; \ClaimStatusConditional]
\label{thm:motzkin-path-model}
\index{Motzkin numbers!path model}
\index{Virasoro algebra!bar cohomology!combinatorial model}
\index{bar complex!combinatorial model}
There exists a bijection
\begin{equation}\label{eq:motzkin-bijection}
\operatorname{Basis}(H^n(\bar{B}(\mathrm{Vir}_c)))
\;\xleftrightarrow{\;\sim\;}
\bigl\{\text{Motzkin paths of length $n{+}1$ starting with $U$}\bigr\}
\end{equation}
respecting the conformal weight filtration
\textup{(}Remark~\textup{\ref{rem:conformal-weight-filtration})},
where $U = (1,1)$, $D = (1,-1)$, $F = (1,0)$
are up, down, and flat steps of a lattice path staying
at height $\geq 0$ and ending at height $0$.
\end{theorem}

\begin{proof}
The proof has three steps: dimension matching,
explicit construction of the bijection, and filtration compatibility.

\medskip
\noindent\emph{Step~1: Dimension matching.}
The Virasoro algebra $\mathrm{Vir}_c$ is chiral Koszul
(Theorem~\ref{thm:virasoro-chiral-koszul}), so
by Theorem~\ref{thm:spectral-sequence-collapse} the
PBW spectral sequence on $\bar{B}(\mathrm{Vir}_c)$
collapses at~$E_2$, giving
$\dim H^n(\bar{B}(\mathrm{Vir}_c)) = M(n{+}1) - M(n)$,
where $M(n)$ denotes the $n$-th Motzkin number.
On the combinatorial side, the standard Motzkin recurrence
\[
M(n{+}1) = M(n) + \sum_{k=0}^{n-1} M(k)\,M(n{-}1{-}k)
\]
decomposes paths of length $n{+}1$ by their first step:
$F$-initial paths biject with all paths of length~$n$
(contributing $M(n)$), while $U$-initial paths biject with
pairs of shorter paths via the first return decomposition
(contributing the convolution sum). Therefore the number
of $U$-initial Motzkin paths of length $n{+}1$ is exactly
$M(n{+}1) - M(n)$.

\medskip
\noindent\emph{Step~2: Construction of the bijection.}
The $E_2$-collapse identifies $H^n(\bar{B}(\mathrm{Vir}_c))$
with the degree-$n$ component of the Koszul dual
$\mathrm{Vir}_c^!$, which is spanned by bar monomials
$[v_1 \otimes \cdots \otimes v_n]$ with each $v_i$ in
the augmentation ideal $\overline{\mathrm{Vir}}_c$
(i.e., descendant modes of the stress tensor~$T$).
We define a map $\Phi$ from such basis elements to
$U$-initial Motzkin paths of length $n{+}1$ by reading
the bar monomial left to right and recording the OPE
type between consecutive tensor factors.

The Virasoro OPE $T(z)T(w)$ has three singular terms:
\begin{itemize}
\item the quartic pole $\frac{c/2}{(z-w)^4}$ (central,
 no new generator produced; record a \emph{flat} step~$F$);
\item the quadratic pole $\frac{2T(w)}{(z-w)^2}$ (produces
 a generator at the same level; record an \emph{up} step~$U$);
\item the simple pole $\frac{\partial T(w)}{(z-w)}$
 (descent to a derivative; record a \emph{down} step~$D$).
\end{itemize}
Concretely, given a bar monomial
$[v_1 \otimes \cdots \otimes v_n]$ representing a class
in $H^n(\bar{B}(\mathrm{Vir}_c))$, set:
\begin{itemize}
\item $s_0 = U$ (the first tensor slot always introduces
 a new generator from the augmentation ideal);
\item for $1 \leq i \leq n$, the step $s_i \in \{U,D,F\}$
 is determined by the leading OPE pole order
 between $v_i$ and $v_{i+1}$ in the bar differential:
 pole order~$2$ gives $U$, pole order~$1$ gives $D$,
 pole order~$4$ (central) gives~$F$.
\end{itemize}
The path $\Phi([v_1 \otimes \cdots \otimes v_n])
= (s_0, s_1, \ldots, s_n)$ has length $n{+}1$ and
starts with~$U$ by construction.

\emph{Non-negativity.}
The path stays at height $\geq 0$ because the conformal
weight filtration on $\bar{B}(\mathrm{Vir}_c)$ is
bounded below: at position~$i$, the running height equals
the number of active generators minus the number of
completed contractions up to that point.
The PBW filtration degree of $v_1 \otimes \cdots \otimes v_i$
is a non-negative integer that equals the path height
after step~$i$ (each $U$ adds a PBW generator, each~$D$
completes a contraction reducing PBW degree by~1,
each~$F$ is a central scalar leaving PBW degree unchanged).
Since PBW degree is non-negative, the path never
drops below the $x$-axis.

\emph{Return to zero.}
The path ends at height~$0$ because elements
of bar cohomology $H^n(\bar{B}(\mathrm{Vir}_c))$
are cocycles: the bar differential closes all open
contractions. In PBW terms, a cohomology class in
bar degree~$n$ has net PBW excess zero (all generators
introduced by $U$-steps are consumed by matching
$D$-steps or neutralized by $F$-steps that contribute
central scalars), so the terminal height is~$0$.

\emph{Bijectivity.}
The map $\Phi$ is injective: distinct bar cohomology
classes produce distinct OPE-type sequences, since the
leading pole order between consecutive tensor factors
is a cohomological invariant (the PBW spectral sequence
collapses, so the OPE type is determined at the $E_2$
page and does not change under further differentials).
Since both sets have cardinality $M(n{+}1) - M(n)$
by Step~1, $\Phi$ is a bijection.

\medskip
\noindent\emph{Step~3: Filtration compatibility.}
The conformal weight filtration on $H^n(\bar{B}(\mathrm{Vir}_c))$
(Remark~\ref{rem:conformal-weight-filtration})
assigns to a bar monomial $[v_1 \otimes \cdots \otimes v_n]$
the total conformal weight $h = \sum_{i=1}^n \Delta(v_i)$,
where $\Delta(v_i)$ is the conformal weight of~$v_i$.

Under the bijection~$\Phi$, this corresponds to the
\emph{area} under the Motzkin path (the sum of heights
$\sum_{i=0}^{n} h_i$ where $h_i$ is the height after
step~$i$): each up step $U$ at position~$i$ contributes
a generator of weight~$2$ (the stress tensor) and
increments all subsequent heights by~$1$; each down
step $D$ records a derivative $\partial T$ of weight~$3$
that reduces subsequent heights by~$1$; each flat
step $F$ records a central term of weight~$0$.

The total conformal weight is therefore
$h = 2(\text{\# of }U\text{-steps}) + 3(\text{\# of }D\text{-steps})$,
which by the constraint that the path starts and ends
at height~$0$ (so $\#U = \#D$) gives
$h = 5 \cdot \#U = 5 \cdot (\text{area}/\text{average height})$.

More precisely, the conformal weight of the bar element
and the area statistic on the path are related by an
explicit affine function of the step counts, so the
filtration $F^{\leq h} H^n(\bar{B}(\mathrm{Vir}_c))$
corresponds to the set of $U$-initial Motzkin paths
of length $n{+}1$ with area $\leq h$.
This establishes the filtration compatibility.
\end{proof}

\begin{remark}[\texorpdfstring{$\beta\gamma$}{beta-gamma} bar dimensions: conventions]\label{rem:betagamma-conventions}
\index{betagamma@$\beta\gamma$ system!bar dimension conventions}
The $\beta\gamma$ bar complex is infinite-dimensional in each bar
degree $n \geq 1$ (Proposition~\ref{prop:bg-bar-coalg}(2)):
without a conformal weight bound,
$\dim(\bar{B}^n) = \infty$, since the augmentation ideal
$\overline{\beta\gamma}$ contains generators $\beta_{-m-\lambda}$
and $\gamma_{-m-1+\lambda}$ for all $m \geq 0$.
The generating function in the bigraded sense is
\[
\sum_{n,h} \dim(\bar{B}^n_h)\, q^h t^n
= \prod_{m \geq 0}
 \frac{1}{(1-q^{m+\lambda}t)(1-q^{m+1-\lambda}t)}.
\]
The values in Table~\ref{tab:bar-dimensions}
are \emph{Koszul dual} dimensions $\dim (\beta\gamma^!)_n$
at the lowest-weight truncation (only the modes
$\beta_{-\lambda}$ and $\gamma_{\lambda-1}$, no conformal descendants),
with closed-form Hilbert series
$\sum_{n \geq 0} \dim (\beta\gamma^!)_n \cdot x^n = \sqrt{(1{+}x)/(1{-}3x)}$
(Theorem~\ref{thm:ds-bar-gf-discriminant}).

The formula $\operatorname{rank}(\bar{B}^n) = 2 \cdot 3^{n-1}$
(Theorem~\ref{thm:betagamma-bar-dim}) gives the \emph{chain space}
dimension, the total number of bar-$n$ elements before passing to
cohomology. This is the $\beta\gamma$ analogue of the Kac--Moody
formula $\dim(\bar{B}^n(\hat{\fg}_k)) = (\dim\fg)^n \cdot (n{-}1)!$
(Lemma~\ref{lem:bar-dims-level-independent}). In both cases,
the chain space is much larger than bar cohomology:
\begin{center}
\begin{tabular}{c|ccccc}
$n$ & 1 & 2 & 3 & 4 & 5 \\
\hline
chain $2 \cdot 3^{n-1}$ & 2 & 6 & 18 & 54 & 162 \\
cohom (Table) & 2 & 4 & 10 & 26 & 70
\end{tabular}
\end{center}
The inequality chain $\geq$ cohomology holds at every degree.
\end{remark}

\begin{corollary}[{\texorpdfstring{$\beta\gamma$}{beta-gamma} generating function via discriminant}; \ClaimStatusConditional]
\label{cor:betagamma-inverse-discriminant}
\index{betagamma@$\beta\gamma$ system!discriminant relation}
The $\beta\gamma$ bar cohomology generating function is the inverse square root of the shared discriminant:
\[
P_{\beta\gamma}(x) = \frac{1}{\sqrt{\Delta(x)}}\,(1+x),
\qquad \Delta(x) = 1-2x-3x^2.
\]
In particular, $P_{\beta\gamma}$ has singularity type $\Delta^{-1/2}$ while $P_{\widehat{\mathfrak{sl}}_2}$ and $P_{\mathrm{Vir}}$ have type $\Delta^{1/2}$, explaining the different sub-exponential corrections ($n^{-1/2}$ vs $n^{-3/2}$).
\end{corollary}

\begin{proof}
From~\eqref{eq:gf-betagamma}: $P_{\beta\gamma}(x) = \sqrt{(1+x)/(1-3x)} = (1+x)/\sqrt{(1+x)(1-3x)} = (1+x)/\sqrt{\Delta(x)}$.
\end{proof}

\section*{Genus-$g$ obstruction coefficients}
\label{sec:genus-obstruction-table}
\index{obstruction!genus-g table|textbf}
\index{Faber--Pandharipande formula!numerical values}

By the genus universality theorem
(Theorem~\ref{thm:genus-universality}), the scalar genus-$g$
coefficient for uniform-weight algebras is
$F_g^{\mathrm{sc}}(\cA)=\kappa(\cA) \cdot \lambda_g^{\mathrm{FP}}$,
while $H_D^K$ gives the signed Hodge character
$\operatorname{ch}_g(\mathfrak O_g^K(\cA))=
(-1)^g\kappa(\cA)\lambda_g$. Here
\[
\lambda_g^{\mathrm{FP}}
 = \frac{2^{2g-1}-1}{2^{2g-1}} \cdot \frac{|B_{2g}|}{(2g)!}\,.
\]
The first ten values are:

\begin{proposition}[Faber--Pandharipande coefficients \cite{FP03}; \ClaimStatusConditional]
\label{prop:fp-coefficients}
\index{Bernoulli numbers!obstruction coefficients}
With the Faber--Pandharipande convention used in
Theorem~\ref{thm:genus-universality}, the scalar coefficients are:
\begin{align}
\lambda_1^{\mathrm{FP}} &= \tfrac{1}{2} \cdot \tfrac{1/6}{2}
 = \tfrac{1}{24}\,,
 \label{eq:lambda1} \\
\lambda_2^{\mathrm{FP}} &= \tfrac{7}{8} \cdot \tfrac{1/30}{24}
 = \tfrac{7}{5760}\,,
 \label{eq:lambda2} \\
\lambda_3^{\mathrm{FP}} &= \tfrac{31}{32} \cdot \tfrac{1/42}{720}
 = \tfrac{31}{967680}\,,
 \label{eq:lambda3} \\
\lambda_4^{\mathrm{FP}} &= \tfrac{127}{128} \cdot \tfrac{1/30}{40320}
 = \tfrac{127}{154828800}\,,
 \label{eq:lambda4} \\
\lambda_5^{\mathrm{FP}} &= \tfrac{511}{512} \cdot \tfrac{5/66}{3628800}
 = \tfrac{73}{3503554560}\,,
 \label{eq:lambda5} \\
\lambda_6^{\mathrm{FP}} &= \tfrac{2047}{2048} \cdot \tfrac{691/2730}{479001600}
 = \tfrac{1414477}{2678117105664000}\,,
 \label{eq:lambda6} \\
\lambda_7^{\mathrm{FP}} &= \tfrac{8191}{8192} \cdot \tfrac{7/6}{87178291200}
 = \tfrac{8191}{612141052723200}\,,
 \label{eq:lambda7} \\
\lambda_8^{\mathrm{FP}} &= \tfrac{32767}{32768} \cdot \tfrac{3617/510}{20922789888000}
 = \tfrac{16931177}{49950709902213120000}\,,
 \label{eq:lambda8} \\
\lambda_9^{\mathrm{FP}} &= \tfrac{131071}{131072} \cdot \tfrac{43867/798}{6402373705728000}
 = \tfrac{5749691557}{669659197233029971968000}\,,
 \label{eq:lambda9} \\
\lambda_{10}^{\mathrm{FP}} &= \tfrac{524287}{524288} \cdot \tfrac{174611/330}{2432902008176640000}
 = \tfrac{91546277357}{420928638260761696665600000}\,.
 \label{eq:lambda10}
\end{align}
Here $|B_{2g}|$ denotes the absolute value of the $2g$-th
Bernoulli number: $B_2 = 1/6$, $B_4 = -1/30$, $B_6 = 1/42$,
$B_8 = -1/30$, $B_{10} = 5/66$, $B_{12} = -691/2730$,
$B_{14} = 7/6$, $B_{16} = -3617/510$, $B_{18} = 43867/798$,
$B_{20} = -174611/330$.
The rapid growth of $|B_{2g}|$ (asymptotically
$|B_{2g}| \sim 2 \cdot (2g)! / (2\pi)^{2g}$) is exactly
compensated by the $(2g)!$ denominator, so
$\lambda_g^{\mathrm{FP}} \sim 2/(2\pi)^{2g}$ decreases
super-exponentially.
\end{proposition}

\begin{proof}
Direct substitution of the Bernoulli numbers into the
Faber--Pandharipande formula; see
Theorem~\ref{thm:genus-universality}(iii).
\end{proof}

Table~\ref{tab:genus-obstructions} gives the explicit values
$F_g^{\mathrm{sc}}(\cA)=\kappa(\cA) \cdot \lambda_g^{\mathrm{FP}}$ for
$g = 1, \ldots, 5$ .

\begin{table}[ht]
\centering
\caption{Scalar genus-$g$ coefficients $F_g^{\mathrm{sc}}(\cA)=\kappa(\cA) \cdot \lambda_g^{\mathrm{FP}}$ through $g = 5$ }
\label{tab:genus-obstructions}
\index{obstruction!numerical values}
\renewcommand{\arraystretch}{1.5}
{\footnotesize
\begin{tabular}{|l|c|c|c|c|c|c|}
\hline
\textbf{Algebra $\cA$}
 & $\boldsymbol{\kappa(\cA)}$
 & $\boldsymbol{F^{\mathrm{sc}}_1}$
 & $\boldsymbol{F^{\mathrm{sc}}_2}$
 & $\boldsymbol{F^{\mathrm{sc}}_3}$
 & $\boldsymbol{F^{\mathrm{sc}}_4}$
 & $\boldsymbol{F^{\mathrm{sc}}_5}$ \\
\hline
\multicolumn{7}{|c|}{\textit{Free fields}} \\
\hline
Free fermion $\psi$ ($c = \tfrac{1}{2}$)
 & $\tfrac{1}{4}$
 & $\tfrac{1}{96}$
 & $\tfrac{7}{23040}$
 & $\tfrac{31}{3870720}$
 & $\tfrac{127}{619315200}$
 & $\tfrac{73}{14014218240}$ \\
\hline
$bc$ ghosts (weight $\lambda$)
 & $c/2$
 & $c/48$
 & $7c/11520$
 & $31c/1935360$
 & $127c/309657600$
 & $73c/7007109120$ \\
\hline
\multicolumn{7}{|c|}{\textit{Heisenberg}} \\
\hline
$\mathcal{H}_\kappa$ (rank~$1$)
 & $\kappa$
 & $\kappa/24$
 & $7\kappa/5760$
 & $31\kappa/967680$
 & $127\kappa/154828800$
 & $73\kappa/3503554560$ \\
\hline
\multicolumn{7}{|c|}{\textit{Affine Kac--Moody ($t = k + h^\vee$)}} \\
\hline
$\widehat{\mathfrak{sl}}_2$ (level~$k$)
 & $\frac{3t}{4}$
 & $\frac{t}{32}$
 & $\frac{7t}{7680}$
 & $\frac{31t}{1290240}$
 & $\frac{127t}{206438400}$
 & $\frac{73t}{4671406080}$ \\[4pt]
\hline
$\widehat{\mathfrak{sl}}_3$ (level~$k$)
 & $\frac{4t}{3}$
 & $\frac{t}{18}$
 & $\frac{7t}{4320}$
 & $\frac{31t}{725760}$
 & $\frac{127t}{116121600}$
 & $\frac{73t}{2627665920}$ \\[4pt]
\hline
$\widehat{\fg}_k$ (general)
 & $\frac{td}{2h^\vee}$
 & $\frac{td}{48 h^\vee}$
 & $\frac{7td}{11520 h^\vee}$
 & $\frac{31td}{1935360 h^\vee}$
 & $\frac{127td}{309657600 h^\vee}$
 & $\frac{73td}{7007109120 h^\vee}$ \\[4pt]
\hline
\multicolumn{7}{|c|}{\textit{$\mathcal{W}$-algebras}} \\
\hline
$\mathrm{Vir}_c$
 & $c/2$
 & $c/48$
 & $7c/11520$
 & $31c/1935360$
 & $127c/309657600$
 & $73c/7007109120$ \\[4pt]
\hline
$\mathcal{W}_3^k$\textsuperscript{$\mathrm{mod/comp}$}
 & $5c/6$
 & $5c/144$
 & $7c/6912$
 & $31c/1161216$
 & $127c/185794560$
 & $73c/4204265472$ \\[4pt]
\hline
\end{tabular}
}
\end{table}

\noindent
The principal $\mathcal W_3$ row in the preceding genus table is read
on $H_{\mathrm{diag}}^{g=1}+H_{W_3}^{\mathrm{DS/bar}}$.

\begin{remark}[Complementarity at each genus]\label{rem:genus-complementarity}
\index{complementarity!genus-g}
On the scalar Hodge component of Theorem~\ref{thm:quantum-complementarity-main},
the sign-normalized characters satisfy
$(-1)^g\operatorname{ch}_g(\mathfrak O_g^K(\cA))
+(-1)^g\operatorname{ch}_g(\mathfrak O_g^K(\cA^!))
=\lambda_g(\kappa+\kappa')$.
For affine Kac--Moody algebras,
$\kappa(\widehat{\fg}_k) + \kappa(\widehat{\fg}_{-k-2h^\vee}) = 0$
(the shifted level $t = k + h^\vee$ maps to $-t$, and
$\kappa = td/(2h^\vee)$ is linear in~$t$).
Hence the sum of the two sign-normalized Hodge characters vanishes.
For $\widehat{\mathfrak{sl}}_2$:
$3(k{+}2)/96 + 3(-k{-}2)/96 = 0$.
At the critical level $k = -h^\vee$, both $\kappa = 0$ and $\kappa' = 0$,
so both Hodge characters vanish.

The Virasoro algebra has the unconditional modular conductor
$\kappa+\kappa'=13$.  On
$H_{\mathrm{diag}}^{g=1}+H_{W_3}^{\mathrm{DS/bar}}$, the principal
$\mathcal W_3$ algebra has the conditional modular conductor
$\kappa+\kappa'=250/3$; its exact central conductor is~$100$
\textup{(}Theorem~\ref{thm:genus-universality}(ii) and
Remark~\ref{rem:vir-vs-km-complementarity}\textup{)}.
\end{remark}

\begin{table}[ht]
\centering
\caption{Higher-genus Faber--Pandharipande coefficients $\lambda_g$ for $g = 6, \ldots, 10$}
\label{tab:genus-obstructions-high}
\index{obstruction!higher genus values}
\renewcommand{\arraystretch}{1.5}
{\footnotesize
\begin{tabular}{|c|c|c|}
\hline
$g$ & $\lambda_g$ (exact) & $\lambda_g$ (decimal) \\
\hline
6 & $\dfrac{1414477}{2678117105664000}$ & $5.28 \times 10^{-10}$ \\[6pt]
\hline
7 & $\dfrac{8191}{612141052723200}$ & $1.34 \times 10^{-11}$ \\[6pt]
\hline
8 & $\dfrac{16931177}{49950709902213120000}$ & $3.39 \times 10^{-13}$ \\[6pt]
\hline
9 & $\dfrac{5749691557}{669659197233029971968000}$ & $8.59 \times 10^{-15}$ \\[6pt]
\hline
10 & $\dfrac{91546277357}{420928638260761696665600000}$ & $2.17 \times 10^{-16}$ \\[6pt]
\hline
\end{tabular}
}
\end{table}

\noindent
For any uniform-weight algebra $\cA$ with curvature
parameter~$\kappa(\cA)$, the $H_D^K$ Hodge character is
$\operatorname{ch}_g(\mathfrak O_g^K(\cA))=
(-1)^g\kappa(\cA)\lambda_g$, and the $H_D^{\mathrm{graph}}$ scalar
coefficient is
$F_g^{\mathrm{sc}}(\cA)=\kappa(\cA)\cdot\lambda_g^{\mathrm{FP}}$
(Table~\ref{tab:genus-obstructions}; the multi-weight full coefficient
requires $\delta F_g^{\mathrm{cross}}$,
cf.~Theorem~\ref{thm:multi-weight-genus-expansion}). The numerators of
$\lambda_g$ involve Bernoulli numbers via
$\lambda_g = \tfrac{2^{2g-1}-1}{2^{2g-1}} \cdot
\tfrac{|B_{2g}|}{(2g)!}$
(Proposition~\ref{prop:fp-coefficients}).
At $g = 6$, the numerator $1414477 = 691 \cdot 2047$ reflects the
appearance of $B_{12} = -691/2730$, the first ``large'' Bernoulli
number (responsible for the Kervaire invariant and Ramanujan's
$\tau$-function).

\section*{Spectral sequence data}
\label{sec:spectral-sequence-data}
\index{spectral sequence!collapse data|textbf}

The bar complex $\barBgeom(\cA)$ carries a filtration by bar degree,
inducing a spectral sequence with
$E_1^{p,q} = H^q(\barBgeom^p(\cA))$
converging to $H^{p+q}(\barBgeom(\cA))$.
The Koszul property of~$\cA$ is equivalent to collapse at~$E_2$
(Theorem~\ref{thm:spectral-sequence-collapse}).
Table~\ref{tab:spectral-data} summarizes the collapse behavior
for each algebra class.

\begin{table}[ht]
\centering
\caption{Spectral sequence collapse data for the bar complex}
\label{tab:spectral-data}
\index{spectral sequence!E2 page@$E_2$ page}
\renewcommand{\arraystretch}{1.4}
{\small
\begin{tabular}{|l|c|l|}
\hline
\textbf{Algebra $\cA$}
 & \textbf{Collapse page}
 & \textbf{Notes} \\
\hline
Free fermion $\psi$
 & $E_1$
 & Exact Koszul; $\barBgeom$ acyclic \\
\hline
$bc$--$\beta\gamma$
 & $E_1$
 & Exact Koszul; bar differential strict \\
\hline
Heisenberg $\mathcal{H}_\kappa$
 & $E_2$
 & Koszul, curved ($m_0 \ne 0$) \\
\hline
$\widehat{\fg}_k$ (generic~$k$)
 & $E_2$
 & Koszul; $d_1$ from OPE residues \\
\hline
$\widehat{\fg}_{-h^\vee}$ (critical)
 & $E_2$ \textup{(}PBW\textup{)}, not diagonal
 & Feigin--Frenkel center; full Koszulness fails \\
\hline
$\mathrm{Vir}_c$ (generic~$c$)
 & $E_2$
 & Koszul; two generators ($T$, $\mathbf{1}$) \\
\hline
$\mathcal{W}_3^k$ (generic~$k$)
 & $E_2$
 & Koszul; three generators ($T$, $W$, $\mathbf{1}$) \\
\hline
Yangian $Y(\fg)$
 & finite DK core; full collapse open
 & RTT ordered-bar duality on the proved core \\
\hline
Lattice VOA $V_\Lambda$
 & $E_1$
 & Heisenberg sector acyclic \\
\hline
\end{tabular}
}
\end{table}

\begin{proposition}[Spectral sequence collapse; \ClaimStatusConditional]
\label{prop:spectral-collapse-summary}
\index{Koszul algebra!spectral sequence characterization}
For the algebras listed in Table~\textup{\ref{tab:spectral-data}}:
\begin{enumerate}
\item The free fermion, $bc$--$\beta\gamma$, and lattice VOA
 bar complexes are exact Koszul: the spectral sequence collapses
 at~$E_1$ and the bar differential is strict ($d^2 = 0$ on generators).
\item The Heisenberg, generic affine Kac--Moody, generic Virasoro,
 and generic $\mathcal{W}_3$ bar complexes are Koszul: the spectral
 sequence collapses at~$E_2$, and the dual algebra is identified
 with $E_2^{*,0}$.
\item At the critical level $k = -h^\vee$, the affine Kac--Moody
 PBW spectral sequence still reaches its $E_2$ page, but that page is
 not diagonally concentrated. The Feigin--Frenkel center contributes
 off-diagonal classes, so the full Koszul criterion fails.
\end{enumerate}
\end{proposition}

\begin{proof}
Claims (1) and (2) are proved in the individual example chapters:
Theorem~\ref{thm:fermion-bar-coalg} (fermion),
Theorem~\ref{thm:betagamma-bc-koszul} ($bc$--$\beta\gamma$),
Theorem~\ref{thm:heisenberg-bar} (Heisenberg),
Theorem~\ref{thm:universal-kac-moody-koszul} (Kac--Moody),
Theorem~\ref{thm:w-algebra-koszul-main} ($\mathcal{W}$-algebras),
and Theorem~\ref{thm:lattice:unimodular-self-dual} (lattice VOA).
Claim (3) is the critical-level distinction recorded in
Remark~\ref{rem:pbw-vs-diagonal-critical}: PBW degeneration is
$k$-independent, but at $k=-h^\vee$ the $E_2$ page contains
$\Omega^*(\mathrm{Op}_{\fg^\vee}(D))$-classes from the
Feigin--Frenkel center. Thus the diagonal Ext criterion required by
Theorem~\ref{thm:spectral-sequence-collapse} fails even though the
PBW page itself is defined.
\end{proof}

\begin{remark}[Yangian full-collapse obligation]\label{rem:yangian-collapse-conj}
\ClaimStatusConjectured
\index{Yangian!spectral sequence collapse}
The RTT presentation (Theorem~\ref{thm:yangian-bar-rtt}) and the
ordered $R\mapsto R^{-1}$ duality theorem
(Theorem~\ref{thm:yangian-koszul-dual}) give the finite DK core used
in the census. The open assertion is stronger: the spectral sequence
for the ordinary-derived/completed/coderived enlargement of
$\barBgeom(Y(\fg))$ should collapse at~$E_2$ on the full
$\Eone$-chiral Yangian landscape, extending the finite core beyond
the evaluation-generated and RTT windows. The degree-$3$ Hilbert data
are consistent with this full-collapse lift but do not prove it.
A proof would require extending the FM compactification arguments of
Theorem~\ref{thm:spectral-sequence-collapse} to the $\Eone$ setting.
The domain of this extension is discussed in
Remark~\ref{rem:yangian-three-theorems}.
(Contributing to Conjecture~\ref{conj:master-dk-kl}.)
\end{remark}

\section*{Completion kinematics of the standard landscape}
\label{sec:completion-kinematics-census}
\index{completion entropy!census}
\index{primitive cumulant!census}
\index{Koszul radius!census}

The completion-tower theorem
(Theorem~\ref{thm:completed-bar-cobar-strong}) resolves the formal
MC4 problem; what remains is a \emph{kinematic} phase diagram for the
completed dual. The primitive generating series
$G_\cA(t) = \tfrac{1-t}{t}\bigl(\chi_\cA(t)-1\bigr)$
(Definition~\ref{def:primitive-generating-series}) and the completion
entropy $h_K(\cA) = \log(\rho_\cA^{-1})$
(Definition~\ref{def:completion-entropy}) are computable from the
vacuum character alone, without curvature data.

\begin{table}[ht]
\centering
\begin{tabular}{lccccl}
\toprule
\textbf{Family} & $\rho_K$ & $h_K$ & \textbf{Shadow depth}
 & $\Delta(t)$ starts & \textbf{Class} \\
\midrule
Heisenberg & $1$ & $0$ & $2$ & $0$ & G \\
$\hat{\mathfrak{g}}_k$ \textup{(}non-critical\textup{)} & $1/\dim\mathfrak{g}$ & $\log\dim\mathfrak{g}$ & $3$ & $0$ & L \\
$\beta\gamma$ & $\tfrac{1}{2}$ & $\log 2$ & $4$ & $0$ & C \\
$\mathrm{Vir}_c$ & $\approx 0.567$ & $\approx 0.567$ & $\infty$ & $t^3 + \cdots$ & M \\
$\Walg_3$ & $\approx 0.462$ & $\approx 0.772$ & $\infty$ & $t^3 + t^4 + \cdots$ & M \\
$\Walg_4$ & $\approx 0.434$ & $\approx 0.835$ & $\infty$ & $t^3 + \cdots$ & M \\
$\Walg_5$ & $\approx 0.424$ & $\approx 0.857$ & $\infty$ & $t^3 + \cdots$ & M \\
$\Walg_\infty$ & $\approx 0.418$ & $\approx 0.872$ & $\infty$ & $t^3 + \cdots$ & M \\
\bottomrule
\end{tabular}
\caption{Completion kinematics of the standard landscape.
 $\rho_K = $ Koszul radius (smallest positive root of $G_\cA(t)=1$);
 $h_K = \log(\rho_K^{-1})$ is the exponential growth rate of
 reduced-weight bar windows; $\Delta(t)$ is the primitive defect
 beyond declared generators.
 Classes: G~(Gaussian, $r_{\max}=2$),
 L~(Lie/tree, $r_{\max}=3$),
 C~(contact/quartic, $r_{\max}=4$),
 M~(mixed, $r_{\max}=\infty$).
 The $\Walg_N$ entropy ladder is monotone increasing with limit
 $h_K(\Walg_\infty) \approx 0.872$; the $\Walg_\infty$ character
 is governed by the MacMahon plane-partition function
 (Proposition~\ref{prop:wn-entropy-ladder}).}
\label{tab:completion-kinematics}
\end{table}

\begin{remark}[Interpretation]
\label{rem:completion-kinematics-interpretation}
The table reveals a sharp kinematic stratification: G/L/C-class
algebras (finite shadow depth) have zero primitive defect and rational
Koszul radius; M-class algebras (infinite tower) have non-zero defect
starting at $t^3$ and irrational Koszul radius. The entropy ladder
$h_K(\Walg_2) < h_K(\Walg_3) < \cdots < h_K(\Walg_\infty)$ quantifies
the completion cost of each additional spin field. This is a genuine
\emph{kinematic} invariant: it is independent of the central charge
$c$, the curvature $\kappa(\cA)$, and all dynamic data. The
$c \mapsto 26-c$ involution changes the dynamic shadow but preserves
the kinematic entropy.
\end{remark}

\section*{Cross-reference guide}
\label{sec:cross-ref-guide}
\index{master table of invariants!cross-references}

\begin{remark}[Locating the computations]\label{rem:cross-ref-guide}
Each row of the Master Table~\textup{\ref{tab:master-invariants}} is
computed in the following locations.
\begin{enumerate}
\item \emph{Free fermion}:
 bar complex (Theorem~\ref{thm:fermion-bar-coalg}),
 Koszul dual $\mathrm{Sym}^{\mathrm{ch}}(\gamma)$
 (Theorem~\ref{thm:betagamma-bc-koszul}),
 $\kappa = \tfrac{1}{4}$ (Table~\ref{tab:master-invariants}),
 three-theorem synthesis
 (Remark~\ref{rem:free-field-three-theorems}).
\item \emph{$bc$--$\beta\gamma$}:
 bar complex (Theorem~\ref{thm:betagamma-bar-complex}),
 Koszul duality
 (Proposition~\ref{prop:bc-betagamma-orthogonality},
 Theorem~\ref{thm:betagamma-bc-koszul}),
 central charge (Theorem~\ref{thm:betagamma-bc-koszul}),
 detailed computation (\S\ref{sec:betagamma-koszul-dual}).
\item \emph{Heisenberg}:
 bar complex (Theorem~\ref{thm:heisenberg-bar}),
 Koszul dual (Theorem~\ref{thm:heisenberg-koszul-dual-early}),
 genus expansion
 (\S\ref{sec:heisenberg-genus-expansion-eisenstein}),
 non-self-duality (Theorem~\ref{thm:heisenberg-not-self-dual}).
\item \emph{$\widehat{\mathfrak{sl}}_2$}:
 bar complex and Koszul dual
 (Theorem~\ref{thm:sl2-koszul-dual}),
	 genus-$1$ computation (\S\ref{sec:sl2-genus-one-computation}),
 three-theorem synthesis (Remark~\ref{rem:sl2-three-theorems}).
\item \emph{$\widehat{\mathfrak{sl}}_3$, $\widehat{E}_8$,
 general~$\widehat{\fg}_k$}:
 universal Koszul duality
 (Theorem~\ref{thm:universal-kac-moody-koszul}),
 complementarity
 (Theorem~\ref{thm:central-charge-complementarity}),
 obstruction coefficient
 (Theorem~\ref{thm:genus-universality}).
\item \emph{Virasoro}:
 $\mathcal{W}$-algebra Koszul duality
 (Theorem~\ref{thm:w-algebra-koszul-main}),
	 genus-$1$ computation (\S\ref{sec:virasoro-genus-one-computation}),
 $c + c' = 26$
 (Theorem~\ref{thm:central-charge-complementarity}).
\item \emph{$\mathcal{W}_3$ and $\mathcal{W}_N$}:
 Koszul duality (Theorem~\ref{thm:w-algebra-koszul-main}),
	 genus-$1$ computation (\S\ref{sec:w3-genus-one-computation}),
 general $\kappa$ formula (Theorem~\ref{thm:wn-obstruction}),
 three-theorem synthesis
 (Remark~\ref{rem:w-algebra-three-theorems}).
\item \emph{Yangian}:
 bar complex (Theorem~\ref{thm:yangian-bar-rtt}),
 Koszul dual (Theorem~\ref{thm:yangian-koszul-dual}),
 explicit computations (\S\ref{sec:yangian-bar-explicit}),
 three-theorem synthesis
 (Remark~\ref{rem:yangian-three-theorems}).
\item \emph{Lattice VOA}:
 bar complex (\S\ref{sec:lattice:bar}),
 Koszul duality (\S\ref{sec:lattice:koszul}),
 self-duality for unimodular lattices
 (Theorem~\ref{thm:lattice:unimodular-self-dual}),
 three-theorem synthesis
 (Remark~\ref{rem:lattice-three-theorems}).
\item \emph{Deformation quantization}:
 chiral Kontsevich theorem
 (Theorem~\ref{thm:chiral-kontsevich}),
 quantization construction
 (Construction~\ref{constr:quantization}),
 three-theorem synthesis
 (Remark~\ref{rem:deformation-three-theorems}).
\item \emph{Toroidal and elliptic}:
 elliptic-curve $\Eone$-chiral realization
 (Conjecture~\ref{V3-conj:toroidal-e1}),
 elliptic bar complex
 (\S\ref{V3-sec:elliptic-bar-heisenberg}).
\end{enumerate}
\end{remark}

\begin{remark}[Explicit non-coverage of the master table]
\label{rem:census-silent-non-coverage}
Five families appear in the shadow-obstruction tower census
(Table~\ref{tab:shadow-tower-census}) but do not receive a complete
row in the master invariant table
(Table~\ref{tab:master-invariants}). They are recorded here as
open frontier rather than folded into completed rows.
\begin{enumerate}[leftmargin=*]
 \item $\mathcal{N}{=}2$ superconformal $\mathrm{SCA}_c$: the
 shadow class is~$M$ with multi-component tower
 (Rem~\ref{rem:n2-shadow-tower}); the self-dual central charge
 is $c^{*}=9$ on the Ito--Tanaka involution
 $c\mapsto 6-c$ rescaled by the $U(1)$-current level. The Koszul
 conductor on the Virasoro subcomponent is $K=18$; on the full
 $\mathrm{SCA}$ involution the expression remains a frontier item.
 \item Logarithmic triplet $\mathcal{W}(p)$ ($p\ge 2$): class~$M$
 with $c = 1 - 6(p-1)^{2}/p$; $C_{2}$-cofinite but non-rational;
 Zhu-bounded Massey proof does not extend (Gurarie~1993,
 Flohr~1996 unbounded amplitudes). Koszul-conductor and
 complementarity entries are pending resolution of the
 character-amplitude bound.
 \item Cosets of the form $\operatorname{Com}(\cH_{\kappa},\cA)$
 beyond GKO and Kac--Roan--Wakimoto hook reductions: the coset
 slot of the fingerprint carries the governing datum, but the
 family-by-family master-table entries are not yet computed.
 \item Non-rational lattice VOAs $V_{L}$ with $L$ indefinite or
 non-positive-definite: the Gaussian-class entry
 (Table~\ref{tab:shadow-invariants-landscape}) assumes $L$
 positive-definite and unimodular for the $V_{L}$ self-duality;
 the indefinite case is open.
 \item Roots-of-unity quantum groups, i.e.\ $V_{k}(\mathfrak{g})$
 at admissible level $k+h^{\vee} = p/q$ with $(p,q)=1$ and $p>1$:
 the periodic-CDG admissible chain
 (Thm~\ref{thm:periodic-cdg-is-koszul-compatible}) secures
 compatibility with chiral Koszul duality on the universal locus. The
 master-table admissible row is exactly that universal restriction;
 simple-quotient admissible entries require separate bar-rank
 computations and are not supplied by the universal formula.
\end{enumerate}
The $\mathcal{N}{=}1,4$ superconformal rows and the super-Yangian
row are covered in parts and appear in dedicated chapters but are not
re-entered in this master table; the $\mathcal{N}{=}2$ row above is
only the scalar Ito--Tanaka locus recorded above. The
Schellekens $71$ rows at
$c=24$ are now governed explicitly by
Theorem~\ref{thm:schellekens-structured-subset}: the master table keeps
the coarse Niemeier / Monster rows, while the full
$24_{\mathrm{Niemeier}}+1_{V_1=0}+46_{\mathrm{nonlat}}$
structured decomposition lives in the
infinite-fingerprint chapter rather than as bare numerology.
\end{remark}

\medskip

\noindent
Part~\ref{part:physics-bridges} supplies the Feynman, BV--BRST,
holomorphic-topological, and 4d/2d comparison data. Part~\ref{part:seven-faces}
then records the concordance and modular Koszul duality
(\S\ref{sec:modular-koszul-theory}).

\section{Synthesis}
\label{sec:landscape-supplement}

\begin{remark}[K3$\times$E: four scalars and one primitive denominator]
\label{rem:landscape-1}
Let $Y=K3\times E$. The census row attached to $Y$ is not a
single $\kappa$-value. Four constructions give four scalars:
\[
 \begin{aligned}
	 \kappa_{\mathrm{cat}}(Y)
	 &= \chi(\cO_{K3})\cdot\chi(\cO_E) = 2\cdot 0 = 0,\\
	 \kappa_{\mathrm{ch}}^{\mathrm{Heis}}(Y)
	 &= 3,\\
	 \kappa_{\mathrm{BKM}}(\Delta_5)
	 &= \mathrm{wt}(\Delta_5)=c_{N=1}(0)/2 = 5,\\
	 \kappa_{\mathrm{fiber}}(K3)
	 &= 24,
	 \end{aligned}
\]
and the separate K3 fibre witness $\kappa_{\mathrm{cat}}(K3)=2$.
The primitive BKM denominator is $\Delta_5$. The reduced
Oberdieck--Pandharipande DT scalar uses the monic Borcherds product
\[
 D_5=64^{-1}\Delta_5,\qquad
 \Phi_{10}^{\mathrm{un}}=\Delta_5^2,\qquad
 \Phi_{10}^{\mathrm{OP}}=D_5^2=4096^{-1}\Phi_{10}^{\mathrm{un}},
\]
so
\[
 Z_{\mathrm{OP}}(K3\times E)
 := Z^{\mathrm{red}}_{\mathrm{DT}}(K3\times E)
 = -(\Phi_{10}^{\mathrm{OP}})^{-1}
 = -D_5^{-2}
 = -4096\,\Delta_5^{-2}.
\]
Thus $\Delta_5$ is the primitive Borcherds denominator,
$\Phi_{10}^{\mathrm{un}}=\Delta_5^2$ is the unnormalised Igusa square,
$\Phi_{10}^{\mathrm{OP}}$ is the OP square in the $D_5$ normalisation,
and $Z_{\mathrm{OP}}$ is the reduced DT scalar.  The Mukai lattice
supplies the exact arithmetic
$2c_+(\widetilde H(K3,\mathbb Z))=8$~\cite{Mukai1987}.  The chiral
identification $K^{\kappa_{\mathrm{ch}}}_{\mathsf B}=8$ and the
selection $\hbar^2=-1/8$ have status \ClaimStatusConjectured under
$H_{\mathsf B}$; after that selection, the scalar equality
$\hbar^2K^{\kappa_{\mathrm{ch}}}=-1$ is immediate.
Cross-ref Vol.~III Chapter~\ref{V3-ch:k3-chiral-bialgebra-platonic}.
\end{remark}

\begin{proposition}[K3/BKM scalar invariants and Hall--Drinfeld taxonomy;
\ClaimStatusConditional]
\label{prop:landscape-k3-bkm-taxonomy}
The K3$\times E$ Borcherds row contributes the scalar invariants
\[
\{\kappa_{\mathrm{cat}},\kappa_{\mathrm{ch}}^{\mathrm{Heis}},
\kappa_{\mathrm{BKM}},\kappa_{\mathrm{fiber}}\}
=\{0,3,5,24\}
\]
on the hypothesis package of
Theorem~\ref{thm:k3e-five-coordinate-computation}.  The Mukai lattice
also contributes the proved arithmetic value
\[
2c_+(\mathrm{Mukai}(K3))=8.
\]
Its interpretation as $K^\kappa_{\mathcal B}=8$ has status
\ClaimStatusConjectured under~$H_{\mathsf B}$.
It occupies the Hall--Borcherds/CoHA lane, structurally distinct from
the standard RTT Yangian $Y(\fg)$ and from the four $G/L/C/M$ shadow
archetypes.
\end{proposition}

\begin{proof}
The four displayed scalars have different definitions: categorical
Euler scalar, Heisenberg chiral scalar, Borcherds weight, and K3 fibre
Euler scalar. The Borcherds value is
$\kappa_{\mathrm{BKM}}(\Delta_5)=\mathrm{wt}(\Delta_5)=5$, while the
Mukai arithmetic uses the positive-signature index
$c_+(\mathrm{Mukai}(K3))=4$.  The passage from this arithmetic to a
chiral conductor is exactly $H_{\mathrm{scalar}}$ inside
$H_{\mathsf B}$. Proposition~\ref{prop:uc-K3-BKM-scalar-separation}
identifies the associated Vol.~III object as a Hall--Drinfeld object
built from CoHA data. A standard RTT Yangian requires a finite-rank
symmetrisable Cartan datum and an RTT presentation; the Borcherds
algebra attached to $\Delta_5$ has imaginary simple roots and a
Hall--Drinfeld double. The taxonomy and scalar invariants are therefore
disjoint.
\end{proof}

\begin{remark}[Scalar typing of rows adjacent to K3$\times$E]
\label{rem:landscape-2}
The neighbouring rows carry their own typed scalar data alongside the
four $K3\times E$ invariants.
\begin{center}
\renewcommand{\arraystretch}{1.25}
\begin{tabular}{p{0.19\textwidth}p{0.37\textwidth}p{0.36\textwidth}}
\textbf{Row} & \textbf{$\kappa$-datum} & \textbf{BKM / CoHA datum}\\
\hline
$E$ & $\kappa_{\mathrm{ch}}(E)=0$ & no $\Delta_5$ denominator\\
$K3$ & $\kappa_{\mathrm{ch}}(K3)=\kappa_{\mathrm{cat}}(K3)=2$
 & $\phi_{0,1}^{K3}$ is Jacobi input, not the $K3\times E$ denominator\\
$K3\times E$ &
$\{\kappa_{\mathrm{cat}}(K3\times E),
\kappa_{\mathrm{ch}}^{\mathrm{Heis}}(K3\times E),
\kappa_{\mathrm{BKM}}(\Delta_5),\kappa_{\mathrm{fiber}}(K3)\}
=\{0,3,5,24\}$ &
$\Delta_5$ is primitive;
$Z_{\mathrm{OP}}=-4096\Delta_5^{-2}$\\
$Q_5$ &
$\kappa_{\mathrm{ch}}^{\mathrm{Hodge}}(Q_5)=0$,
$\kappa_{\mathrm{ch}}^{\mathrm{BCOV}}(Q_5)=\chi(Q_5)/24=-25/3$ &
BCOV component, not $\Delta_5$\\
$\mathrm{Tot}(\omega_{\mathbb P^2})$ &
$\kappa_{\mathrm{ch}}=3/2$ &
local surface McKay locus\\
$\mathrm{Tot}(\cO(-1)^{\oplus2}\to\mathbb P^1)$ &
$\kappa_{\mathrm{ch}}=1$ &
resolved conifold, not a local surface\\
$\mathbb C^3$ &
$\mathrm{CoHA}(\mathbb C^3)=Y^+$ &
$\mathcal W_{1+\infty}$ appears only after the double / centre passage.
\end{tabular}
\end{center}
The universal BKM identity remains
$\kappa_{\mathrm{BKM}}(\Phi_N)=c_N(0)/2$ for
$N\in\{1,2,3,4,6\}$ \cite{Borcherds1998,Gritsenko1999}; it does not
license any additive formula identifying the Borcherds weight with the
chiral scalar plus $\chi(\cO_{\mathrm{fiber}})$
and it does not transfer the $K3\times E$ denominator to the quintic,
the conifold, local $\mathbb P^2$, or $\mathbb C^3$. Cross-ref
Chapter~\ref{V3-ch:k3-chiral-bialgebra-platonic}.
\end{remark}

\begin{remark}[Fake-Monster landscape row: lattice ambient $\mathrm{II}_{25,1}$, weight $12$]
\label{rem:landscape-fake-monster-row}
Parallel to the paramodular K3$\times$E row but with a disjoint lattice
ambient, the Fake-Monster Hall--Drinfeld double
$\mathbf H_{\mathrm{FakeMonster}}$ carries the Borcherds product
$\Phi_{12}$ on the Grassmannian of $\mathrm{II}_{26,2}$ as its
automorphic denominator, with BKM weight
\[
 \kappa_{\mathrm{BKM}}(\Phi_{12}) \;=\; \tfrac{1}{2}\,c^{\mathrm{FM}}(0) \;=\; 12
 \qquad\bigl(\text{Borcherds~\cite{Bor95} Thm~$13.1$;
 lift of $\theta_{\Lambda_{24}}/\eta^{24}$}\bigr).
\]
The Cartan $\mathrm{II}_{25,1} = \Lambda_{24}\oplus\mathrm{II}_{1,1}$
has signature $(25,1)$, rank~$26$; all imaginary-root multiplicities
equal the Leech rank~$24$. This row is \emph{not} a re-indexing of the
paramodular K3$\times$E row of
Remark~\ref{rem:landscape-1}: the Cartans
$\Lambda^{2,1}_{\mathrm{II}}$ and $\mathrm{II}_{25,1}$ admit no
signature-preserving embedding, the K3 lattice $\mathrm{II}_{3,19}$ is
not a sublattice of the Leech lattice
(Gritsenko--Nikulin~\cite{GritsenkoNikulin1998} Prop.~$2.5$), and the
Fourier channels at weight~$12$ do not interfere
(Remark~\ref{rem:nc-fake-monster-non-interference}). The cross-volume
placement is fixed in
Proposition~\ref{prop:kappa-bkm-cross-volume-reconciliation}:
$\kappa_{\mathrm{BKM}}(\Delta_5) = 5$ records the paramodular row at
$d = 3$ on CY$_3$ ambient $K3\times E$, while
$\kappa_{\mathrm{BKM}}(\Phi_{12}) = 12$ records the Fake-Monster row on
the lattice vertex algebra $V_{\mathrm{II}_{25,1}}$. The universal
identity $\kappa_{\mathrm{BKM}}(\Phi) = c^\Phi(0)/2$ applies on each
row with its own Fourier-zero coefficient.
\end{remark}

\section*{Canonical coefficient values}
\label{sec:canonical-coefficient-values}
\index{canonical values!coefficient table|textbf}

The five propositions below fix the coefficient values that appear
across all three volumes at the $c=-214$ K3 base point and at
adjacent CY rows. Each value is a single number determined by a
primary source. A different value is possible only after an explicit
change of convention.

\begin{proposition}[Pope--Romans--Shen projector coefficient at
$c=-214$; \ClaimStatusProvedElsewhere]
\label{prop:canonical-prs-coefficient-cm214}
\index{Pope--Romans--Shen!projector coefficient}
\index{$\mathcal W_\infty[c]$!primary-projector coefficient}
The $\mathcal W_\infty[c]$-primary-projector identity
\[
 W_k \;=\; e_k + \frac{-c}{22}\, \Lambda_Z,
 \qquad \Lambda_Z \;=\; {:}TT{:} - \tfrac{3}{10}\partial^2 T,
\]
of Pope--Romans--Shen \textup{(}\emph{Phys.~Lett.~B} $\mathbf{236}$ $\textup{(}1990\textup{)}$,
pages $173$--$178$, equation~\textup{(}$2.10$\textup{)}\textup{)} evaluates at the K3 base
point $c=-214$ to
\[
 \frac{-c}{22}\bigg|_{c=-214} \;=\; \frac{214}{22} \;=\; \frac{107}{11}.
\]
Consequently, at $c=-214$ the primary basis is
$e_4 = W_4 - \tfrac{107}{11}\,\Lambda_Z$, and the quasi-primary
expansion of $W_4$ in the $(:T\partial^2 T:,\ :(\partial T)^2:,\
\partial^4 T)$ basis carries coefficients $(1,\ -3/2,\ 321/10)$
\textup{(}Wang $1998$ \emph{Commun.~Math.~Phys.~}$\mathbf{195}$, Proposition~$4.2$\textup{)}.
Primary sources: Pope--Romans--Shen $1990$ \emph{Phys.~Lett.~B} $236$
equation~\textup{(}$2.10$\textup{)}; Wang $1998$ \emph{Commun.~Math.~Phys.~}$195$
Proposition~$4.2$.
\end{proposition}

\begin{proof}
Pope--Romans--Shen $1990$ equation $\textup{(}2.10\textup{)}$ identifies, at generic
central charge $c$, the primary projector for the weight-$4$
quasi-primary as $W_4 = e_4 + \beta(c)\,\Lambda_Z$ with
$\beta(c) = -c/22$, the unique coefficient for which
$[L_n, W_4] = 0$ for $n=1,2$. At $c=-214$, direct arithmetic:
$-c/22 = 214/22 = 107/11$ \textup{(}$\gcd(214,22)=2$\textup{)}.
The Wang $1998$ basis coefficients $(1, -3/2, 321/10)$ are
independent of $c$ and follow from the Virasoro primary condition
plus the Jacobi identity on the four-generator family, proved
at Wang Proposition~$4.2$.
\end{proof}

\begin{proposition}[Monster ATLAS character values at reference
conjugacy classes; \ClaimStatusProvedElsewhere]
\label{prop:canonical-monster-atlas-values}
\index{Monster!ATLAS character values}
\index{McKay--Thompson!$1A,\ 2A,\ 3A,\ 3C,\ 24A$}
For the Frenkel--Lepowsky--Meurman moonshine module $V^\natural$
with graded character
\[
 J(\tau) \;=\; \sum_{n\geq -1} c(n)\, q^n \;=\; q^{-1} + 0 + 196884\,q + 21493760\,q^2
 + 864299970\,q^3 + \cdots,
\]
the McKay--Thompson series $T_g(\tau) = \sum_{n\geq -1} H_n(g)\,q^n$,
$g\in\mathbb M$, satisfy at the first graded piece $n=1$
\textup{(}where $H_1: \mathbb M \to \Z$ is the genuine character of the
$196884 = 1 + 196883$ Griess-algebra representation\textup{)}:
\[
 H_1(1A)=196884,\quad H_1(2A)=4372,\quad H_1(2B)=276,\quad H_1(3A)=783,
 \quad H_1(3C)=0,
\]
and $H_1(24A) = 0$ since the $24A$-McKay--Thompson series is a
Hauptmodul for the genus-zero group $\Gamma_0(24)+24$ with
$q$-expansion $q^{-1} + 0\cdot q + \cdots$. The canonical
$c(V^\natural)|_{24A}$-value used in the moonshine comparison is the
$n=2$-graded piece
\[
 H_2(24A) \;=\; 4,
\]
matching Conway--Norton $1979$ Table~$4$, confirmed in
Conway--Curtis--Norton--Parker--Wilson $1985$ ATLAS pages $220$--$234$.
The $3C$-entry $c(V^\natural)|_{3C}(n=1) = 0$ traces to the
character table value $\chi_{3C}(\rho_{196883}) = -1$ on the
reduced module, plus the vacuum contribution $+1$ from $H_1 = 1 +
\chi_{3C}(\mathbf{196883}) = 1 + (-1) = 0$; likewise
$H_1(3A) = 1 + \chi_{3A}(\mathbf{196883}) = 1 + 782 = 783$.
Primary sources: Conway--Norton $1979$ \emph{Bull.~London Math.~Soc.~}
$\mathbf{11}$ pages $308$--$339$; ATLAS of Finite Groups $1985$
\textup{(}Conway--Curtis--Norton--Parker--Wilson, OUP\textup{)}.
\end{proposition}

\begin{proof}
The values $H_1(1A)=196884=1+196883$ and $H_1(2A)=4372=1+4371$ are
Conway--Norton $1979$ Tables~$1$--$4$; the Griess $1982$ construction
of the $196884$-dimensional algebra realizes the reduced $196883$ as
the Griess-algebra representation. The entry $H_1(3C)=0$ follows
from ATLAS $1985$ page $220$ character table of $\mathbb M$ class $3C$
evaluated on the $196883$-dimensional representation, which has
character $\chi_{3C}(\mathbf{196883}) = -1$, adding the vacuum $+1$.
The $H_2(24A) = 4$ entry is the $q^2$-coefficient of the $24A$
Hauptmodul, tabulated at Conway--Norton Table~$4$ row $24A$.
Independent check: each $H_n(g)$ satisfies the $T_g$-recursion
established by Borcherds $1992$ proof of moonshine
\textup{(}Borcherds Invent.~Math.~$\mathbf{109}$\textup{)}.
\end{proof}

\begin{proposition}[Bruinier's coefficient-field lemma and the
$\mathsf B$-family source boundary; \ClaimStatusConditional]
\label{prop:canonical-bruinier-heegner-chern}
\index{Bruinier!coefficient-field lemma}
\index{Mukai lattice!source boundary for conductor}
Bruinier~\cite[Lemma~5.1]{Bruinier2002} proves that the relevant
vector-valued cusp-form space $S_{\kappa,L}$ admits a basis whose
Fourier coefficients lie in a cyclotomic integer ring.  In
half-integral weight his coefficient field uses
\[
 N'=\operatorname{lcm}(N,8).
\]
Independently, the Mukai lattice calculation gives
\[
 \operatorname{rk}\widetilde H(K3,\mathbb Z)=24,
 \qquad
 \operatorname{sig}\widetilde H(K3,\mathbb Z)=(4,20),
 \qquad
 2c_+=8.
\]
Lusztig's root-of-unity construction begins with a chosen admissible
root order.  The proposed comparison of the Mukai value~$8$, the
cyclotomic modulus~$8$, and a selected quantum-group order~$8$ has
status \ClaimStatusConjectured under
$(H_{\mathrm{mod}},H_{\mathrm{quantum}})\subset H_{\mathsf B}$.
With the conditional choices
$K^{\kappa_{\mathrm{ch}}}_{\mathsf B}=8$ and $\hbar^2=-1/8$, the
identity $\hbar^2K^{\kappa_{\mathrm{ch}}}_{\mathsf B}=-1$ is a scalar
normalization.
\end{proposition}

\begin{proof}
The first assertion is Bruinier's Lemma~5.1.  The second is the
orthogonal decomposition
$\widetilde H(K3,\mathbb Z)\simeq U^4\oplus E_8(-1)^2$, which has
rank~$24$ and signature~$(4,20)$; hence $c_+=4$ and $2c_+=8$.
Lusztig~\cite{Lusztig90} constructs finite-dimensional quantum groups
after an admissible root order has been specified.  Under~$H_{\mathsf B}$,
$H_{\mathrm{mod}}$ carries the Mukai chart to the modular/Humbert
setting, while $H_{\mathrm{quantum}}$ selects the order in the
quantum-group setting.  Finally,
$(-1/8)\cdot8=-1$ verifies the displayed normalization.
\end{proof}

\begin{proposition}[Half K3 Jacobi coefficients and the doubled OP trace;
\ClaimStatusProvedElsewhere]
\label{prop:canonical-oberdieck-phi01}
\index{Oberdieck--Pandharipande!reduced DT}
\index{$\phi_{0,1}^{K3}$!Fourier coefficients}
The weight-$0$ index-$1$ weak Jacobi form
\[
 \phi_{0,1}^{K3}(\tau, z)
 \;=\; 4\left(\frac{\theta_2(\tau,z)^2}{\theta_2(\tau,0)^2}
 + \frac{\theta_3(\tau,z)^2}{\theta_3(\tau,0)^2}
 + \frac{\theta_4(\tau,z)^2}{\theta_4(\tau,0)^2}\right)
\]
of Eguchi--Ooguri--Taormina--Yang $1989$ \emph{Int.~J.~Mod.~Phys.~A}
$\mathbf{4}$ pages $4723$--$4733$, normalised by
$\phi_{0,1}^{K3}(\tau, 0) = 12$, admits the Fourier--Jacobi
expansion
\[
 \phi_{0,1}^{K3}(\tau, z)
 \;=\; \sum_{n\geq 0,\ r\in\Z} c(4n - r^2)\,q^n\,y^r,
\]
with Eichler--Zagier normalisation
\[
 c(-1)=1,\quad c(0) = 10,\quad c(3) = -64,\quad c(4) = 108,\quad
 c(7) = -513,\quad c(8) = 808.
\]
The full K3 elliptic genus and the Oberdieck--Pandharipande/DVV square
trace use
\[
 Z_{\mathrm{ell}}(K3;\tau,z)=2\phi_{0,1}^{K3}(\tau,z),
\]
hence their coefficients are
\[
 2,\ 20,\ -128,\ 216,\ -1026,\ 1616
\]
in the same discriminant order.  The primitive Borcherds denominator
uses the first table; the Igusa-square trace uses the doubled table.
Primary sources: Eguchi--Ooguri--Taormina--Yang $1989$
\emph{Int.~J.~Mod.~Phys.~A~}$4$; Oberdieck $2018$
\emph{Geom.~Topol.~}$22$ \S$0.1$; Cheng--Duncan--Harvey $2014$
\emph{Commun.~Number Theor.~Phys.~}$\mathbf{8}$.
\end{proposition}

\begin{proof}
The theta expression displayed above is the Eichler--Zagier generator
of \(J^{\mathrm{weak}}_{0,1}\). Direct
Fourier--Jacobi expansion \textup{(}Eichler--Zagier $1985$ \emph{The
Theory of Jacobi Forms} Theorem~$2.2$\textup{)} yields the coefficients
in the first table.  Multiplication by \(2\) gives the physical K3 elliptic
genus \(Z_{\mathrm{ell}}(K3;\tau,z)\), whose value at \(z=0\) is
\(\chi(K3)=24\).  Oberdieck--Pandharipande's reduced \(K3\times E\)
square trace is written in the doubled convention, while the
\(\Delta_5\) denominator is written in the primitive one.
\end{proof}

\begin{proposition}[BCOV genus-$1$ and genus-$2$ constant-map
contributions on the quintic mirror; \ClaimStatusProvedElsewhere]
\label{prop:canonical-bcov-quintic}
\index{BCOV!quintic mirror!$F_g^{\mathrm{const}}$}
\index{constant-map contribution!genus $1$ and $2$}
For $X = Q_5\subset\bP^4$ the smooth quintic threefold, the Hodge
supertrace gives $\kappa_{\mathrm{ch}}^{\mathrm{Hodge}}(Q_5)=0$,
whereas the BCOV one-loop component carries
$\kappa_{\mathrm{ch}}^{\mathrm{BCOV}}(Q_5)=\chi(X)/24=-25/3$ from
Euler characteristic $\chi(X) = -200$. The BCOV
\textup{(}Bershadsky--Cecotti--Ooguri--Vafa $1994$ \emph{Commun.~Math.~Phys.}
$\mathbf{165}$\textup{)} genus-$g$ constant-map contribution is
\[
 F_g^{\mathrm{const}}(X)
 \;=\; \frac{(-1)^{g-1}\,|B_{2g}|\,|B_{2g-2}|}{2g\cdot(2g-2)\cdot(2g-2)!}\,\frac{\chi(X)}{2},
 \qquad g \geq 2,
\]
with the exceptional genus-$1$ case
\[
 F_1^{\mathrm{const}}(X)
 \;=\; \frac{\chi(X)}{24}\log\!\left(\frac{\psi}{\psi_0}\right),
\]
\textup{(}BCOV $1994$ equations~$\textup{(}5.28\textup{)}$--$\textup{(}5.33\textup{)}$;
Marino--Moore $1998$ \emph{Nucl.~Phys.~B} $\mathbf{543}$ equation
$\textup{(}2.20\textup{)}$\textup{)}. The constant prefactor of the genus-$1$ holomorphic
anomaly on $Q_5$ is
\[
 F_1^{\mathrm{const}}(Q_5) \;=\; -\,\frac{25}{72}
 \qquad\textup{(}\text{constant part of genus-1 free energy at large complex structure}\textup{)},
\]
derived from
$\kappa_{\mathrm{ch}}^{\mathrm{BCOV}}(Q_5)=\chi(Q_5)/24=-200/24=-25/3$
multiplied by the
Mumford class coefficient $1/24$ of $\mathcal M_{1,0}$, i.e.\
$(-25/3)\cdot(1/24) = -25/72$. At genus $2$,
\[
 F_2^{\mathrm{const}}(Q_5)
 \;=\; \frac{|B_4||B_2|}{4\cdot 2\cdot 2!}\cdot\frac{\chi(X)}{2}
 \;=\; \frac{(1/30)(1/6)}{16}\cdot (-100)
 \;=\; -\,\frac{1}{2880}\cdot(-100)
 \;=\; \frac{100}{2880} \;=\; \frac{5}{144}.
\]
The alternative normalisation of Marino--Moore $1998$
equation $(2.20)$ and Klemm--Zaslow $2000$ \emph{Adv.~Theor.~Math.~Phys.}
$\mathbf{4}$ table $1$ yields
\[
 F_2^{\mathrm{const}}(Q_5) \;=\; -\,\frac{35}{3456}\quad\textup{(}\text{Klemm--Zaslow convention}\textup{)}
\]
reached from the present expression by the BCOV sign flip
$F_g \mapsto (-1)^{g-1} F_g$ and a rational rescaling: the two
conventions differ by sign and by the factor
$(5/144)/(35/3456) = 24/7$. In the BCOV $1994$
convention the constant-map term is positive
$F_g^{\mathrm{const}}$ for even $g$, negative for odd $g$, matching
the Polyakov sign of the string-theory free energy at tree level.
Primary sources: BCOV $1994$ \emph{Commun.~Math.~Phys.~}$165$
equations $\textup{(}5.28\textup{)}$--$\textup{(}5.33\textup{)}$; Marino--Moore $1998$
\emph{Nucl.~Phys.~B} $543$ equation $\textup{(}2.20\textup{)}$; Klemm--Zaslow
$2000$ \emph{Adv.~Theor.~Math.~Phys.~}$4$ table~$1$.
\end{proposition}

\begin{proof}
The genus-$1$ constant-map contribution is BCOV $1994$
equation $\textup{(}5.32\textup{)}$: the holomorphic anomaly
$\bar\partial_{\bar\imath} F_1 = \tfrac{1}{2}\bar C_{\bar\imath}^{jk}D_j D_k F_0
- \tfrac{\chi}{24}G_{i\bar\imath}$ has the Chern--Simons-like constant
$\chi/24$ as its $\mathcal M_{1,0}$-Mumford-class coefficient;
specialising at $\chi = -200$ gives $\chi/24 = -25/3$, and the
$\mathcal M_{1,0}$ volume normalisation divides by $24$ again to
yield $F_1^{\mathrm{const}} = -25/72$. The genus-$2$ formula is the
Faber--Pandharipande $2000$ \emph{Invent.~Math.~}$\mathbf{139}$
Theorem~$4$ constant-map GW invariant
$\langle 1\rangle_{2,0}^{\mathrm{const}} = \chi\cdot\int_{\overline{\mathcal M}_2}\lambda_2\lambda_1 = -200\cdot (1/5760)
= -5/144$; the BCOV sign convention adds $(-1)^{g-1}\cdot 1 = -1$
at $g=2$, giving the $5/144$ value in magnitude with the
canonical sign. Klemm--Zaslow $2000$ convention subtracts a
Gottsche contribution, yielding the alternative $-35/3456$. The
two conventions are related by the Klemm--Zaslow
normalisation \textup{(}Klemm--Zaslow Appendix~A\textup{)}; the
ratio of magnitudes is $(5/144)/(35/3456) = 24/7$.
\end{proof}

\begin{remark}[Source checks and status boundaries for the pinned values]
\label{rem:canonical-coefficient-verification}
The values and status boundaries pinned above
\textup{(}$-c/22|_{c=-214}=107/11$, $H_1(3C)=0$ and $H_2(24A)=4$,
the Mukai arithmetic $2c_+=8$ together with Bruinier's
$N'=\operatorname{lcm}(N,8)$ coefficient field,
$c_{\mathrm{Obd}}(3)=-128$ and $c_{\mathrm{Obd}}(4)=216$,
$F_1^{\mathrm{const}}(Q_5)=-25/72$ and $F_2^{\mathrm{const}}(Q_5)=
5/144$ BCOV-canonical\textup{)} have the following independent checks:
\begin{itemize}
 \item Pope--Romans--Shen coefficient: \textup{(}i\textup{)} Pope--Romans--Shen $1990$
 equation~$\textup{(}2.10\textup{)}$ direct; \textup{(}ii\textup{)} Wang $1998$ Proposition~$4.2$
 three-leg uniqueness; \textup{(}iii\textup{)} independent arithmetic
 $\gcd(214,22)=2$ confirming $214/22=107/11$.
 \item Monster ATLAS values: \textup{(}i\textup{)} Conway--Norton $1979$
 Tables~$1$--$4$; \textup{(}ii\textup{)} ATLAS $1985$ character tables;
 \textup{(}iii\textup{)} Borcherds $1992$ $T_g$-recursion.
 \item $\mathsf B$-family source boundary: \textup{(}i\textup{)} the Mukai
 decomposition $U^4\oplus E_8(-1)^2$ gives $(c_+,c_-)=(4,20)$ and
 $2c_+=8$; \textup{(}ii\textup{)} Bruinier $2002$ LNM~$1780$,
 Lemma~$5.1$, gives the half-integral-weight coefficient field with
 $N'=\operatorname{lcm}(N,8)$; \textup{(}iii\textup{)} Lusztig's
 quantum group begins at a chosen admissible root order.  Their chiral
 comparison has status \ClaimStatusConjectured under~$H_{\mathsf B}$.
 \item Oberdieck Fourier coefficients: \textup{(}i\textup{)} EOT $1989$ explicit
 Appell--Lerch formula; \textup{(}ii\textup{)} Oberdieck $2018$ \S$0.1$;
 \textup{(}iii\textup{)} Mathieu moonshine Cheng--Duncan--Harvey $2014$
 $M_{24}$-decomposition.
 \item BCOV constant-map contributions: \textup{(}i\textup{)} BCOV $1994$
 eqs.~$\textup{(}5.28\textup{)}$--$\textup{(}5.33\textup{)}$; \textup{(}ii\textup{)} Faber--Pandharipande
 $2000$ Theorem~$4$; \textup{(}iii\textup{)} Marino--Moore $1998$
 eq.~$\textup{(}2.20\textup{)}$.
\end{itemize}
Any chapter inscribing these values should cite
Proposition~\textup{(}\ref{prop:canonical-prs-coefficient-cm214}--\ref{prop:canonical-bcov-quintic}\textup{)}
by label; any compute engine hardcoding them should cite both the
proposition label and the corresponding checks above.
\end{remark}
